\chapter{Brave-New Formal Lie Groupoids and Infinitesimal Quotients}
\label{ch:formal-lie-groupoids}


\section{Purpose, scope, and the nonlinear starting point}
\label{sec:flg-purpose}

The preceding differential and jet chapters constructed an intrinsic
infinitesimal geometry in the native big affine spectral Nisnevich
$\infty$-topos.  For every morphism
\[
    p:X\longrightarrow S
\]
there is a relative de Rham unit
\[
    \eta_{X/S}:X\longrightarrow X_{\mathrm{dR}/S}
\]
with effective-image factorization
\[
    X
    \xrightarrow{\epsilon_{X/S}}
    Q_{X/S}
    \xrightarrow{\jmath_{X/S}}
    X_{\mathrm{dR}/S},
\]
where $\epsilon_{X/S}$ is an effective epimorphism and
$\jmath_{X/S}$ is a monomorphism.  The associated effective
infinitesimal relation is
\[
    R_{X/S,\bullet}
    \simeq
    \check C_\bullet
    \left(
        X\xrightarrow{\epsilon_{X/S}}Q_{X/S}
    \right),
\]
with
\[
    R_{X/S}
    \simeq
    X\times_{Q_{X/S}}X
    \simeq
    X\times_{X_{\mathrm{dR}/S}}X,
    \qquad
    \lvert R_{X/S,\bullet}\rvert\simeq Q_{X/S}.
\]
The preceding jet chapter further proved that the formal-disk monad,
jet comonad, Cartan descent, and co-Kleisli differential-operator
calculus depend only on the effective quotient
$X\twoheadrightarrow Q_{X/S}$.

The present chapter asks for coherent nonlinear subsystems of this full
infinitesimal transport.  The first construction is more elementary than a
groupoid: for every centred object $X\to Y$ we form its intrinsic relative
de Rham completion
\[
    C_X(Y)
    :=
    Y\times_{Y_{\mathrm{dR}/S}}X_{\mathrm{dR}/S}.
\]
The central observation is that $C_X$ preserves pullbacks, fixes the
initial centred object $(X,\operatorname{id}_X)$, and is the right adjoint to
the inclusion of its fixed points.  Its associated comonad is idempotent.
Applying this pullback-preserving coreflection degreewise to the fully
degenerate unit simplices of a fixed-object groupoid preserves the Segal
pullbacks and hence the groupoid structure.

There are then two categorical layers.  Before formality is imposed, an
\emph{infinitesimally anchored groupoid} on $X/S$ is a fixed-object groupoid
equipped with a morphism
\[
    \mathfrak G_\bullet\longrightarrow R_{X/S,\bullet}.
\]
By groupoid effectivity these are equivalent to quotient factorizations
\[
    X\xrightarrow{q} Q\xrightarrow{r} Q_{X/S},
    \qquad
    rq\simeq\epsilon_{X/S},
\]
with $q$ an effective epimorphism.  Before formality, the morphism $r$ is
part of the anchor data; once the factorization is specified, its being an
effective epimorphism is forced by the effectivity of
$\epsilon_{X/S}$ and is therefore not an additional hypothesis.

A \emph{formal Lie groupoid}, by contrast, is defined without an anchor: it is
a fixed-object groupoid which is already complete along its fully degenerate
unit simplices.  The maximal relation $R_{X/S,\bullet}$ is the completion of
the pair groupoid and is terminal among such formal fixed-object groupoids.
Consequently every formal Lie groupoid acquires a contractibly unique
infinitesimal anchor as a theorem.  The anchored category is retained only for
arbitrary, not-yet-formal groupoids and for quotient comparisons.

The adjective \emph{formal} means that the arrow geometry is already complete
along the fully degenerate unit simplices for the inherited relative de Rham
reflector.  This is a nonlinear condition.  No tangent module, bracket,
spectral Lie algebra, or partition-Lie presentation enters the definition.

The isotropy theory requires one further distinction.  The raw automorphism
object
\[
    I^{\mathrm{raw}}_{\mathfrak G}
    =
    X\times_{X\times_SX}\mathfrak G_1
\]
is a genuine group object over $X$, but left exactness alone does not force it
to be complete along its identity section.  The \emph{formal isotropy group}
is therefore its centred completion
\[
    I^{\mathrm{for}}_{\mathfrak G}
    :=
    C_X(I^{\mathrm{raw}}_{\mathfrak G}).
\]
For a formal groupoid this admits the intrinsic anchor-kernel description
\[
    I^{\mathrm{for}}_{\mathfrak G}
    \simeq
    \mathfrak G_1\times_{R_{X/S}}X.
\]
This is the isotropy object used for the formal bitorsor, adjoint transport,
and nonlinear commutator constructions below.

This organization is compatible with two complementary traditions.  In higher
topos theory, groupoid objects are effective and are therefore controlled by
their quotients \cite{LurieHTT}.  In characteristic-zero derived deformation
theory over a commutative dg-algebra, relative formal moduli problems are
controlled by dg Lie algebroids \cite{Nuiten2019}.  Partition Lie algebroids
admit formal integration to formal leaf spaces over the coherent bases and
under the finiteness hypotheses of the corresponding formal-integration
theory \cite{BrantnerMagidsonNuiten2025}.  Formal groupoids have also long
appeared as nonlinear integrations of Lie algebroids in algebraic geometry;
a particularly relevant example is Kapranov's formal path groupoid
\cite{Kapranov2007}.  None of those algebraic presentations is imported into
the definition here.  The construction remains internal to the spectral
de Rham geometry already established.

\paragraph{Standing ambient convention.}
Let $\mathcal X$ denote the native big affine spectral Nisnevich
$\infty$-topos.  Fix $S\in\mathcal X$, put
\[
    \mathcal E:=\mathcal X_{/S},
\]
and work in this slice.  All limits, colimits, groupoid objects, effective
epimorphisms, monomorphisms, realizations, relative de Rham objects, and
dependent adjoints in the intrinsic part of the chapter are formed in
$\mathcal E$ unless explicitly indicated otherwise.  For brevity write
\[
    L_S(Y):=Y_{\mathrm{dR}/S},
    \qquad
    \eta_Y:Y\longrightarrow L_S(Y).
\]
We retain the universe convention and enlargement fixed in the preceding jet
chapter.  In particular, the groupoid, slice, functor, Eilenberg--Moore, and
coalgebra $\infty$-categories below are formed in that same enlarged universe.

\paragraph{Terminology convention.}
The phrase \emph{formal Lie groupoid} is anticipatory terminology for the
nonlinear objects constructed here.  No representability, smoothness of a
source map, finite-presentation condition, tangent perfectness, or
Lie-algebraic presentation is included in the definition.  Representability
enters only in the final first-order handoff where the manuscript's cotangent
complex has been constructed; smoothness there upgrades the constructed
internal-dual anchor to the finite-locally-free dualizable range.  Later
differentiation or recognition theorems may impose their own geometric
scope.  In particular, the one-object examples of the intrinsic
formal-groupoid package begin with arbitrary group objects in
$\mathcal X_{/S}$.

\paragraph{Scope convention.}
No tangent object, ordinary Lie bracket, spectral Lie algebra, partition Lie
algebra, partition Lie algebroid, Chevalley--Eilenberg algebra, enveloping
algebra, Spencer normalization, or PBW comparison is part of the definition
of a formal Lie groupoid.  Those objects belong to differentiation and
recognition.  The intrinsic formal-groupoid package through
Theorem~\ref{thm:formal-lie-groupoid-package} is entirely nonlinear.  The
final handoff section records only the first represented linearization of the
canonical nonlinear anchor, after that nonlinear package is complete.

\paragraph{Chapter-level output.}
For every $X\to S$ the chapter constructs:
\begin{enumerate}
    \item the pullback-preserving centred relative de Rham completion
    $C_X$, its recognition theorem, coreflection, and associated idempotent
    comonad structure;
    \item the induced idempotent unitwise completion of fixed-object
    groupoids on $X$;
    \item the canonical effective quotient map
    $\epsilon_{X/S}:X\twoheadrightarrow Q_{X/S}$ as a relative de Rham
    equivalence, together with the canonical relation groupoid
    $R_{X/S,\bullet}$ as the completion of the pair groupoid and its
    terminality among formal fixed-object groupoids;
    \item the category of infinitesimally anchored groupoids and its
    classification by quotient factorizations
    \[
        X\twoheadrightarrow Q\twoheadrightarrow Q_{X/S},
    \]
    with effectivity of the second map derived from the composite
    $\epsilon_{X/S}$;
    \item formal Lie groupoids defined solely as the fixed points of unitwise
    completion, together with their contractibly unique constructed
    infinitesimal anchors;
    \item the quotient-side formalization
    $X\twoheadrightarrow\widehat Q_q\to Q$ and the equivalent recognition
    criteria
    \[
        \check C_\bullet(q)\text{ formal}
        \Longleftrightarrow
        q_{\mathrm{dR}/S}\text{ monic}
        \Longleftrightarrow
        q_{\mathrm{dR}/S}\text{ an equivalence}
        \Longleftrightarrow
        \widehat Q_q\xrightarrow{\sim}Q;
    \]
    \item the effective formal leaf space and formal orbit disks;
    \item the raw isotropy group, its completed formal isotropy group, the
    anchor-kernel description, and the descended relative inertia kernel;
    \item raw and formal isotropy bitorsor theorems over the corresponding
    effective arrow images;
    \item the nonlinear formal-isotropy commutator, the formal adjoint group
    constructed as the contractibly unique effective descent of the formal
    inertia kernel, and coherent adjoint transport derived from that descent;
    \item the effective-image relation disk--jet calculus attached to an
    arbitrary morphism $q:X\to B$, including the canonical reduction to its
    effective image, monad, comonad, adjunction, transport equivalence, and
    effective descent, together with its specialization to every formal
    quotient;
    \item functoriality, arbitrary spectral base change, fixed-object
    quotient invariance, the explicit failure of arbitrary-atlas invariance
    of unitwise formality, and the calculation showing that canonical
    unitwise formalization does not restore that invariance;
    \item the canonical formalization of arbitrary internal groupoids;
    \item pair and action groupoid examples, together with the universal
    classifying object $S\twoheadrightarrow BG$, construction and recognition
    of principal right $G$-objects, the classification equivalence
    \[
        \operatorname{Map}_{\mathcal E}(X,BG)
        \simeq
        \operatorname{Prin}_G(X),
    \]
    and the resulting gauge groupoids;
    \item group objects and their constructed formal identity groups as the
    one-object special case;
    \item the translation trivialization and integrated group-valued
    Maurer--Cartan cocycle as consequences of the one-object specialization;
    \item in the represented range, representability of the full simplicial
    groupoid from degrees zero and one, the first structured neighbourhood
    of the unit, the first-order cotangent anchor, its internal-dual tangent
    anchor without smoothness, and the finite-locally-free smooth upgrade.
\end{enumerate}

The logical progression is
\[
\boxed{
\begin{aligned}
\text{relative de Rham geometry}
&\longrightarrow
\text{centred pullback-preserving completion}
\\
&\longrightarrow
\text{formal fixed-object groupoids}
\\
&\longrightarrow
\text{canonical infinitesimal anchor and formal quotients}
\\
&\longrightarrow
\text{formal leaves, formal isotropy, and Cartan descent}.
\end{aligned}}
\]
Differentiation begins only after this nonlinear package is complete.

\section{Effective groupoids in the ambient $\infty$-topos}
\label{sec:flg-effective-groupoids}

\subsection{Groupoid objects with fixed object space}

Let
\[
    \operatorname{Gpd}(\mathcal X_{/S})
    \subset
    \operatorname{Fun}(\Delta^{\mathrm{op}},\mathcal X_{/S})
\]
denote the full $\infty$-category of groupoid objects.  Thus a simplicial object
$\mathfrak G_\bullet$ is a groupoid object when it satisfies the Segal
conditions and every degree-one arrow is invertible.  We write
\[
    \mathfrak G_1
    \mathop{\rightrightarrows}^{s}_{t}
    \mathfrak G_0
\]
for its arrow object, with source $s$, target $t$, unit
$u:\mathfrak G_0\to\mathfrak G_1$, inversion
$\iota:\mathfrak G_1\simeq\mathfrak G_1$, and composition
\[
    m:
    \mathfrak G_1\times_{t,\mathfrak G_0,s}\mathfrak G_1
    \longrightarrow
    \mathfrak G_1.
\]
All coherence is supplied by the simplicial structure.

\begin{definition}[Fixed-object groupoids]
\label{def:fixed-object-groupoids}
For $X\in\mathcal X_{/S}$, let
\[
    \operatorname{Gpd}_X(\mathcal X_{/S})
\]
denote the fibre of evaluation in degree zero
\[
    \operatorname{ev}_0:
    \operatorname{Gpd}(\mathcal X_{/S})
    \longrightarrow
    \mathcal X_{/S}
\]
over $X$, with a specified identification $\mathfrak G_0\simeq X$.
Morphisms in $\operatorname{Gpd}_X(\mathcal X_{/S})$ are therefore
simplicial morphisms whose degree-zero component is identified with
$\operatorname{id}_X$ through the specified fibre data.
\end{definition}

\begin{remark}
The fixed-object formulation is useful because an algebroid is eventually an
infinitesimal structure \emph{over a fixed base object} $X$.  It also makes the
one-object case transparent: when $X=S$, a group object in
$\mathcal X_{/S}$ is exactly a one-object groupoid in
$\operatorname{Gpd}_S(\mathcal X_{/S})$.
\end{remark}

\subsection{Effectivity and quotient objects}

The basic categorical input is the effectivity of groupoid objects in an
$\infty$-topos \cite{LurieHTT}.

\begin{theorem}[Effectivity of internal groupoids]
\label{thm:groupoid-effectivity}
Let $\mathfrak G_\bullet\in
\operatorname{Gpd}(\mathcal X_{/S})$ and put
\[
    Q_{\mathfrak G}:=\lvert\mathfrak G_\bullet\rvert .
\]
Then the canonical map
\[
    q_{\mathfrak G}:
    \mathfrak G_0\longrightarrow Q_{\mathfrak G}
\]
is an effective epimorphism and there is a canonical equivalence of simplicial
objects
\[
    \mathfrak G_\bullet
    \simeq
    \check C_\bullet(q_{\mathfrak G}).
\]
In particular,
\[
    \mathfrak G_n
    \simeq
    \underbrace{
    \mathfrak G_0
    \times_{Q_{\mathfrak G}}\cdots
    \times_{Q_{\mathfrak G}}
    \mathfrak G_0}_{n+1\text{ factors}}.
\]
\end{theorem}

\begin{proof}
This is the effectivity theorem for groupoid objects in an $\infty$-topos.
The map from the degree-zero object to the geometric realization is an effective
epimorphism, and its Cech nerve is canonically equivalent to the original
groupoid object.  The degreewise formula is the Cech description.
\end{proof}

\begin{definition}[Effective quotient of a groupoid]
\label{def:groupoid-quotient}
For $\mathfrak G_\bullet\in
\operatorname{Gpd}_X(\mathcal X_{/S})$, the object
\[
    Q_{\mathfrak G}:=\lvert\mathfrak G_\bullet\rvert
\]
is its effective quotient, and
\[
    q_{\mathfrak G}:X\twoheadrightarrow Q_{\mathfrak G}
\]
is its quotient presentation.
\end{definition}

\begin{proposition}[Fixed-object groupoids and effective quotients]
\label{prop:groupoids-effective-quotients-equivalence}
Let
\[
    \operatorname{Epi}(X)
    \subset
    (\mathcal X_{/S})_{X/}
\]
denote the full $\infty$-subcategory spanned by the effective epimorphisms
$q:X\twoheadrightarrow Q$.  Thus its mapping spaces and all higher coherence
are inherited from the undercategory.  Realization and Cech nerve induce
inverse equivalences
\[
    \operatorname{Gpd}_X(\mathcal X_{/S})
    \simeq
    \operatorname{Epi}(X).
\]
\end{proposition}

\begin{proof}
Theorem~\ref{thm:groupoid-effectivity} sends a fixed-object groupoid to the
effective epimorphism $X\to Q_{\mathfrak G}$.  Conversely, an effective
epimorphism $q:X\to Q$ has Cech nerve
\[
    \check C_n(q)
    =
    \underbrace{X\times_Q\cdots\times_QX}_{n+1},
\]
which is a groupoid object with degree zero $X$.  Since $q$ is effective, its
Cech nerve realizes to $Q$.  The two constructions are inverse by effectivity
and the universal property of geometric realization.
\end{proof}

\begin{remark}[Why quotient language is primary]
\label{rem:quotient-language-primary}
The preceding proposition means that there is no loss of higher groupoid
coherence in passing from $\mathfrak G_\bullet$ to
$X\twoheadrightarrow Q_{\mathfrak G}$.  The quotient description is therefore
not a truncation or a coarse quotient.  It is an equivalent presentation of the
entire internal groupoid in the ambient $\infty$-topos.
\end{remark}


\section{Centred and unitwise relative de Rham completion}
\label{sec:flg-unitwise-formalization}

The formalization of groupoids is controlled by a simpler construction on
objects under a fixed centre.  We first develop that construction in the
undercategory
\[
    \mathcal E_{X/}
    =
    (\mathcal X_{/S})_{X/}.
\]

\subsection{Centred completion}

\begin{definition}[Centred relative de Rham completion]
\label{def:centred-dr-completion}
Let
\[
    u:X\longrightarrow Y
\]
be an object of $\mathcal E_{X/}$.  Define
\[
    C_X(Y,u)
    :=
    Y\times_{L_S(Y)}L_S(X),
\]
where the map
\[
    L_S(X)\longrightarrow L_S(Y)
\]
is $L_S(u)$.  The induced centre
\[
    \widetilde u
    :=
    (u,\eta_X):
    X\longrightarrow C_X(Y,u)
\]
and first projection
\[
    \kappa_{Y,u}:
    C_X(Y,u)\longrightarrow Y
\]
satisfy
\[
    \kappa_{Y,u}\widetilde u\simeq u.
\]
When the centre is understood, write simply $C_X(Y)$ and $\kappa_Y$.
\end{definition}

\begin{proposition}[Pullback preservation]
\label{prop:centred-completion-left-exact}
The assignment
\[
    C_X:\mathcal E_{X/}\longrightarrow\mathcal E_{X/}
\]
is functorial, preserves pullbacks, and fixes the initial centred object:
\[
    C_X(X,\operatorname{id}_X)
    \simeq
    (X,\operatorname{id}_X).
\]
The morphisms $\kappa_Y:C_X(Y)\to Y$ assemble into a natural transformation
\[
    \kappa:C_X\longrightarrow\operatorname{id}_{\mathcal E_{X/}}.
\]
\end{proposition}

\begin{proof}
A morphism
\[
    f:(X\xrightarrow{u}Y)\longrightarrow(X\xrightarrow{v}Z)
\]
in the undercategory satisfies $fu\simeq v$.  Naturality of the relative
de Rham unit gives a commutative square
\[
\begin{tikzcd}
L_S(X) \arrow[r,"L_S(u)"] \arrow[d,equal] &
L_S(Y) \arrow[d,"L_S(f)"]\\
L_S(X) \arrow[r,"L_S(v)"'] &
L_S(Z),
\end{tikzcd}
\]
and the universal property of the defining pullbacks gives
$C_X(f):C_X(Y)\to C_X(Z)$.  The construction is functorial by functoriality
of pullbacks.  The identity
\[
    C_X(X,\operatorname{id}_X)
    =
    X\times_{L_S(X)}L_S(X)
    \simeq X
\]
is compatible with the identity centre.

It remains to prove pullback preservation.  Let
\[
\begin{tikzcd}
Y \arrow[r] \arrow[d] &
Y_1 \arrow[d]\\
Y_2 \arrow[r] &
Y_0
\end{tikzcd}
\]
be a pullback square in $\mathcal E_{X/}$.  Its underlying square in
$\mathcal E$ is Cartesian, and the four centres are induced by the same
object $X$.  Since $L_S$ is left exact,
\[
    L_S(Y)
    \simeq
    L_S(Y_1)\times_{L_S(Y_0)}L_S(Y_2).
\]
Pullback pasting applied to the defining squares of the four centred
completions gives
\[
\begin{aligned}
    C_X(Y)
    &=
    Y\times_{L_S(Y)}L_S(X)\\
    &\simeq
    C_X(Y_1)\times_{C_X(Y_0)}C_X(Y_2).
\end{aligned}
\]
This is the canonical pullback comparison map in the undercategory.
Naturality of $\kappa$ is immediate from the same defining pullback diagrams.
\end{proof}

\begin{proposition}[Localized centre and recognition]
\label{prop:centred-completion-recognition}
For every centred object $u:X\to Y$ there is a canonical equivalence
\[
    L_S(C_X(Y))
    \xrightarrow{\sim}
    L_S(X)
\]
under which the localized completed centre
\[
    L_S(\widetilde u):
    L_S(X)\longrightarrow L_S(C_X(Y))
\]
is an equivalence.  Moreover the following are equivalent:
\begin{enumerate}
    \item $\kappa_Y:C_X(Y)\to Y$ is an equivalence;
    \item $L_S(u):L_S(X)\to L_S(Y)$ is an equivalence.
\end{enumerate}
\end{proposition}

\begin{proof}
Left exactness and idempotence give
\[
\begin{aligned}
    L_S(C_X(Y))
    &\simeq
    L_S(Y)
    \times_{L_S(L_S(Y))}
    L_S(L_S(X))\\
    &\simeq
    L_S(Y)
    \times_{L_S(Y)}
    L_S(X)\\
    &\simeq
    L_S(X).
\end{aligned}
\]
Tracing the centre through this equivalence identifies
$L_S(\widetilde u)$ with the identity of $L_S(X)$.

If $L_S(u)$ is an equivalence, then $\kappa_Y$ is its base change along
$\eta_Y:Y\to L_S(Y)$ and is therefore an equivalence.  Conversely, if
$\kappa_Y$ is an equivalence, then applying $L_S$ and using
\[
    \kappa_Y\widetilde u\simeq u
\]
shows that $L_S(u)$ is a composite of equivalences.
\end{proof}

\begin{definition}[Complete centred object]
\label{def:complete-centred-object}
A centred object $u:X\to Y$ is \emph{relative de Rham-complete along $X$}
if $\kappa_Y:C_X(Y)\to Y$ is an equivalence.  Write
\[
    \mathcal E^{\mathrm{for}}_{X/}
    \subset
    \mathcal E_{X/}
\]
for the full subcategory of complete centred objects.
\end{definition}

\begin{theorem}[Centred formalization coreflection]
\label{thm:centred-formalization-coreflection}
The inclusion
\[
    \mathcal E^{\mathrm{for}}_{X/}
    \hookrightarrow
    \mathcal E_{X/}
\]
admits $C_X$ as a right adjoint.  Equivalently, if $F$ is complete, then
composition with $\kappa_Y$ induces a canonical equivalence
\[
    \operatorname{Map}_{\mathcal E_{X/}}
    (F,C_X(Y))
    \xrightarrow{\sim}
    \operatorname{Map}_{\mathcal E_{X/}}
    (F,Y).
\]
\end{theorem}

\begin{proof}
Proposition~\ref{prop:centred-completion-recognition} applied to the completed
centre shows first that $C_X(Y)$ is complete for every centred object $Y$.

Let
\[
    v:X\longrightarrow F
\]
be the centre of a complete object $F$.  By
Proposition~\ref{prop:centred-completion-recognition},
\[
    L_S(v):L_S(X)\xrightarrow{\sim}L_S(F)
\]
is an equivalence.  Hence for every relative de Rham-local object
$Z\in\mathcal E$, precomposition with $v$ is an equivalence
\[
    v^*:
    \operatorname{Map}_{\mathcal E}(F,Z)
    \xrightarrow{\sim}
    \operatorname{Map}_{\mathcal E}(X,Z).
\]
Indeed,
\[
\begin{aligned}
    \operatorname{Map}_{\mathcal E}(F,Z)
    &\simeq
    \operatorname{Map}_{\mathcal E}(L_S(F),Z)\\
    &\simeq
    \operatorname{Map}_{\mathcal E}(L_S(X),Z)\\
    &\simeq
    \operatorname{Map}_{\mathcal E}(X,Z).
\end{aligned}
\]

Regard $L_S(X)$ as a centred object through $\eta_X$ and
$L_S(Y)$ as a centred object through
\[
    X\xrightarrow{u}Y\xrightarrow{\eta_Y}L_S(Y).
\]
Since both targets are relative de Rham-local, the corresponding
undercategory mapping spaces are the homotopy fibres of the preceding
restriction equivalences over their specified centres.  Consequently
\[
    \operatorname{Map}_{\mathcal E_{X/}}(F,L_S(X))
    \simeq *,
    \qquad
    \operatorname{Map}_{\mathcal E_{X/}}(F,L_S(Y))
    \simeq *.
\]

The defining pullback
\[
    C_X(Y)
    =
    Y\times_{L_S(Y)}L_S(X)
\]
is also a pullback in $\mathcal E_{X/}$.  Mapping the complete centred object
$F$ into this pullback therefore gives
\[
\begin{aligned}
    \operatorname{Map}_{\mathcal E_{X/}}(F,C_X(Y))
    &\simeq
    \operatorname{Map}_{\mathcal E_{X/}}(F,Y)
    \times_{\operatorname{Map}_{\mathcal E_{X/}}(F,L_S(Y))}
    \operatorname{Map}_{\mathcal E_{X/}}(F,L_S(X))\\
    &\simeq
    \operatorname{Map}_{\mathcal E_{X/}}(F,Y).
\end{aligned}
\]
Under this equivalence the projection to the first factor is composition with
$\kappa_Y$.  Hence $C_X$ is right adjoint to the inclusion of complete
centred objects.
\end{proof}


\begin{theorem}[Idempotent centred formalization]
\label{thm:centred-completion-idempotent}
The coreflection of
Theorem~\ref{thm:centred-formalization-coreflection} equips $C_X$ with a
canonical idempotent comonad structure whose counit is
$\kappa:C_X\to\operatorname{id}$.  In particular there is a canonical natural
equivalence
\[
    C_XC_X
    \xrightarrow{\sim}
    C_X.
\]
\end{theorem}

\begin{proof}
Let
\[
    i:
    \mathcal E^{\mathrm{for}}_{X/}
    \hookrightarrow
    \mathcal E_{X/}
\]
denote the fully faithful inclusion.  The preceding theorem identifies
$C_X$ with the composite endofunctor of the adjunction
\[
    i\dashv C_X
\]
after regarding the right adjoint as landing in
$\mathcal E^{\mathrm{for}}_{X/}$.  The associated adjunction therefore
supplies the comonad structure on the endofunctor $C_X$, with counit
$\kappa$.

Because $i$ is fully faithful, the unit of the adjunction is an equivalence.
Equivalently, every object in the image of $C_X$ is fixed by $C_X$.  Hence
the comultiplication is an equivalence and the comonad is idempotent, giving
the displayed natural equivalence.
\end{proof}

\begin{proposition}[Functoriality in the centre]
\label{prop:centred-completion-varying-centre}
Suppose
\[
\begin{tikzcd}
X \arrow[r,"u"] \arrow[d,"f"'] &
Y \arrow[d,"g"]\\
X' \arrow[r,"u'"'] &
Y'
\end{tikzcd}
\]
is a commutative square in $\mathcal E$.  Then there is a canonical morphism
\[
    C_X(Y,u)
    \longrightarrow
    C_{X'}(Y',u')
\]
whose components are $g:Y\to Y'$ and
$L_S(f):L_S(X)\to L_S(X')$.  These morphisms preserve identities,
composition, and all higher functorial coherences.
\end{proposition}

\begin{proof}
Naturality of $\eta$ and commutativity of the given square supply a morphism
between the two defining pullback diagrams.  Functoriality of finite limits
gives the asserted map and all composition coherences.
\end{proof}

\begin{proposition}[Base change of centred completion]
\label{prop:centred-completion-base-change}
Let $c:S'\to S$, put
\[
    X':=X\times_SS',
    \qquad
    Y':=Y\times_SS',
\]
and let $u':X'\to Y'$ be the pullback of $u$.  Then there is a canonical
equivalence
\[
    C_{X'}(Y',u')
    \simeq
    C_X(Y,u)\times_SS',
\]
compatible with $\kappa$, the comonad structure, identities, iterated base
change, and higher pullback-pasting coherences.
\end{proposition}

\begin{proof}
The relative de Rham base-change theorem gives
\[
    L_{S'}(X')\simeq L_S(X)\times_SS',
    \qquad
    L_{S'}(Y')\simeq L_S(Y)\times_SS'.
\]
Substitute these into the defining pullback and commute finite limits with
base change.  Every displayed structure is constructed from these same
pullbacks, so the compatibility statements follow formally.
\end{proof}

\subsection{Unitwise completion of fixed-object groupoids}

Fix
\[
    \mathfrak G_\bullet
    \in
    \operatorname{Gpd}_X(\mathcal X_{/S}).
\]
For every $n\geq0$, let
\[
    u_n:X\longrightarrow\mathfrak G_n
\]
be the fully degenerate unit simplex.

\begin{definition}[Unitwise relative de Rham completion]
\label{def:unitwise-groupoid-completion}
Define
\[
    \widehat{\mathfrak G}^{\,\mathrm{dR}}_n
    :=
    C_X(\mathfrak G_n,u_n)
    =
    \mathfrak G_n
    \times_{(\mathfrak G_n)_{\mathrm{dR}/S}}
    X_{\mathrm{dR}/S}.
\]
Write
\[
    \kappa_{\mathfrak G,n}:
    \widehat{\mathfrak G}^{\,\mathrm{dR}}_n
    \longrightarrow
    \mathfrak G_n
\]
for the first projection.
\end{definition}

\begin{proposition}[Unitwise completion is a groupoid]
\label{prop:unitwise-completion-groupoid}
The objects
$\widehat{\mathfrak G}^{\,\mathrm{dR}}_n$ assemble canonically into a
groupoid object
\[
    \widehat{\mathfrak G}^{\,\mathrm{dR}}_\bullet
    \rightrightarrows X.
\]
The projections assemble into a natural fixed-object groupoid morphism
\[
    \kappa_{\mathfrak G}:
    \widehat{\mathfrak G}^{\,\mathrm{dR}}_\bullet
    \longrightarrow
    \mathfrak G_\bullet.
\]
Thus unitwise completion defines an endofunctor
\[
    \widehat{(-)}^{\,\mathrm{dR}}:
    \operatorname{Gpd}_X(\mathcal X_{/S})
    \longrightarrow
    \operatorname{Gpd}_X(\mathcal X_{/S}).
\]
For $n\ge2$ there are canonical Segal equivalences
\[
    \widehat{\mathfrak G}^{\,\mathrm{dR}}_n
    \simeq
    \underbrace{
    \widehat{\mathfrak G}^{\,\mathrm{dR}}_1
    \times_X\cdots\times_X
    \widehat{\mathfrak G}^{\,\mathrm{dR}}_1
    }_{n\text{ factors}}.
\]
\end{proposition}

\begin{proof}
Every simplicial operator $\alpha:[m]\to[n]$ satisfies
\[
    \alpha^*u_n\simeq u_m.
\]
Proposition~\ref{prop:centred-completion-varying-centre}, with the centre map
equal to $\operatorname{id}_X$, therefore gives the simplicial structure and
functoriality.

For $n\geq2$, the Segal equivalence expresses $\mathfrak G_n$ as the
iterated pullback
\[
    \mathfrak G_n
    \simeq
    \underbrace{
    \mathfrak G_1\times_X\cdots\times_X\mathfrak G_1
    }_{n\text{ factors}}
\]
in $\mathcal E_{X/}$.  Proposition~\ref{prop:centred-completion-left-exact}
preserves these pullbacks and fixes $(X,\operatorname{id}_X)$, giving the
displayed completed Segal equivalences.  The inversion map preserves the unit
centre and therefore completes to an inversion.  Composition is supplied by
the completed Segal structure, equivalently by applying $C_X$ to the original
composition diagram.  All groupoid identities and higher coherences are images
of the original ones under the functor $C_X$.
\end{proof}

\begin{theorem}[Idempotence of unitwise completion]
\label{thm:unitwise-completion-idempotent}
Unitwise completion is an idempotent comonad on
$\operatorname{Gpd}_X(\mathcal X_{/S})$.  In particular there is a canonical
natural equivalence
\[
    \widehat{
        \widehat{\mathfrak G}^{\,\mathrm{dR}}
    }^{\,\mathrm{dR}}_\bullet
    \simeq
    \widehat{\mathfrak G}^{\,\mathrm{dR}}_\bullet,
\]
and
\[
    \left(
        \widehat{\mathfrak G}^{\,\mathrm{dR}}_n
    \right)_{\mathrm{dR}/S}
    \simeq
    X_{\mathrm{dR}/S}
\]
for every $n$.
\end{theorem}

\begin{proof}
Apply Theorem~\ref{thm:centred-completion-idempotent} degreewise.  Since the
comonad maps are natural for morphisms of centred objects, they commute with
all simplicial operators and therefore define the stated groupoid-level
comonad.
\end{proof}

\begin{definition}[Unitwise de Rham-complete groupoid]
\label{def:unitwise-complete-groupoid}
A fixed-object groupoid $\mathfrak G_\bullet$ is
\emph{unitwise relative de Rham-complete}, or simply \emph{formal}, if
\[
    \kappa_{\mathfrak G}:
    \widehat{\mathfrak G}^{\,\mathrm{dR}}_\bullet
    \xrightarrow{\sim}
    \mathfrak G_\bullet
\]
is an equivalence.  Write
\[
    \operatorname{Gpd}^{\mathrm{for}}_X(\mathcal X_{/S})
    \subset
    \operatorname{Gpd}_X(\mathcal X_{/S})
\]
for the full subcategory of formal fixed-object groupoids.
\end{definition}

\begin{proposition}[Degree-one recognition of unitwise formality]
\label{prop:degree-one-formality}
For a fixed-object groupoid $\mathfrak G_\bullet$, the following are
equivalent:
\begin{enumerate}
    \item $\mathfrak G_\bullet$ is unitwise relative de Rham-complete;
    \item the canonical degree-one projection
    \[
        \mathfrak G_1
        \times_{(\mathfrak G_1)_{\mathrm{dR}/S}}
        X_{\mathrm{dR}/S}
        \longrightarrow
        \mathfrak G_1
    \]
    is an equivalence;
    \item the localized unit
    \[
        (u_1)_{\mathrm{dR}/S}:
        X_{\mathrm{dR}/S}
        \xrightarrow{\sim}
        (\mathfrak G_1)_{\mathrm{dR}/S}
    \]
    is an equivalence.
\end{enumerate}
\end{proposition}

\begin{proof}
The equivalence of (2) and (3) is
Proposition~\ref{prop:centred-completion-recognition}.  Both
$\mathfrak G_\bullet$ and
$\widehat{\mathfrak G}^{\,\mathrm{dR}}_\bullet$ are Segal groupoids with
degree zero $X$, and under their Segal equivalences
$\kappa_{\mathfrak G,n}$ is the $n$-fold fibre product over $X$ of
$\kappa_{\mathfrak G,1}$.  Hence (1) and (2) are equivalent.
\end{proof}

\begin{theorem}[Formalization coreflection]
\label{thm:formalization-coreflection}
The inclusion
\[
    \operatorname{Gpd}^{\mathrm{for}}_X(\mathcal X_{/S})
    \hookrightarrow
    \operatorname{Gpd}_X(\mathcal X_{/S})
\]
admits
\[
    \widehat{(-)}^{\,\mathrm{dR}}
\]
as a right adjoint.  Equivalently, for every formal
$\mathfrak F_\bullet$ and every fixed-object groupoid
$\mathfrak G_\bullet$, composition with $\kappa_{\mathfrak G}$ induces
\[
    \operatorname{Map}
    \left(
        \mathfrak F_\bullet,
        \widehat{\mathfrak G}^{\,\mathrm{dR}}_\bullet
    \right)
    \xrightarrow{\sim}
    \operatorname{Map}
    \left(
        \mathfrak F_\bullet,
        \mathfrak G_\bullet
    \right).
\]
\end{theorem}

\begin{proof}
Apply Theorem~\ref{thm:centred-formalization-coreflection} pointwise in the
functor category
\[
    \operatorname{Fun}
    \left(
        \Delta^{\mathrm{op}},
        \mathcal E_{X/}
    \right).
\]
Pointwise adjunction gives a right adjoint obtained by applying $C_X$ in every
simplicial degree, with counit the degreewise maps $\kappa$.  By
Proposition~\ref{prop:unitwise-completion-groupoid}, this right adjoint carries
fixed-object groupoid objects to fixed-object groupoid objects, and a formal
fixed-object groupoid is pointwise complete by
Proposition~\ref{prop:degree-one-formality} and the Segal equivalences.
Since groupoid objects form a full subcategory of the simplicial functor
category, the pointwise mapping-space adjunction restricts to
\[
    \operatorname{Map}
    \left(
        \mathfrak F_\bullet,
        \widehat{\mathfrak G}^{\,\mathrm{dR}}_\bullet
    \right)
    \simeq
    \operatorname{Map}
    \left(
        \mathfrak F_\bullet,
        \mathfrak G_\bullet
    \right)
\]
for formal $\mathfrak F_\bullet$.  This is the required coreflection.
\end{proof}

\begin{remark}
The natural transformation points from the completed groupoid to the original
one, so formalization is a coreflection.  The comonad and its coreflection are
not additional axioms: they are forced by the previously constructed
left-exact relative de Rham reflector and the defining pullbacks.
\end{remark}


\section{The canonical relative infinitesimal groupoid}
\label{sec:flg-canonical-relation}

Retain
\[
    \eta_{X/S}:X\to X_{\mathrm{dR}/S}
\]
and its effective-image factorization
\[
    X\xrightarrow{\epsilon_{X/S}}Q_{X/S}
      \xrightarrow{\jmath_{X/S}}X_{\mathrm{dR}/S}.
\]

\begin{definition}[Canonical relative infinitesimal groupoid]
\label{def:canonical-relative-infinitesimal-groupoid}
Define
\[
    R_{X/S,\bullet}
    :=
    \check C_\bullet(\epsilon_{X/S}),
    \qquad
    R_{X/S}:=R_{X/S,1}.
\]
\end{definition}

\begin{lemma}[The effective infinitesimal quotient is a relative de Rham equivalence]
\label{lem:effective-quotient-dr-equivalence}
The effective epimorphism
\[
    \epsilon_{X/S}:X\twoheadrightarrow Q_{X/S}
\]
is inverted by the relative de Rham localization:
\[
    L_S(\epsilon_{X/S}):
    L_S(X)\xrightarrow{\sim}L_S(Q_{X/S}).
\]
Moreover
\[
    L_S(\jmath_{X/S}):
    L_S(Q_{X/S})
    \xrightarrow{\sim}
    L_S(L_S(X))
    \simeq
    L_S(X)
\]
is an equivalence.
\end{lemma}

\begin{proof}
Apply $L_S$ to the effective-epimorphism--monomorphism factorization
\[
    X\xrightarrow{\epsilon_{X/S}}Q_{X/S}
      \xrightarrow{\jmath_{X/S}}L_S(X)
\]
of the localization unit.  The preceding de Rham chapter proves separately
that relative de Rham localization carries effective epimorphisms to
effective epimorphisms, so
\[
    L_S(\epsilon_{X/S}):
    L_S(X)\longrightarrow L_S(Q_{X/S})
\]
is an effective epimorphism.  Since $L_S$ is left exact, it preserves
monomorphisms, and therefore
\[
    L_S(\jmath_{X/S}):
    L_S(Q_{X/S})\longrightarrow L_S(L_S(X))
\]
is a monomorphism.

Their composite is
\[
    L_S(\eta_{X/S}):
    L_S(X)\longrightarrow L_S(L_S(X)),
\]
which is an equivalence by idempotence of the localization.  Hence the two
displayed morphisms form the effective-epimorphism--monomorphism
factorization of an equivalence.  By uniqueness of that factorization, both
are equivalences.
\end{proof}

The preceding de Rham chapter gives canonical equivalences
\[
    R_{X/S,n}
    \simeq
    \underbrace{
    X\times_{Q_{X/S}}\cdots\times_{Q_{X/S}}X
    }_{n+1}
    \simeq
    \underbrace{
    X\times_{X_{\mathrm{dR}/S}}\cdots
      \times_{X_{\mathrm{dR}/S}}X
    }_{n+1}.
\]
For degree one we identify the Chapter~11 notation by
\[
    s=c=\operatorname{pr}_1,
    \qquad
    t=e=\operatorname{pr}_2.
\]
Thus source is the centre projection and target is the evaluation projection.

\subsection{The pair groupoid and maximal formal relation}

Define the relative pair groupoid by
\[
    \operatorname{Pair}(X/S)_n
    :=
    X^{\times_S(n+1)}.
\]
It is the relative $0$-coskeleton of $X$ and hence is terminal among
simplicial objects with specified degree-zero object $X$.

\begin{theorem}[Pair-groupoid completion]
\label{thm:pair-groupoid-completion}
There is a canonical equivalence of fixed-object formal groupoids
\[
    \widehat{\operatorname{Pair}(X/S)}^{\,\mathrm{dR}}_\bullet
    \xrightarrow{\sim}
    R_{X/S,\bullet}.
\]
\end{theorem}

\begin{proof}
In degree $n$, the centre of the pair groupoid is the $(n+1)$-fold diagonal
\[
    \Delta_n:X\longrightarrow X^{\times_S(n+1)}.
\]
Left exactness gives
\[
    L_S\!\left(X^{\times_S(n+1)}\right)
    \simeq
    L_S(X)^{\times_S(n+1)}.
\]
Therefore
\[
\begin{aligned}
    C_X\!\left(X^{\times_S(n+1)},\Delta_n\right)
    &\simeq
    X^{\times_S(n+1)}
    \times_{L_S(X)^{\times_S(n+1)}}
    L_S(X)\\
    &\simeq
    \underbrace{
    X\times_{L_S(X)}\cdots\times_{L_S(X)}X
    }_{n+1}\\
    &\simeq
    R_{X/S,n}.
\end{aligned}
\]
The construction uses only the product projections and diagonal, so the
equivalences commute with every face and degeneracy map.
\end{proof}

\begin{theorem}[Terminality of the maximal formal relation]
\label{thm:maximal-formal-terminal}
The groupoid $R_{X/S,\bullet}$ is terminal in
\[
    \operatorname{Gpd}^{\mathrm{for}}_X(\mathcal X_{/S}).
\]
Equivalently, for every formal fixed-object groupoid $\mathfrak F_\bullet$,
the mapping space
\[
    \operatorname{Map}_{\operatorname{Gpd}_X}
    \left(
        \mathfrak F_\bullet,
        R_{X/S,\bullet}
    \right)
\]
is contractible.
\end{theorem}

\begin{proof}
The relative pair groupoid
$\operatorname{Pair}(X/S)_\bullet$ is the relative $0$-coskeleton of $X$ and
is therefore terminal in
$\operatorname{Gpd}_X(\mathcal X_{/S})$.
Theorem~\ref{thm:formalization-coreflection} identifies unitwise completion
as the right adjoint to the inclusion of formal fixed-object groupoids.
A right adjoint carries terminal objects to terminal objects.  Hence
\[
    \widehat{\operatorname{Pair}(X/S)}^{\,\mathrm{dR}}_\bullet
\]
is terminal in
$\operatorname{Gpd}^{\mathrm{for}}_X(\mathcal X_{/S})$.
By Theorem~\ref{thm:pair-groupoid-completion} this completed pair groupoid is
canonically equivalent to $R_{X/S,\bullet}$.
\end{proof}


\begin{corollary}[Canonical infinitesimal anchor of a formal groupoid]
\label{cor:canonical-formal-anchor}
Every formal fixed-object groupoid
$\mathfrak F_\bullet\in
\operatorname{Gpd}^{\mathrm{for}}_X(\mathcal X_{/S})$
admits a contractibly unique morphism
\[
    \phi_{\mathfrak F}:
    \mathfrak F_\bullet
    \longrightarrow
    R_{X/S,\bullet}.
\]
We call it the \emph{canonical infinitesimal anchor}.
\end{corollary}

\begin{proof}
This is Theorem~\ref{thm:maximal-formal-terminal}.
\end{proof}

\begin{proposition}[The canonical anchor is the completed source--target map]
\label{prop:canonical-anchor-completed-source-target}
Let
\[
    \mathfrak G_\bullet
    \in
    \operatorname{Gpd}^{\mathrm{for}}_X(\mathcal X_{/S}).
\]
In degree one, the canonical anchor of
Corollary~\ref{cor:canonical-formal-anchor} is canonically equivalent to the
completion of the source--target map:
\[
\begin{tikzcd}
C_X(\mathfrak G_1,u)
    \arrow[r,"{C_X(s,t)}"]
    \arrow[d,"{\kappa_{\mathfrak G,1}}"',"\sim"]
&
C_X(X\times_SX,\Delta_X)
    \arrow[d,"\sim"]
\\
\mathfrak G_1
    \arrow[r,"{\phi_{\mathfrak G,1}}"']
&
R_{X/S}.
\end{tikzcd}
\]
Thus, after the canonical identifications supplied by formality and
Theorem~\ref{thm:pair-groupoid-completion},
\[
    \phi_{\mathfrak G,1}
    \simeq
    C_X(s,t).
\]
\end{proposition}

\begin{proof}
The canonical vertex morphism
\[
    \mathfrak G_\bullet
    \longrightarrow
    \operatorname{Pair}(X/S)_\bullet
\]
has degree-one component $(s,t):\mathfrak G_1\to X\times_SX$ and carries the
unit centre to the diagonal.  Applying unitwise completion gives
\[
    \widehat{\mathfrak G}^{\,\mathrm{dR}}_\bullet
    \longrightarrow
    \widehat{\operatorname{Pair}(X/S)}^{\,\mathrm{dR}}_\bullet.
\]
The source identifies with $\mathfrak G_\bullet$ because $\mathfrak G$ is
formal, and the target identifies with $R_{X/S,\bullet}$ by
Theorem~\ref{thm:pair-groupoid-completion}.  This is the canonical anchor
constructed in Theorem~\ref{thm:maximal-formal-terminal}; its degree-one
component is therefore the displayed completed source--target map.
\end{proof}

\begin{proposition}[Extremal formal groupoids]
\label{prop:extremal-infinitesimal-groupoids}
There are canonical formal groupoids on $X$:
\begin{enumerate}
    \item the identity or zero groupoid
    \[
        0_{X,\bullet}
        :=
        \check C_\bullet(\operatorname{id}_X)
        \simeq
        \operatorname{const}(X),
    \]
    whose quotient is $X$;
    \item the maximal infinitesimal groupoid
    $R_{X/S,\bullet}$, whose quotient is $Q_{X/S}$.
\end{enumerate}
For every formal fixed-object groupoid $\mathfrak G_\bullet$ there are
canonical morphisms
\[
    0_{X,\bullet}
    \longrightarrow
    \mathfrak G_\bullet
    \longrightarrow
    R_{X/S,\bullet}.
\]
The first is initial and the second terminal in the formal fixed-object
category.
\end{proposition}

\begin{proof}
The zero groupoid is formal because its degree-one centre is
$\operatorname{id}_X$.  A fixed-object functor from the constant identity
groupoid to any fixed-object groupoid is forced by the unit maps, so the zero
groupoid is initial.  The maximal relation is formal by
Theorem~\ref{thm:pair-groupoid-completion} and idempotence of formalization,
and is terminal by Theorem~\ref{thm:maximal-formal-terminal}.
\end{proof}

\begin{remark}
The word ``zero'' refers to the absence of additional relation arrows between
distinct objects.  Because the ambient category is an $\infty$-topos, this
terminology must not be read as a claim that all homotopy-theoretic inertia of
the object $X$ itself vanishes.  The raw and formal isotropy objects below
make this distinction explicit.
\end{remark}


\section{Formal Lie groupoids and their quotient classification}
\label{sec:flg-definition-classification}

\subsection{Anchored and formal groupoids}

\begin{definition}[Infinitesimally anchored groupoid]
\label{def:infinitesimally-anchored-groupoid}
An \emph{infinitesimally anchored groupoid on $X/S$} is an object
\[
    \mathfrak G_\bullet
    \in
    \operatorname{Gpd}_X(\mathcal X_{/S})
\]
equipped with a morphism
\[
    \phi_{\mathfrak G}:
    \mathfrak G_\bullet
    \longrightarrow
    R_{X/S,\bullet}.
\]
The $\infty$-category of anchored groupoids is
\[
    \operatorname{AnchGpd}(X/S)
    :=
    \left(
        \operatorname{Gpd}_X(\mathcal X_{/S})
    \right)_{/R_{X/S,\bullet}}.
\]
\end{definition}

\begin{definition}[Brave-new formal Lie groupoid]
\label{def:formal-lie-groupoid}
A \emph{brave-new formal Lie groupoid on $X/S$} is a fixed-object groupoid
\[
    \mathfrak G_\bullet
    \in
    \operatorname{Gpd}_X(\mathcal X_{/S})
\]
which is unitwise relative de Rham-complete.  Define
\[
    \operatorname{FormLieGpd}(X/S)
    :=
    \operatorname{Gpd}^{\mathrm{for}}_X(\mathcal X_{/S}).
\]
No infinitesimal anchor is part of this definition.
\end{definition}

\begin{theorem}[The anchor is constructed, not additional formal data]
\label{thm:formal-anchor-forgetful-equivalence}
The canonical anchor of
Corollary~\ref{cor:canonical-formal-anchor} defines a fully faithful functor
\[
    \operatorname{can}_{X/S}:
    \operatorname{FormLieGpd}(X/S)
    \longrightarrow
    \operatorname{AnchGpd}(X/S).
\]
Its essential image is the full subcategory of anchored groupoids whose
underlying fixed-object groupoid is formal.  Forgetting the anchor on this
essential image is inverse to $\operatorname{can}_{X/S}$.
\end{theorem}

\begin{proof}
For every formal fixed-object groupoid $\mathfrak F_\bullet$,
Theorem~\ref{thm:maximal-formal-terminal} gives
\[
    \operatorname{Map}_{\operatorname{Gpd}_X}
    \left(
        \mathfrak F_\bullet,
        R_{X/S,\bullet}
    \right)
    \simeq *.
\]
Hence the space of anchors on $\mathfrak F_\bullet$ is contractible.

Let $\mathfrak F_\bullet,\mathfrak G_\bullet$ be formal.  Any fixed-object
groupoid morphism
\[
    F:\mathfrak F_\bullet\longrightarrow\mathfrak G_\bullet
\]
gives two maps
\[
    \phi_{\mathfrak G}F,\ \phi_{\mathfrak F}:
    \mathfrak F_\bullet\rightrightarrows R_{X/S,\bullet}.
\]
The mapping space from $\mathfrak F_\bullet$ to $R_{X/S,\bullet}$ is
contractible, so the space of compatible homotopies
\[
    \phi_{\mathfrak G}F\simeq\phi_{\mathfrak F}
\]
is contractible.  Therefore passing from fixed-object morphisms to anchored
morphisms changes no mapping space.  This proves full faithfulness and shows
that every anchored formal object lies, through a contractible space of
identifications, in the essential image.
\end{proof}

\begin{remark}[No monomorphism condition]
\label{rem:no-mono-formal-gpd}
The canonical anchor
\[
    \mathfrak G_1\longrightarrow R_{X/S}
\]
is not required to be a monomorphism.  Its homotopy fibres carry genuine
isotropy information.  Requiring a monomorphism would erase those fibres and,
in particular, would exclude one-object formal groups and completed gauge
groupoids.  The formality condition is instead completion along the unit
section.
\end{remark}

\subsection{Canonical formalization of arbitrary internal groupoids}

\begin{theorem}[Unitwise Formal Groupoid Completion Theorem]
\label{thm:unitwise-formal-groupoid-completion}
Let
\[
    \mathcal H_\bullet\rightrightarrows X
\]
be any groupoid object in $\mathcal X_{/S}$.  Its unitwise completion
\[
    \widehat{\mathcal H}^{\,\mathrm{dR}}_\bullet
\]
is formal and admits a canonical anchor
\[
    \phi_{\widehat{\mathcal H}}:
    \widehat{\mathcal H}^{\,\mathrm{dR}}_\bullet
    \longrightarrow
    R_{X/S,\bullet}.
\]
The construction is functorial in fixed-object groupoid morphisms and
compatible with arbitrary base change in $S$.
\end{theorem}

\begin{proof}
There is a contractibly unique fixed-object simplicial morphism
\[
    \mathcal H_\bullet
    \longrightarrow
    \operatorname{Pair}(X/S)_\bullet
\]
given degreewise by the ordered tuple of vertex maps.  Applying unitwise
completion and Theorem~\ref{thm:pair-groupoid-completion} gives
\[
    \widehat{\mathcal H}^{\,\mathrm{dR}}_\bullet
    \longrightarrow
    \widehat{\operatorname{Pair}(X/S)}^{\,\mathrm{dR}}_\bullet
    \simeq
    R_{X/S,\bullet}.
\]
The source is formal by
Theorem~\ref{thm:unitwise-completion-idempotent}.  Functoriality follows from
functoriality of the map to the relative pair groupoid and of completion.
Base-change compatibility follows from
Proposition~\ref{prop:centred-completion-base-change}.
\end{proof}

\begin{proposition}[Coreflection of anchored groupoids]
\label{prop:anchored-formal-coreflection}
The canonical-anchor functor
\[
    \operatorname{can}_{X/S}:
    \operatorname{FormLieGpd}(X/S)
    \longrightarrow
    \operatorname{AnchGpd}(X/S)
\]
admits unitwise completion as a right adjoint
\[
    \widehat{(-)}^{\,\mathrm{dR}}:
    \operatorname{AnchGpd}(X/S)
    \longrightarrow
    \operatorname{FormLieGpd}(X/S).
\]
For an anchored groupoid
\[
    \phi_{\mathfrak G}:
    \mathfrak G_\bullet\to R_{X/S,\bullet},
\]
the canonical anchor of
$\widehat{\mathfrak G}^{\,\mathrm{dR}}_\bullet$ is equivalently represented by
either composite
\[
    \widehat{\mathfrak G}^{\,\mathrm{dR}}_\bullet
    \xrightarrow{\kappa_{\mathfrak G}}
    \mathfrak G_\bullet
    \xrightarrow{\phi_{\mathfrak G}}
    R_{X/S,\bullet}
\]
or
\[
    \widehat{\mathfrak G}^{\,\mathrm{dR}}_\bullet
    \xrightarrow{\widehat{\phi_{\mathfrak G}}^{\,\mathrm{dR}}}
    \widehat{R_{X/S,\bullet}}^{\,\mathrm{dR}}
    \xrightarrow{\kappa_{R}}
    R_{X/S,\bullet}.
\]
\end{proposition}

\begin{proof}
Naturality of the counit $\kappa$ gives
\[
    \kappa_R
    \widehat{\phi_{\mathfrak G}}^{\,\mathrm{dR}}
    \simeq
    \phi_{\mathfrak G}\kappa_{\mathfrak G}.
\]
Because $R_{X/S,\bullet}$ is formal, $\kappa_R$ is an equivalence, so the two
displayed anchors agree canonically.  The completed groupoid is formal by
Theorem~\ref{thm:unitwise-completion-idempotent}; its anchor is therefore the
canonical one by Theorem~\ref{thm:maximal-formal-terminal}.

Now let $\mathfrak F_\bullet$ be formal and let
$(\mathfrak G_\bullet,\phi_{\mathfrak G})$ be anchored.  Since
\[
    \operatorname{Map}_{\operatorname{Gpd}_X}
    \left(
        \mathfrak F_\bullet,
        R_{X/S,\bullet}
    \right)
    \simeq *,
\]
the anchor-compatibility condition on a morphism
$\mathfrak F_\bullet\to\mathfrak G_\bullet$ has a contractible space of
solutions.  Hence
\[
\begin{aligned}
    \operatorname{Map}_{\operatorname{AnchGpd}}
    \left(
        \operatorname{can}_{X/S}(\mathfrak F_\bullet),
        (\mathfrak G_\bullet,\phi_{\mathfrak G})
    \right)
    &\simeq
    \operatorname{Map}_{\operatorname{Gpd}_X}
    \left(
        \mathfrak F_\bullet,
        \mathfrak G_\bullet
    \right)\\
    &\simeq
    \operatorname{Map}_{\operatorname{Gpd}_X}
    \left(
        \mathfrak F_\bullet,
        \widehat{\mathfrak G}^{\,\mathrm{dR}}_\bullet
    \right)\\
    &\simeq
    \operatorname{Map}_{\operatorname{FormLieGpd}}
    \left(
        \mathfrak F_\bullet,
        \widehat{\mathfrak G}^{\,\mathrm{dR}}_\bullet
    \right),
\end{aligned}
\]
where the middle equivalence is
Theorem~\ref{thm:formalization-coreflection}.  This is the claimed adjunction.
\end{proof}

\subsection{Effective formal leaf spaces}

\begin{definition}[Effective formal leaf space]
\label{def:formal-leaf-space}
For
$\mathfrak G_\bullet\in\operatorname{FormLieGpd}(X/S)$ define
\[
    Q_{\mathfrak G}
    :=
    \lvert\mathfrak G_\bullet\rvert.
\]
Its quotient presentation is
\[
    q_{\mathfrak G}:
    X\twoheadrightarrow Q_{\mathfrak G}.
\]
Realizing the canonical anchor gives
\[
    r_{\mathfrak G}:
    Q_{\mathfrak G}
    \longrightarrow
    Q_{X/S}
\]
with
\[
    r_{\mathfrak G}q_{\mathfrak G}
    \simeq
    \epsilon_{X/S}.
\]
We call $Q_{\mathfrak G}$ the \emph{effective formal leaf space} or
\emph{effective formal quotient}.
\end{definition}

\subsection{Anchored quotient classification}

\begin{definition}[Infinitesimal quotient factorizations]
\label{def:infinitesimal-quotient-factorizations}
Let
\[
    \iota_{02}:\Delta^1\longrightarrow\Delta^2,
    \qquad
    0\longmapsto0,\quad 1\longmapsto2,
\]
be the long-edge inclusion.  Define the $\infty$-category of factorizations of
$\epsilon_{X/S}$ by the homotopy fibre
\[
    \operatorname{Fact}(\epsilon_{X/S})
    :=
    \operatorname{fib}_{\epsilon_{X/S}}
    \left(
        \iota_{02}^*:
        \operatorname{Fun}(\Delta^2,\mathcal E)
        \longrightarrow
        \operatorname{Fun}(\Delta^1,\mathcal E)
    \right).
\]
Thus its objects are coherently commutative factorizations
\[
    X\xrightarrow{q}Q\xrightarrow{r}Q_{X/S},
    \qquad
    rq\simeq\epsilon_{X/S},
\]
and all mapping-space coherence is inherited from this homotopy fibre.

Let
\[
    \operatorname{Fact}_{\mathrm{eff}}(\epsilon_{X/S})
    \subset
    \operatorname{Fact}(\epsilon_{X/S})
\]
be the full subcategory on the factorizations for which
$q:X\twoheadrightarrow Q$ is an effective epimorphism.
\end{definition}

\begin{lemma}[Effectivity of the second quotient map]
\label{lem:second-factor-effective}
Let
\[
    A\xrightarrow{f}B\xrightarrow{g}C
\]
be morphisms in an $\infty$-topos.  If $gf$ is an effective epimorphism, then
$g$ is an effective epimorphism.  Consequently every object of
$\operatorname{Fact}_{\mathrm{eff}}(\epsilon_{X/S})$ is canonically a chain
\[
    X\twoheadrightarrow Q\twoheadrightarrow Q_{X/S}
\]
of effective epimorphisms.
\end{lemma}

\begin{proof}
Factor $g$ by the effective-epimorphism--monomorphism factorization
\[
    B\twoheadrightarrow I_g\lhook\joinrel\longrightarrow C.
\]
The composite $gf$ factors through the monomorphism $I_g\to C$.  Since the
effective image of the effective epimorphism $gf$ is all of $C$, this
monomorphism is an equivalence.  Hence $g$ is an effective epimorphism.
Apply this to
\[
    X\xrightarrow{q}Q\xrightarrow{r}Q_{X/S},
    \qquad
    rq\simeq\epsilon_{X/S}.
\]
\end{proof}

\begin{theorem}[Anchored Quotient Classification Theorem]
\label{thm:anchored-gpd-quotient-classification}
Realization and \v{C}ech nerve induce inverse equivalences
\[
    \operatorname{AnchGpd}(X/S)
    \simeq
    \operatorname{Fact}_{\mathrm{eff}}(\epsilon_{X/S}).
\]
Under this equivalence,
\[
    \mathfrak G_\bullet
    \longleftrightarrow
    \left(
        X\twoheadrightarrow Q_{\mathfrak G}
        \twoheadrightarrow Q_{X/S}
    \right),
\]
and
\[
    \mathfrak G_n
    \simeq
    \underbrace{
    X\times_{Q_{\mathfrak G}}\cdots
    \times_{Q_{\mathfrak G}}X
    }_{n+1}.
\]
\end{theorem}

\begin{proof}
For an anchored groupoid, realization of the anchor gives
\[
    r_{\mathfrak G}:
    Q_{\mathfrak G}
    \longrightarrow
    Q_{X/S}.
\]
Because the anchor is the identity in degree zero,
\[
    r_{\mathfrak G}q_{\mathfrak G}
    \simeq
    \epsilon_{X/S}.
\]
The map $q_{\mathfrak G}$ is effective by groupoid effectivity, and
Lemma~\ref{lem:second-factor-effective} therefore makes
$r_{\mathfrak G}$ effective as well.  This produces an object of
$\operatorname{Fact}_{\mathrm{eff}}(\epsilon_{X/S})$.

Conversely, let
\[
    X\xrightarrow{q}Q\xrightarrow{r}Q_{X/S}
\]
represent an object of
$\operatorname{Fact}_{\mathrm{eff}}(\epsilon_{X/S})$.  The \v{C}ech nerve
$\check C_\bullet(q)$ is a fixed-object groupoid.  The identity maps on the
$X$-factors and the specified $2$-simplex
$rq\simeq\epsilon_{X/S}$ give degreewise maps
\[
    \underbrace{X\times_Q\cdots\times_QX}_{n+1}
    \longrightarrow
    \underbrace{
    X\times_{Q_{X/S}}\cdots\times_{Q_{X/S}}X
    }_{n+1},
\]
compatible with all simplicial operators.  These define the anchor.

Groupoid effectivity identifies a fixed-object groupoid with the \v{C}ech
nerve of its quotient presentation and an effective epimorphism with the
realization of its \v{C}ech nerve.  Since the factorization category was
defined as a homotopy fibre, the same constructions retain the complete
mapping-space coherence.  Realization and \v{C}ech nerve are therefore inverse
equivalences.
\end{proof}

\subsection{Quotient-side formalization and recognition}

\begin{definition}[Centred quotient completion]
\label{def:quotient-formalization}
Let
\[
    q:X\twoheadrightarrow Q
\]
be an effective epimorphism.  Its centred completion is
\[
    P_q
    :=
    C_X(Q,q)
    =
    Q\times_{Q_{\mathrm{dR}/S}}X_{\mathrm{dR}/S}.
\]
The completed centre is
\[
    \delta_q
    :=
    (q,\eta_{X/S}):
    X\longrightarrow P_q.
\]
Factor $\delta_q$ by the effective-epimorphism--monomorphism factorization:
\[
    X
    \xrightarrow{\widehat q}
    \widehat Q_q
    \xrightarrow{i_q}
    P_q.
\]
Let $p_q:P_q\to Q$ be the first projection and put
\[
    c_q
    :=
    p_qi_q:
    \widehat Q_q\longrightarrow Q.
\]
Then
\[
    c_q\widehat q\simeq q.
\]
\end{definition}

\begin{lemma}[The quotient comparison is effective]
\label{lem:quotient-formalization-effective}
The morphism
\[
    c_q:\widehat Q_q\longrightarrow Q
\]
is an effective epimorphism.
\end{lemma}

\begin{proof}
The composite
\[
    X\xrightarrow{\widehat q}\widehat Q_q
      \xrightarrow{c_q}Q
\]
is the effective epimorphism $q$.  The image of the composite factors through
the image of $c_q$, so the image of $c_q$ is all of $Q$.
\end{proof}

\begin{theorem}[Quotient formula for unitwise completion]
\label{thm:quotient-formula-unitwise-completion}
For an effective epimorphism $q:X\twoheadrightarrow Q$ there is a canonical
equivalence
\[
    \widehat{\check C_\bullet(q)}^{\,\mathrm{dR}}
    \simeq
    \check C_\bullet(\widehat q).
\]
Under this equivalence the realization of the completion counit is
\[
    c_q:\widehat Q_q\longrightarrow Q.
\]
\end{theorem}

\begin{proof}
In degree $n$,
\[
    \check C_n(q)
    =
    \underbrace{X\times_Q\cdots\times_QX}_{n+1}
\]
with diagonal centre $X\to\check C_n(q)$.  Left exactness gives
\[
    L_S(\check C_n(q))
    \simeq
    \underbrace{
    L_S(X)\times_{L_S(Q)}\cdots\times_{L_S(Q)}L_S(X)
    }_{n+1},
\]
and the localized centre is the full diagonal from $L_S(X)$.  Substituting
this into the defining pullback for centred completion and regrouping finite
limits gives
\[
\begin{aligned}
    C_X(\check C_n(q))
    &\simeq
    \underbrace{
    X\times_{P_q}\cdots\times_{P_q}X
    }_{n+1},
\end{aligned}
\]
where
\[
    P_q
    =
    Q\times_{L_S(Q)}L_S(X)
\]
and each map $X\to P_q$ is the completed centre $\delta_q$.

Because $i_q:\widehat Q_q\to P_q$ is a monomorphism,
\[
    \widehat Q_q
    \simeq
    \widehat Q_q\times_{P_q}\widehat Q_q.
\]
Since $\delta_q=i_q\widehat q$, finite-limit pasting therefore gives
\[
    \underbrace{X\times_{P_q}\cdots\times_{P_q}X}_{n+1}
    \simeq
    \underbrace{
    X\times_{\widehat Q_q}\cdots
    \times_{\widehat Q_q}X
    }_{n+1}.
\]
The equivalences are compatible with deletion and repetition of factors, so
they assemble simplicially.  Since $\widehat q$ is effective, the right-hand
\v{C}ech nerve realizes to $\widehat Q_q$.  The counit is induced by
$c_q\widehat q\simeq q$, hence realizes to $c_q$.
\end{proof}

\begin{theorem}[Recognition of formal quotients]
\label{thm:formal-quotient-monomorphism-recognition}
For an effective epimorphism
\[
    q:X\twoheadrightarrow Q,
\]
the following are equivalent:
\begin{enumerate}
    \item $\check C_\bullet(q)$ is a formal fixed-object groupoid;
    \item
    \[
        q_{\mathrm{dR}/S}:
        X_{\mathrm{dR}/S}
        \longrightarrow
        Q_{\mathrm{dR}/S}
    \]
    is a monomorphism;
    \item $q_{\mathrm{dR}/S}$ is an equivalence;
    \item
    \[
        c_q:\widehat Q_q\longrightarrow Q
    \]
    is an equivalence.
\end{enumerate}
\end{theorem}

\begin{proof}
The localized degree-one unit of $\check C_\bullet(q)$ is the diagonal
\[
    \Delta_{q_{\mathrm{dR}/S}}:
    X_{\mathrm{dR}/S}
    \longrightarrow
    X_{\mathrm{dR}/S}
    \times_{Q_{\mathrm{dR}/S}}
    X_{\mathrm{dR}/S}.
\]
By Proposition~\ref{prop:degree-one-formality}, the \v{C}ech groupoid is formal
if and only if this diagonal is an equivalence.  A morphism in an
$\infty$-topos is a monomorphism precisely when its diagonal is an
equivalence.  This proves the equivalence of (1) and (2).

The preceding de Rham chapter separately proves that relative de Rham
localization carries effective epimorphisms in the slice to effective
epimorphisms in the ambient slice; this conclusion is not being inferred from
left exactness alone.  Since $q$ is an effective epimorphism,
$q_{\mathrm{dR}/S}$ is therefore an effective epimorphism.  In an
$\infty$-topos a morphism which is both an effective epimorphism and a
monomorphism is an equivalence.  Hence (2) and (3) are equivalent.

By Theorem~\ref{thm:quotient-formula-unitwise-completion}, the completion
counit realizes to $c_q$.  A morphism of effective groupoid presentations is
an equivalence if and only if its realization is an equivalence, proving the
equivalence of (1) and (4).
\end{proof}

\begin{definition}[Formal effective quotients]
\label{def:formal-effective-quotients}
Let
\[
    \operatorname{Epi}^{\mathrm{for}}(X)
    \subset
    \operatorname{Epi}(X)
\]
be the full subcategory spanned by effective epimorphisms
$q:X\twoheadrightarrow Q$ for which
$q_{\mathrm{dR}/S}$ is a monomorphism, equivalently an equivalence.
\end{definition}

\begin{theorem}[Quotient formalization coreflection]
\label{thm:quotient-formalization-coreflection}
The inclusion
\[
    \operatorname{Epi}^{\mathrm{for}}(X)
    \hookrightarrow
    \operatorname{Epi}(X)
\]
admits
\[
    q\longmapsto\widehat q
\]
as a right adjoint.  Its counit at $q$ is
\[
    c_q:\widehat Q_q\to Q.
\]
\end{theorem}

\begin{proof}
Transport the groupoid formalization coreflection
of Theorem~\ref{thm:formalization-coreflection} through the equivalence
\[
    \operatorname{Gpd}_X(\mathcal X_{/S})
    \simeq
    \operatorname{Epi}(X)
\]
of Proposition~\ref{prop:groupoids-effective-quotients-equivalence}.
Theorem~\ref{thm:quotient-formula-unitwise-completion} identifies the image
of $q$ under the transported right adjoint with $\widehat q$, and identifies
the transported counit with $c_q$.  The fixed points are
$\operatorname{Epi}^{\mathrm{for}}(X)$ by
Theorem~\ref{thm:formal-quotient-monomorphism-recognition}.
\end{proof}

\begin{proposition}[Canonical factorization through the maximal quotient]
\label{prop:formal-quotient-canonical-anchor}
Every formal effective quotient
\[
    q:X\twoheadrightarrow Q
\]
admits a contractibly unique morphism
\[
    r_q:Q\longrightarrow Q_{X/S}
\]
such that
\[
    r_qq\simeq\epsilon_{X/S}.
\]
The morphism $r_q$ is an effective epimorphism and a relative de Rham
equivalence:
\[
    L_S(r_q):
    L_S(Q)\xrightarrow{\sim}L_S(Q_{X/S}).
\]
\end{proposition}

\begin{proof}
The \v{C}ech groupoid $\check C_\bullet(q)$ is formal, so it has a
contractibly unique canonical anchor
\[
    \check C_\bullet(q)\longrightarrow R_{X/S,\bullet}.
\]
Realization gives $r_q$, and the degree-zero identity gives
$r_qq\simeq\epsilon_{X/S}$.  Conversely, a map $r_q$ with this property
determines the corresponding anchored \v{C}ech morphism, so uniqueness of the
anchor gives contractible uniqueness of $r_q$.

The composite $r_qq=\epsilon_{X/S}$ is effective, hence
Lemma~\ref{lem:second-factor-effective} makes $r_q$ effective.

By Theorem~\ref{thm:formal-quotient-monomorphism-recognition},
formality of $\check C_\bullet(q)$ gives an equivalence
\[
    L_S(q):L_S(X)\xrightarrow{\sim}L_S(Q).
\]
Lemma~\ref{lem:effective-quotient-dr-equivalence} gives independently
\[
    L_S(\epsilon_{X/S}):
    L_S(X)\xrightarrow{\sim}L_S(Q_{X/S}).
\]
Applying $L_S$ to $r_qq\simeq\epsilon_{X/S}$ yields
\[
    L_S(r_q)L_S(q)\simeq L_S(\epsilon_{X/S}).
\]
Thus
\[
    L_S(r_q)
    \simeq
    L_S(\epsilon_{X/S})L_S(q)^{-1}
\]
is an equivalence.
\end{proof}

\begin{definition}[Formal infinitesimal quotient factorizations]
\label{def:formal-infinitesimal-quotient-factorizations}
Let
\[
    \operatorname{Fact}^{\mathrm{for}}_{\mathrm{eff}}
    (\epsilon_{X/S})
    \subset
    \operatorname{Fact}_{\mathrm{eff}}(\epsilon_{X/S})
\]
be the full subcategory on factorizations
\[
    X\xrightarrow{q}Q\xrightarrow{r}Q_{X/S}
\]
for which $q_{\mathrm{dR}/S}$ is a monomorphism, equivalently an
equivalence.  Equivalently, $c_q:\widehat Q_q\to Q$ is an equivalence.
For every such factorization both arrows are effective epimorphisms by
Lemma~\ref{lem:second-factor-effective}, and
Proposition~\ref{prop:formal-quotient-canonical-anchor} shows that the second
arrow is also a relative de Rham equivalence.
\end{definition}

\begin{theorem}[Formal Quotient Classification Theorem]
\label{thm:formal-gpd-quotient-classification}
There are canonical equivalences
\[
    \operatorname{FormLieGpd}(X/S)
    \simeq
    \operatorname{Gpd}^{\mathrm{for}}_X(\mathcal X_{/S})
    \simeq
    \operatorname{Epi}^{\mathrm{for}}(X)
    \simeq
    \operatorname{Fact}^{\mathrm{for}}_{\mathrm{eff}}
    (\epsilon_{X/S}).
\]
Thus a brave-new formal Lie groupoid is equivalently an effective quotient
\[
    q:X\twoheadrightarrow Q
\]
such that
\[
    q_{\mathrm{dR}/S}:
    X_{\mathrm{dR}/S}\to Q_{\mathrm{dR}/S}
\]
is monic, equivalently an equivalence.  The map
$Q\to Q_{X/S}$ is then recovered canonically and is both an effective
epimorphism and a relative de Rham equivalence.
\end{theorem}

\begin{proof}
The first equivalence is the definition
\[
    \operatorname{FormLieGpd}(X/S)
    =
    \operatorname{Gpd}^{\mathrm{for}}_X(\mathcal X_{/S}).
\]
The second is the fixed-object groupoid/effective quotient equivalence
restricted by
Theorem~\ref{thm:formal-quotient-monomorphism-recognition}.

It remains to compare formal effective quotients with formal infinitesimal
factorizations on mapping spaces.  For a formal effective quotient
$q:X\twoheadrightarrow Q$, let
\[
    \mathcal A(q)
    :=
    \left\{
        r:Q\to Q_{X/S}
        \ \middle|\ 
        rq\simeq\epsilon_{X/S}
    \right\}
\]
denote the corresponding homotopy fibre of mapping spaces.  Passing between a
map $r$ in this fibre and the induced morphism of \v{C}ech nerves identifies
$\mathcal A(q)$ with
\[
    \operatorname{Map}_{\operatorname{Gpd}_X}
    \left(
        \check C_\bullet(q),
        R_{X/S,\bullet}
    \right).
\]
The groupoid $\check C_\bullet(q)$ is formal, and
$R_{X/S,\bullet}$ is terminal among formal fixed-object groupoids.  Hence
\[
    \mathcal A(q)\simeq *.
\]
This recovers Proposition~\ref{prop:formal-quotient-canonical-anchor} and also
controls morphisms.

Indeed, let
\[
    a:Q\longrightarrow Q'
\]
be a morphism under $X$ between formal effective quotients, so that
$aq\simeq q'$.  Then
\[
    r_{q'}a,\ r_q:Q\rightrightarrows Q_{X/S}
\]
both lie in $\mathcal A(q)$, because
\[
    (r_{q'}a)q
    \simeq
    r_{q'}q'
    \simeq
    \epsilon_{X/S}
    \simeq
    r_qq.
\]
Since $\mathcal A(q)$ is contractible, the space of compatible homotopies
\[
    r_{q'}a\simeq r_q
\]
is contractible.  Thus every morphism in
$\operatorname{Epi}^{\mathrm{for}}(X)$ admits a contractibly unique enhancement
to a morphism in
$\operatorname{Fact}^{\mathrm{for}}_{\mathrm{eff}}(\epsilon_{X/S})$.
The forgetful functor from the factorization category is therefore fully
faithful and essentially surjective, proving the final equivalence.
\end{proof}

\begin{corollary}[Initial and terminal formal Lie groupoids]
\label{cor:initial-terminal-formal-gpd}
The category $\operatorname{FormLieGpd}(X/S)$ has canonical initial and
terminal objects
\[
    0_{X,\bullet}
    \longrightarrow
    \mathfrak G_\bullet
    \longrightarrow
    R_{X/S,\bullet}.
\]
Their quotient factorizations are
\[
    X\xrightarrow{\operatorname{id}}X
      \xrightarrow{\epsilon_{X/S}}Q_{X/S}
\]
and
\[
    X\xrightarrow{\epsilon_{X/S}}Q_{X/S}
      \xrightarrow{\operatorname{id}}Q_{X/S}.
\]
\end{corollary}

\begin{proof}
Apply Proposition~\ref{prop:extremal-infinitesimal-groupoids} and
Definition~\ref{def:formal-lie-groupoid}.
\end{proof}

\begin{remark}[Nonlinear anchor]
\label{rem:nonlinear-anchor}
Before formalization, an anchor is genuine additional structure.  After
formalization it is the contractibly unique morphism to the maximal formal
relation.  On quotients it becomes
\[
    r_{\mathfrak G}:Q_{\mathfrak G}\to Q_{X/S}.
\]
The final handoff section below records the represented first-order
linearization of the canonical nonlinear anchor in the geometric range where
the manuscript has constructed the necessary cotangent objects.  No bracket,
Lie-algebraic presentation, or recognition theorem is introduced in the
intrinsic formal-groupoid package.
\end{remark}


\section{Formal orbit disks and formal leaves}
\label{sec:flg-formal-disks}

Fix
\[
    \mathfrak G_\bullet
    \in
    \operatorname{FormLieGpd}(X/S)
\]
and write
\[
    q:=q_{\mathfrak G}:X\twoheadrightarrow Q_{\mathfrak G},
    \qquad
    r:=r_{\mathfrak G}:Q_{\mathfrak G}\twoheadrightarrow Q_{X/S}.
\]
By Theorem~\ref{thm:formal-quotient-monomorphism-recognition},
\[
    q_{\mathrm{dR}/S}:
    X_{\mathrm{dR}/S}
    \longrightarrow
    (Q_{\mathfrak G})_{\mathrm{dR}/S}
\]
is an equivalence, and in particular a monomorphism.

\begin{definition}[Formal orbit disk]
\label{def:formal-orbit-disk}
Let $a:T\to X$ be an object of $\mathcal X_{/X}$.  Define the
$\mathfrak G$-formal orbit disk centred at $a$ by
\[
    \mathbb D^{\mathfrak G}_a
    :=
    T\times_{Q_{\mathfrak G}}X.
\]
Write
\[
    d^{\mathfrak G}_a:
    \mathbb D^{\mathfrak G}_a\to T
\]
for the first projection and
\[
    \operatorname{ev}^{\mathfrak G}_a:
    \mathbb D^{\mathfrak G}_a\to X
\]
for the second.
\end{definition}

\begin{proposition}[Arrow-fibre presentation]
\label{prop:formal-disk-arrow-fibre}
There is a canonical equivalence
\[
    \mathbb D^{\mathfrak G}_a
    \simeq
    T\times_{a,X,s}\mathfrak G_1,
\]
under which $d^{\mathfrak G}_a$ is the pullback of the source and
$\operatorname{ev}^{\mathfrak G}_a$ is the pulled-back target.
\end{proposition}

\begin{proof}
Groupoid effectivity gives
\[
    \mathfrak G_1
    \simeq
    X\times_{Q_{\mathfrak G}}X
\]
with source and target the two projections.  Pulling source back along
$a:T\to X$ yields the claimed equivalence.
\end{proof}

The groupoid unit induces a centre section
\[
    e^{\mathfrak G}_a:
    T\longrightarrow
    \mathbb D^{\mathfrak G}_a.
\]

\begin{proposition}[Intrinsic completeness of formal orbit disks]
\label{prop:formal-disk-unitwise-complete}
The formal orbit disk is relative de Rham-complete along its centre:
\[
    \mathbb D^{\mathfrak G}_a
    \times_{
        (\mathbb D^{\mathfrak G}_a)_{\mathrm{dR}/S}
    }
    T_{\mathrm{dR}/S}
    \xrightarrow{\sim}
    \mathbb D^{\mathfrak G}_a.
\]
Equivalently,
\[
    (e^{\mathfrak G}_a)_{\mathrm{dR}/S}:
    T_{\mathrm{dR}/S}
    \xrightarrow{\sim}
    (\mathbb D^{\mathfrak G}_a)_{\mathrm{dR}/S}
\]
is an equivalence.
\end{proposition}

\begin{proof}
Left exactness gives
\[
    (\mathbb D^{\mathfrak G}_a)_{\mathrm{dR}/S}
    \simeq
    T_{\mathrm{dR}/S}
    \times_{(Q_{\mathfrak G})_{\mathrm{dR}/S}}
    X_{\mathrm{dR}/S}.
\]
By Theorem~\ref{thm:formal-quotient-monomorphism-recognition},
\[
    q_{\mathrm{dR}/S}:
    X_{\mathrm{dR}/S}
    \xrightarrow{\sim}
    (Q_{\mathfrak G})_{\mathrm{dR}/S}
\]
is an equivalence.  The projection
\[
    T_{\mathrm{dR}/S}
    \times_{(Q_{\mathfrak G})_{\mathrm{dR}/S}}
    X_{\mathrm{dR}/S}
    \longrightarrow
    T_{\mathrm{dR}/S}
\]
is therefore its base change along
$(qa)_{\mathrm{dR}/S}:T_{\mathrm{dR}/S}\to
(Q_{\mathfrak G})_{\mathrm{dR}/S}$ and is an equivalence.  Its inverse is the
localized centre section
\[
    (e^{\mathfrak G}_a)_{\mathrm{dR}/S}:
    T_{\mathrm{dR}/S}
    \xrightarrow{\sim}
    (\mathbb D^{\mathfrak G}_a)_{\mathrm{dR}/S}.
\]
Apply Proposition~\ref{prop:centred-completion-recognition}.
\end{proof}

\begin{proposition}[Effectivity of formal orbit disks]
\label{prop:formal-disk-effective}
The projection
\[
    d^{\mathfrak G}_a:
    \mathbb D^{\mathfrak G}_a\twoheadrightarrow T
\]
is an effective epimorphism.
\end{proposition}

\begin{proof}
It is the base change of
$q:X\twoheadrightarrow Q_{\mathfrak G}$
along $qa:T\to Q_{\mathfrak G}$.
\end{proof}

\begin{proposition}[Comparison with the ambient relative de Rham disk]
\label{prop:formal-disk-ambient-comparison}
The quotient morphism
\[
    r:Q_{\mathfrak G}\to Q_{X/S}
\]
induces a canonical morphism over $T$ and $X$
\[
    \chi^{\mathrm{disk}}_{\mathfrak G,a}:
    \mathbb D^{\mathfrak G}_a
    \longrightarrow
    \mathbb D^\infty_{X/S,a}
    \simeq
    T\times_{Q_{X/S}}X.
\]
These maps are natural in $a$ and in $\mathfrak G$.
\end{proposition}

\begin{proof}
Apply $r$ to the base of the defining fibre product
$T\times_{Q_{\mathfrak G}}X$.  Naturality is functoriality of pullback.
\end{proof}

\begin{remark}[Formal leaves]
The fibre of
\[
    q:X\twoheadrightarrow Q_{\mathfrak G}
\]
through $a:T\to X$ is $\mathbb D^{\mathfrak G}_a$.  Thus
$Q_{\mathfrak G}$ parametrizes the effective formal leaves, and every such
leaf is intrinsically complete along its chosen centre.  This is geometrically
parallel to the formal-leaf viewpoint in formal integration of partition Lie
algebroids, while the algebraic input and coefficient categories there are
different \cite{BrantnerMagidsonNuiten2025}.
\end{remark}

\begin{proposition}[Cartesianity of formal disks]
\label{prop:formal-disk-cartesian}
For $b:T'\to T$ and $a:T\to X$ there is a canonical equivalence
\[
    T'\times_T\mathbb D^{\mathfrak G}_a
    \simeq
    \mathbb D^{\mathfrak G}_{ab},
\]
compatible with evaluation, centre sections, identities, composition, and
higher pullback pasting.
\end{proposition}

\begin{proof}
Both sides are canonically
\[
    T'\times_{Q_{\mathfrak G}}X.
\]
\end{proof}


\section{Raw and formal isotropy}
\label{sec:flg-isotropy}

A general internal groupoid carries a raw automorphism object over its object
space.  In an $\infty$-topos this raw object also sees ambient higher inertia
of the object space itself.  For the formal theory we therefore distinguish it
from its centred de Rham completion.

\subsection{Raw isotropy and inertia}

In this subsection let
\[
    \mathfrak G_\bullet\in
    \operatorname{Gpd}_X(\mathcal X_{/S})
\]
be arbitrary; no formality is used in the raw isotropy constructions.

\begin{definition}[Raw isotropy group]
\label{def:raw-isotropy-group}
Let
\[
    (s,t):
    \mathfrak G_1\longrightarrow X\times_SX
\]
be the source-target map.  Define
\[
    I^{\mathrm{raw}}_{\mathfrak G}
    :=
    X\times_{\Delta_X,\;X\times_SX,\;(s,t)}
    \mathfrak G_1.
\]
Write
\[
    p_{\mathrm{raw}}:
    I^{\mathrm{raw}}_{\mathfrak G}\to X
\]
for its structural morphism and
\[
    e_{\mathrm{raw}}:
    X\longrightarrow I^{\mathrm{raw}}_{\mathfrak G}
\]
for the section induced by the groupoid unit.
\end{definition}

\begin{proposition}[Raw isotropy is a group object]
\label{prop:raw-isotropy-group}
The object
$I^{\mathrm{raw}}_{\mathfrak G}\to X$
carries a canonical group-object structure in $\mathcal X_{/X}$.
Multiplication, unit, and inverse are the restrictions of groupoid
composition, unit, and inversion.
\end{proposition}

\begin{proof}
The composable-pair object
\[
    I^{\mathrm{raw}}_{\mathfrak G}
    \times_X
    I^{\mathrm{raw}}_{\mathfrak G}
\]
maps canonically to
$\mathfrak G_1\times_{t,X,s}\mathfrak G_1$, and groupoid composition carries
it back to the source-target diagonal fibre.  Thus composition restricts to
multiplication on $I^{\mathrm{raw}}_{\mathfrak G}$.  Inversion preserves the
same fibre, and the groupoid unit factors through it.  The group identities
and all higher coherences are the corresponding restrictions of the
simplicial groupoid identities.
\end{proof}

\begin{definition}[Relative inertia object]
\label{def:relative-inertia-object}
For an object $Q\to S$, regard the diagonal
\[
    \Delta_Q:Q\longrightarrow Q\times_SQ
\]
as a pointed object of the slice $\mathcal X_{/Q}$ through the first
projection.  Define its relative inertia object by
\[
    \mathcal I_{Q/S}
    :=
    Q\times_{\Delta_Q,\;Q\times_SQ,\;\Delta_Q}Q.
\]
\end{definition}

\begin{proposition}[Canonical group structure on relative inertia]
\label{prop:relative-inertia-group}
The object
\[
    \mathcal I_{Q/S}\longrightarrow Q
\]
carries a canonical group-object structure in $\mathcal X_{/Q}$.  For every
morphism $f:Q\to Q'$ over $S$, functoriality supplies a canonical homomorphism
of group objects over $Q$
\[
    \mathcal I_{Q/S}
    \longrightarrow
    f^*\mathcal I_{Q'/S}.
\]
These morphisms are functorial in $f$ and compatible with arbitrary base
change in $S$.
\end{proposition}

\begin{proof}
Form in the slice $\mathcal X_{/Q}$ the \v{C}ech nerve of the section
\[
    \Delta_Q:
    Q\longrightarrow
    (Q\times_SQ\xrightarrow{\operatorname{pr}_1}Q).
\]
Its degree-zero object is the terminal object $Q$ of the slice and its
degree-one object is
\[
    Q\times_{Q\times_SQ}Q
    =
    \mathcal I_{Q/S}.
\]
Every \v{C}ech nerve is a groupoid object.  A groupoid object whose object
space is terminal is equivalently a group object, so the unit, multiplication,
and inverse of $\mathcal I_{Q/S}$ are the degree-zero, degree-two, and inversion
maps of this \v{C}ech groupoid.  Functoriality follows from functoriality of
diagonals and finite limits.  Pullback in an $\infty$-topos preserves the
entire defining \v{C}ech nerve, giving the asserted base-change compatibility.
\end{proof}

\begin{proposition}[Raw isotropy as pulled-back inertia]
\label{prop:isotropy-inertia}
There is a canonical equivalence of group objects over $X$
\[
    I^{\mathrm{raw}}_{\mathfrak G}
    \simeq
    X\times_{Q_{\mathfrak G}}
    \mathcal I_{Q_{\mathfrak G}/S}.
\]
\end{proposition}

\begin{proof}
Using
\[
    \mathfrak G_1
    \simeq
    X\times_{Q_{\mathfrak G}}X,
\]
substitute the effective relation presentation into the defining pullback for
$I^{\mathrm{raw}}_{\mathfrak G}$.  Finite-limit pasting identifies the
result with the pullback to $X$ of
\[
    Q_{\mathfrak G}
    \times_{Q_{\mathfrak G}\times_SQ_{\mathfrak G}}
    Q_{\mathfrak G}.
\]
Composition and inversion agree with loop concatenation and reversal, so the
equivalence respects the group structures.
\end{proof}

\subsection{Formal isotropy as completed anchor kernel}

From now through the formal-isotropy constructions, let
\[
    \mathfrak G_\bullet\in\operatorname{FormLieGpd}(X/S).
\]

\begin{definition}[Formal isotropy group]
\label{def:isotropy-formal-group}
Define the \emph{formal isotropy group} by centred completion along the raw
identity section:
\[
    I^{\mathrm{for}}_{\mathfrak G}
    :=
    C_X
    \left(
        I^{\mathrm{raw}}_{\mathfrak G},
        e_{\mathrm{raw}}
    \right)
    =
    I^{\mathrm{raw}}_{\mathfrak G}
    \times_{
        (I^{\mathrm{raw}}_{\mathfrak G})_{\mathrm{dR}/S}
    }
    X_{\mathrm{dR}/S}.
\]
Write
\[
    \kappa_I:
    I^{\mathrm{for}}_{\mathfrak G}
    \longrightarrow
    I^{\mathrm{raw}}_{\mathfrak G}
\]
for the completion counit.
\end{definition}

\begin{theorem}[Formal isotropy group theorem]
\label{thm:isotropy-group}
The object
\[
    I^{\mathrm{for}}_{\mathfrak G}\to X
\]
is canonically a group object in $\mathcal X_{/X}$, the morphism
$\kappa_I$ is a homomorphism, and its identity section is relative
de Rham-complete.  Moreover there is a canonical equivalence of group objects
over $X$
\[
    I^{\mathrm{for}}_{\mathfrak G}
    \xrightarrow{\sim}
    \mathfrak G_1
    \times_{R_{X/S}}
    X,
\]
where $\mathfrak G_1\to R_{X/S}$ is the canonical degree-one anchor and
$X\to R_{X/S}$ is the unit section.
\end{theorem}

\begin{proof}
The raw multiplication is a morphism of centred objects
\[
    I^{\mathrm{raw}}_{\mathfrak G}
    \times_X
    I^{\mathrm{raw}}_{\mathfrak G}
    \longrightarrow
    I^{\mathrm{raw}}_{\mathfrak G},
\]
with centre the product of identity sections.  Since $C_X$ preserves pullbacks and fixes
$(X,\operatorname{id}_X)$,
\[
    C_X
    \left(
        I^{\mathrm{raw}}_{\mathfrak G}
        \times_X
        I^{\mathrm{raw}}_{\mathfrak G}
    \right)
    \simeq
    I^{\mathrm{for}}_{\mathfrak G}
    \times_X
    I^{\mathrm{for}}_{\mathfrak G}.
\]
Applying $C_X$ to multiplication, unit, and inversion therefore gives a
group-object structure on $I^{\mathrm{for}}_{\mathfrak G}$.  Naturality of
$\kappa$ makes $\kappa_I$ a homomorphism.  Completeness of the identity
section is Theorem~\ref{thm:centred-completion-idempotent}.

For the kernel formula, the raw isotropy square is a pullback diagram of
centred objects:
\[
    I^{\mathrm{raw}}_{\mathfrak G}
    \simeq
    \mathfrak G_1
    \times_{X\times_SX}
    X,
\]
where the centres are respectively the groupoid unit, the diagonal, and
$\operatorname{id}_X$.  Apply $C_X$.  Since $\mathfrak G$ is formal,
\[
    C_X(\mathfrak G_1)\simeq\mathfrak G_1.
\]
By Theorem~\ref{thm:pair-groupoid-completion},
\[
    C_X(X\times_SX,\Delta_X)\simeq R_{X/S},
\]
and
\[
    C_X(X,\operatorname{id}_X)\simeq X.
\]
Pullback preservation therefore yields
\[
    I^{\mathrm{for}}_{\mathfrak G}
    \simeq
    \mathfrak G_1\times_{R_{X/S}}X.
\]
The completed map $\mathfrak G_1\to R_{X/S}$ is the canonical anchor by
Corollary~\ref{cor:canonical-formal-anchor}.  The multiplication on raw
isotropy is the restriction of groupoid composition, and its domain is the
pullback
\[
    I^{\mathrm{raw}}_{\mathfrak G}
    \times_X
    I^{\mathrm{raw}}_{\mathfrak G}.
\]
Since $C_X$ preserves this pullback and the completed multiplication is
obtained by applying $C_X$ to the same composition map, the displayed
anchor-kernel equivalence intertwines multiplication, unit, and inversion.
It is therefore an equivalence of group objects.
\end{proof}

\begin{remark}[Why raw and formal isotropy are distinct]
\label{rem:raw-formal-isotropy-distinct}
The formality of $\mathfrak G_\bullet$ does not imply, from left exactness
alone, that $I^{\mathrm{raw}}_{\mathfrak G}$ is complete along its identity.
This is a genuine categorical issue rather than a proof artifact.

For example, in the $\infty$-topos of spaces take the accessible left-exact
localization $L=\tau_{\le1}$ and
\[
    X=B\mathbb Z\times K(\mathbb Z,2),
    \qquad
    LX\simeq B\mathbb Z.
\]
The maximal $L$-formal relation
\[
    X\times_{B\mathbb Z}X\rightrightarrows X
\]
is formal.  Since
\[
    \mathcal I_{B\mathbb Z}
    \simeq
    B\mathbb Z\times\mathbb Z,
\]
its raw isotropy is
\[
    I^{\mathrm{raw}}
    \simeq
    X\times\mathbb Z.
\]
The identity section lies in the zero component, and centred completion gives
\[
    C_X(I^{\mathrm{raw}})
    \simeq
    (X\times\mathbb Z)
    \times_{B\mathbb Z\times\mathbb Z}
    B\mathbb Z
    \simeq X,
\]
which is not equivalent to $X\times\mathbb Z$.  Thus the separate formal
isotropy construction is essential at the level of general left-exact
localizations.  The present spectral de Rham package contains no theorem collapsing this
distinction.
\end{remark}

\subsection{The descended formal inertia kernel}

Put
\[
    Q:=Q_{\mathfrak G},
    \qquad
    Q^\infty:=Q_{X/S},
    \qquad
    r:=r_{\mathfrak G}:Q\to Q^\infty.
\]
Functoriality of relative inertia gives a homomorphism of group objects over
$Q$
\[
    \lambda_r:
    \mathcal I_{Q/S}
    \longrightarrow
    r^*\mathcal I_{Q^\infty/S}.
\]

\begin{definition}[Formal inertia kernel and formal adjoint group]
\label{def:formal-inertia-kernel}
Define
\[
    \mathcal K_{\mathfrak G}
    :=
    \mathcal I_{Q/S}
    \times_{r^*\mathcal I_{Q^\infty/S}}
    Q,
\]
where
\[
    Q\longrightarrow r^*\mathcal I_{Q^\infty/S}
\]
is the identity-loop section.  This is a group object over $Q$.  Define the
\emph{formal adjoint group} by the constructed inertia kernel
\[
    \operatorname{Ad}^{\mathrm{for}}(\mathfrak G)
    :=
    \mathcal K_{\mathfrak G}
    \longrightarrow Q.
\]
No independent descent object or adjoint-group datum is introduced.
\end{definition}

\begin{lemma}[Effective descent for internal group objects]
\label{lem:effective-descent-group-objects}
Let
\[
    e:U\twoheadrightarrow V
\]
be an effective epimorphism in $\mathcal E$, and let
$U_\bullet=\check C_\bullet(e)$.  Pullback induces a canonical equivalence
\[
    \operatorname{Grp}(\mathcal X_{/V})
    \xrightarrow{\sim}
    \operatorname*{Tot}_{[n]\in\Delta}
    \operatorname{Grp}(\mathcal X_{/U_n}),
\]
where $\operatorname{Grp}$ denotes internal group objects.  The equivalence is
compatible with the forgetful functors to the corresponding slice
$\infty$-categories and with arbitrary base change.
\end{lemma}

\begin{proof}
Effective descent in an $\infty$-topos gives
\[
    \mathcal X_{/V}
    \xrightarrow{\sim}
    \operatorname*{Tot}_{[n]\in\Delta}
    \mathcal X_{/U_n},
\]
with every transition functor given by pullback and therefore preserving
finite limits.

An internal group object is a model of the finite-product algebraic theory
of groups, equivalently a reduced Segal groupoid object.  Functor categories
commute with limits in the target, and the condition that a functor preserve
the finite products appearing in that theory is detected componentwise
because every transition functor in the displayed totalization preserves
finite limits.  Passing to the corresponding full subcategories gives
\[
    \operatorname{Grp}(\mathcal X_{/V})
    \simeq
    \operatorname*{Tot}_{[n]\in\Delta}
    \operatorname{Grp}(\mathcal X_{/U_n}).
\]
The underlying-object and base-change compatibilities are inherited from the
effective-descent equivalence of slices.
\end{proof}

\begin{theorem}[Descent formula for formal isotropy]
\label{thm:formal-isotropy-descent}
There are canonical equivalences of group objects
\[
    q^*\mathcal I_{Q/S}
    \simeq
    I^{\mathrm{raw}}_{\mathfrak G},
    \qquad
    q^*\mathcal K_{\mathfrak G}
    \simeq
    I^{\mathrm{for}}_{\mathfrak G}.
\]
The second equivalence transports the canonical \v{C}ech descent datum on
$q^*\mathcal K_{\mathfrak G}$ to
$I^{\mathrm{for}}_{\mathfrak G}$.  Under
Lemma~\ref{lem:effective-descent-group-objects},
$\mathcal K_{\mathfrak G}$ is, through a contractible space of compatible
choices, the unique group object over $Q$ realizing this descended formal
isotropy.  Thus formal isotropy descends canonically, rather than by an
additional descent datum, to the formal leaf space as the kernel of the
inertia morphism induced by
$r:Q\to Q_{X/S}$.
\end{theorem}

\begin{proof}
The first equivalence is
Proposition~\ref{prop:isotropy-inertia}.  For the second, pull the defining
finite limit of $\mathcal K_{\mathfrak G}$ back along $q$.  Since
$r q\simeq\epsilon_{X/S}$,
\[
\begin{aligned}
    q^*\mathcal K_{\mathfrak G}
    &\simeq
    q^*\mathcal I_{Q/S}
    \times_{q^*r^*\mathcal I_{Q^\infty/S}}
    X\\
    &\simeq
    I^{\mathrm{raw}}_{\mathfrak G}
    \times_{\epsilon_{X/S}^*\mathcal I_{Q^\infty/S}}
    X.
\end{aligned}
\]
The homomorphism in the second line is the pullback of the inertia morphism
induced by $r:Q\to Q^\infty$.  Under
\[
    \mathfrak G_\bullet\simeq\check C_\bullet(q),
    \qquad
    R_{X/S,\bullet}\simeq\check C_\bullet(\epsilon_{X/S}),
\]
it is exactly the homomorphism on raw isotropy induced by the canonical anchor
$\mathfrak G_\bullet\to R_{X/S,\bullet}$.

The identity-loop section of
$\epsilon_{X/S}^*\mathcal I_{Q^\infty/S}$ is the isotropy of the unit section
$X\to R_{X/S}$.  Taking its fibre therefore gives the finite-limit identity
\[
    I^{\mathrm{raw}}_{\mathfrak G}
    \times_{\epsilon_{X/S}^*\mathcal I_{Q^\infty/S}}
    X
    \simeq
    \mathfrak G_1\times_{R_{X/S}}X.
\]
Theorem~\ref{thm:isotropy-group} identifies the right-hand side with
$I^{\mathrm{for}}_{\mathfrak G}$ as a group object over $X$.

Since $q:X\twoheadrightarrow Q$ is an effective epimorphism,
Lemma~\ref{lem:effective-descent-group-objects} identifies group objects over
$Q$ with group objects over the complete \v{C}ech nerve
$\check C_\bullet(q)\simeq\mathfrak G_\bullet$ equipped with their coherent
descent data.  The pullback $q^*\mathcal K_{\mathfrak G}$ carries the
tautological descent datum of the group object
$\mathcal K_{\mathfrak G}\to Q$.  Transporting that datum through the
canonical equivalence
\[
    q^*\mathcal K_{\mathfrak G}
    \xrightarrow{\sim}
    I^{\mathrm{for}}_{\mathfrak G}
\]
constructs coherent descent data on formal isotropy.  Because the descent
functor is an equivalence of $\infty$-categories, the homotopy fibre over this
specified descent object is contractible.  Hence
$\mathcal K_{\mathfrak G}$ is the contractibly unique descended group object
realizing that formal-isotropy descent datum.
\end{proof}

\subsection{The nonlinear formal-isotropy commutator}

\begin{definition}[Nonlinear formal-isotropy commutator]
\label{def:isotropy-commutator}
Define
\[
    \operatorname{Comm}_{\mathfrak G}:
    I^{\mathrm{for}}_{\mathfrak G}
    \times_X
    I^{\mathrm{for}}_{\mathfrak G}
    \longrightarrow
    I^{\mathrm{for}}_{\mathfrak G}
\]
internally by
\[
    \operatorname{Comm}_{\mathfrak G}(a,b)
    :=
    aba^{-1}b^{-1}.
\]
\end{definition}

\begin{proposition}[Normalization and functoriality of the commutator]
\label{prop:isotropy-commutator-properties}
The commutator is canonically trivial when either input is the identity.
A morphism of formal Lie groupoids carries formal-isotropy commutators to
formal-isotropy commutators.
\end{proposition}

\begin{proof}
The normalization identities are group identities in
$I^{\mathrm{for}}_{\mathfrak G}$.  Functoriality follows because a groupoid
functor preserves composition, unit, inversion, and the canonical anchors,
hence preserves the anchor kernels.
\end{proof}

\begin{remark}[Not yet a bracket]
\label{rem:commutator-not-bracket}
The map $\operatorname{Comm}_{\mathfrak G}$ is multiplicative and nonlinear.
No bilinearity, antisymmetry, Jacobi identity, or additive truncation is
asserted.  Later brackets and higher partition-Lie operations arise only
after differentiation and recognition.  In general characteristic the
partition-Lie operations retain more information than an ordinary binary
dg Lie bracket \cite{BrantnerMathew2025,BrantnerCamposNuiten2025}.
\end{remark}


\section{Arrow images, bitorsors, and coherent translation}
\label{sec:flg-bitorsor}

For a one-object group, multiplication identifies every arrow fibre with the
identity fibre.  For a general groupoid two different effective arrow images
must be distinguished: the image of the source-target map and the image of
the infinitesimal anchor.

\begin{theorem}[Kernel bitorsor theorem for fixed-object groupoid morphisms]
\label{thm:kernel-bitorsor}
Let
\[
    \psi_\bullet:
    \mathfrak G_\bullet
    \longrightarrow
    \mathfrak H_\bullet
\]
be a morphism in
$\operatorname{Gpd}_X(\mathcal X_{/S})$.  Write
$u_{\mathfrak H}:X\to\mathfrak H_1$ for the unit of the target groupoid and
define the arrow-kernel
\[
    K_\psi
    :=
    \mathfrak G_1
    \times_{\psi_1,\;\mathfrak H_1,\;u_{\mathfrak H}}
    X.
\]
Then $K_\psi\to X$ is canonically a group object.  Factor the degree-one map
by its effective image
\[
    \mathfrak G_1
    \xrightarrow{\pi_\psi}
    M_\psi
    \xrightarrow{j_\psi}
    \mathfrak H_1,
\]
with $\pi_\psi$ an effective epimorphism and $j_\psi$ a monomorphism.  The
kernel acts on $\mathfrak G_1$ by target and source composition, and the two
associated shear morphisms are canonical equivalences
\[
\begin{aligned}
    \beta_t:
    t^*K_\psi
    &\xrightarrow{\sim}
    \mathfrak G_1\times_{M_\psi}\mathfrak G_1,
    &
    (g,h)&\longmapsto(h\circ g,g),\\
    \beta_s:
    s^*K_\psi
    &\xrightarrow{\sim}
    \mathfrak G_1\times_{M_\psi}\mathfrak G_1,
    &
    (g,k)&\longmapsto(g\circ k,g).
\end{aligned}
\]
The inverse morphisms are represented internally by
\[
    (g_1,g_2)
    \longmapsto
    (g_2,g_1\circ g_2^{-1})
\]
and
\[
    (g_1,g_2)
    \longmapsto
    (g_2,g_2^{-1}\circ g_1),
\]
where the required morphisms into the homotopy pullbacks defining
$t^*K_\psi$ and $s^*K_\psi$ are supplied canonically by the specified
homotopy in
$\mathfrak G_1\times_{\mathfrak H_1}\mathfrak G_1$.
\end{theorem}

\begin{proof}
Because $\psi_\bullet$ is a fixed-object groupoid functor, it induces a
homomorphism of raw isotropy group objects
\[
    I^{\mathrm{raw}}_{\mathfrak G}
    \longrightarrow
    I^{\mathrm{raw}}_{\mathfrak H}.
\]
Finite-limit pasting identifies the arrow-kernel with the group-object fibre
over the identity section:
\[
    K_\psi
    \simeq
    I^{\mathrm{raw}}_{\mathfrak G}
    \times_{I^{\mathrm{raw}}_{\mathfrak H}}
    X.
\]
Limits of internal group objects are created in the underlying slice, so this
pullback gives the canonical group-object structure on $K_\psi\to X$.

Since $j_\psi$ is a monomorphism, finite-limit pasting gives a canonical
equivalence
\[
    P_\psi
    :=
    \mathfrak G_1\times_{M_\psi}\mathfrak G_1
    \xrightarrow{\sim}
    \mathfrak G_1\times_{\mathfrak H_1}\mathfrak G_1.
\]
Let
\[
    p_1,p_2:P_\psi\rightrightarrows\mathfrak G_1
\]
be the projections.  Through the right-hand homotopy-pullback presentation,
$P_\psi$ carries its universal specified homotopy
\[
    \gamma:
    \psi_1p_1
    \simeq
    \psi_1p_2.
\]
Applying source and target to $\gamma$, and using that $\psi_\bullet$ is the
identity in degree zero, supplies the specified endpoint identifications
needed to form the internal difference morphisms
\[
    d_t:=p_1\circ p_2^{-1},
    \qquad
    d_s:=p_2^{-1}\circ p_1.
\]
Here and below the displayed composition notation denotes the morphisms
obtained from the projections, inversion, the composable-arrow pullback, and
groupoid composition; no strict equality of generalized elements is being
used.

Functoriality of composition and inversion under $\psi_\bullet$ carries the
homotopy $\gamma$ to canonical homotopies
\[
\begin{aligned}
    \psi_1d_t
    &\simeq
    \psi_1p_1\circ(\psi_1p_2)^{-1}
    \simeq
    u_{\mathfrak H}\,t p_2,\\
    \psi_1d_s
    &\simeq
    (\psi_1p_2)^{-1}\circ\psi_1p_1
    \simeq
    u_{\mathfrak H}\,s p_2.
\end{aligned}
\]
These specified homotopies are exactly the data required by the two homotopy
pullbacks.  Hence the universal properties of those pullbacks give canonical
morphisms
\[
    \delta_t:
    P_\psi\longrightarrow t^*K_\psi,
    \qquad
    \delta_s:
    P_\psi\longrightarrow s^*K_\psi
\]
whose first components are $p_2$ and whose kernel components are $d_t$ and
$d_s$, respectively.

Conversely, composition with a kernel arrow has unchanged image under
$\psi_1$.  More precisely, the defining homotopy
$\psi_1(h)\simeq u_{\mathfrak H}$ for a target-kernel arrow gives
\[
    \psi_1(h\circ g)
    \simeq
    \psi_1(h)\circ\psi_1(g)
    \simeq
    \psi_1(g),
\]
and the analogous source-kernel calculation gives
$\psi_1(g\circ k)\simeq\psi_1(g)$.  Thus the two action maps, together with
the unchanged arrow $g$, define the shear morphisms
$\beta_t$ and $\beta_s$ into $P_\psi$.

The simplicial groupoid unit, inverse, and associativity coherences give
\[
    (p_1\circ p_2^{-1})\circ p_2\simeq p_1,
    \qquad
    p_2\circ(p_2^{-1}\circ p_1)\simeq p_1,
\]
in the two relevant composable-arrow diagrams, and similarly
\[
    (h\circ g)\circ g^{-1}\simeq h,
    \qquad
    g^{-1}\circ(g\circ k)\simeq k.
\]
Tracing the specified pullback homotopies shows that these are identities in
the corresponding homotopy pullbacks, not merely identities after forgetting
to $\mathfrak G_1$.  Therefore
\[
    \beta_t\delta_t\simeq\operatorname{id},
    \qquad
    \delta_t\beta_t\simeq\operatorname{id},
    \qquad
    \beta_s\delta_s\simeq\operatorname{id},
    \qquad
    \delta_s\beta_s\simeq\operatorname{id}.
\]
Hence both shear morphisms are equivalences with the stated inverses.  All
higher compatibility is inherited functorially from the groupoid coherences
and the universal properties of the finite homotopy pullbacks.
\end{proof}

\subsection{Raw source-target bitorsors}

Factor the source-target map as
\[
    \mathfrak G_1
    \xrightarrow{\pi_{\mathrm{raw}}}
    M^{\mathrm{raw}}_{\mathfrak G}
    \xrightarrow{j_{\mathrm{raw}}}
    X\times_SX,
\]
where $\pi_{\mathrm{raw}}$ is an effective epimorphism and
$j_{\mathrm{raw}}$ is a monomorphism.

\begin{proposition}[Raw isotropy actions]
\label{prop:isotropy-actions}
Groupoid composition induces commuting left and right actions
\[
    a_L:
    t^*I^{\mathrm{raw}}_{\mathfrak G}
    \longrightarrow
    \mathfrak G_1,
    \qquad
    a_L(g,h)=h\circ g,
\]
and
\[
    a_R:
    s^*I^{\mathrm{raw}}_{\mathfrak G}
    \longrightarrow
    \mathfrak G_1,
    \qquad
    a_R(g,k)=g\circ k.
\]
\end{proposition}

\begin{proof}
The action maps are restrictions of groupoid composition.  Associativity gives
the action laws and
\[
    (h\circ g)\circ k
    \simeq
    h\circ(g\circ k).
\]
All higher compatibility is inherited from the simplicial associativity
coherence.
\end{proof}

\begin{theorem}[Raw Isotropy Bitorsor Theorem]
\label{thm:raw-bitorsor}
There are canonical equivalences
\[
\begin{aligned}
    t^*I^{\mathrm{raw}}_{\mathfrak G}
    &\xrightarrow{\sim}
    \mathfrak G_1
    \times_{M^{\mathrm{raw}}_{\mathfrak G}}
    \mathfrak G_1,
    &
    (g,h)&\longmapsto(h\circ g,g),\\
    s^*I^{\mathrm{raw}}_{\mathfrak G}
    &\xrightarrow{\sim}
    \mathfrak G_1
    \times_{M^{\mathrm{raw}}_{\mathfrak G}}
    \mathfrak G_1,
    &
    (g,k)&\longmapsto(g\circ k,g).
\end{aligned}
\]
Their inverses are
\[
    (g_1,g_2)
    \longmapsto
    (g_2,g_1\circ g_2^{-1})
\]
and
\[
    (g_1,g_2)
    \longmapsto
    (g_2,g_2^{-1}\circ g_1),
\]
respectively.
\end{theorem}

\begin{proof}
Apply Theorem~\ref{thm:kernel-bitorsor} to the canonical fixed-object vertex
morphism
\[
    \psi_\bullet:
    \mathfrak G_\bullet
    \longrightarrow
    \operatorname{Pair}(X/S)_\bullet.
\]
Its degree-one component is
\[
    \psi_1=(s,t):
    \mathfrak G_1\longrightarrow X\times_SX.
\]
The arrow-kernel of $\psi_1$ over the diagonal unit is exactly
$I^{\mathrm{raw}}_{\mathfrak G}$, and the effective image
$M_\psi$ is $M^{\mathrm{raw}}_{\mathfrak G}$.  The two shear equivalences and
their internally constructed inverse morphisms are therefore precisely the
displayed maps.
\end{proof}

\begin{remark}
Fibrewise, this says that every \emph{inhabited} source-target arrow fibre is
a left torsor for target raw isotropy and a right torsor for source raw
isotropy.  Passing through
$M^{\mathrm{raw}}_{\mathfrak G}$ rather than all of $X\times_SX$ records the
possible absence of arrows between two objects.
\end{remark}

\subsection{Formal anchor bitorsors}

Factor the canonical degree-one anchor as
\[
    \mathfrak G_1
    \xrightarrow{\pi_{\mathrm{for}}}
    M^{\mathrm{for}}_{\mathfrak G}
    \xrightarrow{j_{\mathrm{for}}}
    R_{X/S},
\]
with $\pi_{\mathrm{for}}$ effective and $j_{\mathrm{for}}$ monic.

\begin{proposition}[Formal isotropy actions preserve anchor fibres]
\label{prop:formal-isotropy-actions}
The left and right raw isotropy actions, precomposed with the completion
homomorphism
\[
    \kappa_I:
    I^{\mathrm{for}}_{\mathfrak G}
    \longrightarrow
    I^{\mathrm{raw}}_{\mathfrak G},
\]
induce canonically commuting actions
\[
    t^*I^{\mathrm{for}}_{\mathfrak G}
    \longrightarrow
    \mathfrak G_1,
    \qquad
    s^*I^{\mathrm{for}}_{\mathfrak G}
    \longrightarrow
    \mathfrak G_1
\]
which preserve the fibres of
$\pi_{\mathrm{for}}:\mathfrak G_1\to
M^{\mathrm{for}}_{\mathfrak G}$.
\end{proposition}

\begin{proof}
By Theorem~\ref{thm:isotropy-group}, formal isotropy is canonically the
homotopy fibre
\[
    I^{\mathrm{for}}_{\mathfrak G}
    \simeq
    \mathfrak G_1\times_{R_{X/S}}X
\]
of the canonical anchor over the unit section.  Under this identification,
the composite with
$\kappa_I:I^{\mathrm{for}}_{\mathfrak G}\to
I^{\mathrm{raw}}_{\mathfrak G}$
is the canonical map from this homotopy fibre to raw isotropy.  Since
$\phi_{\mathfrak G}$ is a groupoid morphism, composing an arrow $g$ with a
formal-isotropy arrow $h$ gives
\[
    \phi(h\circ g)
    \simeq
    \phi(h)\circ\phi(g)
    \simeq
    1\circ\phi(g)
    \simeq
    \phi(g),
\]
and similarly on the right.  Thus the actions obtained by precomposition with
$\kappa_I$ preserve anchor fibres.  The action coherences are inherited from
the raw actions.
\end{proof}

\begin{theorem}[Formal Bitorsor Theorem]
\label{thm:formal-bitorsor}
There are canonical equivalences
\[
\begin{aligned}
    t^*I^{\mathrm{for}}_{\mathfrak G}
    &\xrightarrow{\sim}
    \mathfrak G_1
    \times_{M^{\mathrm{for}}_{\mathfrak G}}
    \mathfrak G_1,
    &
    (g,h)&\longmapsto(h\circ g,g),\\
    s^*I^{\mathrm{for}}_{\mathfrak G}
    &\xrightarrow{\sim}
    \mathfrak G_1
    \times_{M^{\mathrm{for}}_{\mathfrak G}}
    \mathfrak G_1,
    &
    (g,k)&\longmapsto(g\circ k,g).
\end{aligned}
\]
Equivalently, because $j_{\mathrm{for}}$ is monic,
\[
    t^*I^{\mathrm{for}}_{\mathfrak G}
    \simeq
    \mathfrak G_1\times_{R_{X/S}}\mathfrak G_1
    \simeq
    s^*I^{\mathrm{for}}_{\mathfrak G}.
\]
Thus every inhabited fibre of the infinitesimal anchor is canonically a
left torsor for target formal isotropy and a right torsor for source formal
isotropy, with the two torsor structures commuting coherently.
\end{theorem}

\begin{proof}
Apply Theorem~\ref{thm:kernel-bitorsor} to the canonical infinitesimal anchor
\[
    \phi_{\mathfrak G}:
    \mathfrak G_\bullet
    \longrightarrow
    R_{X/S,\bullet}.
\]
Its arrow-kernel is
\[
    K_{\phi_{\mathfrak G}}
    =
    \mathfrak G_1\times_{R_{X/S}}X
    \simeq
    I^{\mathrm{for}}_{\mathfrak G}
\]
by Theorem~\ref{thm:isotropy-group}, and its effective arrow image is
$M^{\mathrm{for}}_{\mathfrak G}$.  The general kernel-bitorsor theorem
therefore gives both displayed equivalences, including the specified
homotopy-pullback witnesses in the inverse maps.
\end{proof}

\begin{remark}[Why the one-object case is exceptional]
\label{rem:group-case-exceptional}
For a one-object groupoid all arrows have the same source and target and the
formal isotropy group is the arrow group itself.  The bitorsor therefore
admits a global translation trivialization by the single identity fibre.
For a general groupoid formal isotropy varies over $X$, and the intrinsic
replacement for a global translation coordinate is the formal anchor
bitorsor above.
\end{remark}


\section{Conjugation and the formal adjoint local system}
\label{sec:flg-adjoint}

\begin{definition}[Formal adjoint transport]
\label{def:formal-adjoint-transport}
Under the anchor-kernel equivalence
\[
    I^{\mathrm{for}}_{\mathfrak G}
    \simeq
    \mathfrak G_1\times_{R_{X/S}}X,
\]
groupoid conjugation preserves the homotopy fibre of the canonical anchor over
the unit section and therefore induces a morphism of group objects over
$\mathfrak G_1$
\[
    \operatorname{Ad}_{\mathfrak G}:
    s^*I^{\mathrm{for}}_{\mathfrak G}
    \longrightarrow
    t^*I^{\mathrm{for}}_{\mathfrak G}.
\]
In generalized-element notation, for an arrow $g:x\to y$ and
$h\in I^{\mathrm{for}}_{\mathfrak G,x}$ this induced morphism is
\[
    \operatorname{Ad}_g(h)
    :=
    g h g^{-1}.
\]
\end{definition}

\begin{theorem}[Coherent formal adjoint transport]
\label{thm:coherent-adjoint-transport}
The morphism
$\operatorname{Ad}_{\mathfrak G}$ is an equivalence and is the degree-one
transport of the canonical effective descent datum on
$I^{\mathrm{for}}_{\mathfrak G}$ constructed in
Theorem~\ref{thm:formal-isotropy-descent}.  In degree zero it is the identity,
and for composable arrows
\[
    x\xrightarrow{g}y\xrightarrow{h}z
\]
there is a canonical coherent equivalence
\[
    \operatorname{Ad}_{h\circ g}
    \simeq
    \operatorname{Ad}_h\operatorname{Ad}_g.
\]
All higher cocycle compatibilities are the canonical \v{C}ech descent
coherences transported from
$\mathcal K_{\mathfrak G}\to Q_{\mathfrak G}$.
\end{theorem}

\begin{proof}
Put
\[
    q:=q_{\mathfrak G}:X\twoheadrightarrow Q_{\mathfrak G}.
\]
By Theorem~\ref{thm:formal-isotropy-descent}, the canonical equivalence
\[
    q^*\mathcal K_{\mathfrak G}
    \xrightarrow{\sim}
    I^{\mathrm{for}}_{\mathfrak G}
\]
transports the tautological \v{C}ech descent datum of the group object
$\mathcal K_{\mathfrak G}\to Q_{\mathfrak G}$ to a coherent descent datum on
formal isotropy along
\[
    \mathfrak G_\bullet
    \simeq
    \check C_\bullet(q).
\]
Its degree-one term is an equivalence
\[
    \theta:
    s^*I^{\mathrm{for}}_{\mathfrak G}
    \xrightarrow{\sim}
    t^*I^{\mathrm{for}}_{\mathfrak G},
\]
and its degree-two and higher terms already supply the complete cocycle
coherences.  It remains only to identify $\theta$ with the conjugation map of
Definition~\ref{def:formal-adjoint-transport}.

Under
\[
    q^*\mathcal I_{Q_{\mathfrak G}/S}
    \simeq
    I^{\mathrm{raw}}_{\mathfrak G},
\]
the tautological \v{C}ech descent transport of the pulled-back inertia object
is raw groupoid conjugation.  Indeed, over
\[
    \mathfrak G_1
    \simeq
    X\times_{Q_{\mathfrak G}}X
\]
the defining homotopy of the fibre product identifies the two pullbacks of
the same loop over $Q_{\mathfrak G}$, and in the effective groupoid this
identification is represented internally by
\[
    (g,h)\longmapsto ghg^{-1}.
\]
The inertia homomorphism
\[
    \mathcal I_{Q_{\mathfrak G}/S}
    \longrightarrow
    r_{\mathfrak G}^*\mathcal I_{Q_{X/S}/S}
\]
and the identity-loop section are morphisms of objects over
$Q_{\mathfrak G}$, so their pullbacks preserve this canonical descent
transport.  Taking the fibre over the identity-loop section therefore
induces on the homotopy fibre $q^*\mathcal K_{\mathfrak G}$ the corresponding
conjugation transport.  Under the anchor-kernel identification
\[
    I^{\mathrm{for}}_{\mathfrak G}
    \simeq
    \mathfrak G_1\times_{R_{X/S}}X,
\]
that induced map is precisely
\[
    h\longmapsto ghg^{-1}
    =
    \operatorname{Ad}_g(h).
\]
Hence $\theta\simeq\operatorname{Ad}_{\mathfrak G}$.

Since $\theta$ is a descent equivalence, so is
$\operatorname{Ad}_{\mathfrak G}$.  The identity, composition, and all higher
coherences are now the transported \v{C}ech descent coherences; in particular,
on composable arrows they give
\[
    \operatorname{Ad}_{h\circ g}
    \simeq
    \operatorname{Ad}_h\operatorname{Ad}_g.
\]
No additional adjoint local-system datum is imposed.
\end{proof}

\begin{proposition}[Inertia descent is conjugation]
\label{prop:inertia-descent-conjugation}
Put
\[
    q:=q_{\mathfrak G}:X\twoheadrightarrow Q_{\mathfrak G}.
\]
Under the canonical equivalence
\[
    q^*\mathcal I_{Q_{\mathfrak G}/S}
    \simeq
    I^{\mathrm{raw}}_{\mathfrak G}
\]
of Proposition~\ref{prop:isotropy-inertia}, the canonical \v{C}ech descent
transport of the pulled-back inertia group along
\[
    \mathfrak G_\bullet
    \simeq
    \check C_\bullet(q)
\]
is the raw conjugation transport
\[
    h\longmapsto ghg^{-1}.
\]
Under the canonical equivalence
\[
    q^*\mathcal K_{\mathfrak G}
    \simeq
    I^{\mathrm{for}}_{\mathfrak G}
\]
of Theorem~\ref{thm:formal-isotropy-descent}, the induced descent transport on
the formal inertia kernel is the formal adjoint transport of
Theorem~\ref{thm:coherent-adjoint-transport}.
\end{proposition}

\begin{proof}
Over
\[
    \mathfrak G_1
    \simeq
    X\times_{Q_{\mathfrak G}}X
\]
the two composites to $Q_{\mathfrak G}$ are canonically identified by the
defining pullback.  The canonical descent transport on
$q^*\mathcal I_{Q_{\mathfrak G}/S}$ carries a loop based at the first lift
across this identification to the corresponding loop based at the second
lift.  In the effective groupoid
$\check C_\bullet(q)$ this transport is represented internally by
\[
    (g,h)\longmapsto ghg^{-1}.
\]
Under the equivalence
$q^*\mathcal I_{Q_{\mathfrak G}/S}\simeq
I^{\mathrm{raw}}_{\mathfrak G}$ this is exactly groupoid conjugation.
On
\[
    \mathfrak G_2
    \simeq
    X\times_{Q_{\mathfrak G}}X\times_{Q_{\mathfrak G}}X
\]
the \v{C}ech cocycle identity is the canonical equivalence
\[
    (h\circ g)k(h\circ g)^{-1}
    \simeq
    h(gkg^{-1})h^{-1},
\]
and the higher descent coherences are the corresponding higher simplicial
coherences.

The inertia morphism
\[
    \mathcal I_{Q_{\mathfrak G}/S}
    \longrightarrow
    r_{\mathfrak G}^*
    \mathcal I_{Q_{X/S}/S}
\]
is functorial in the quotient map and therefore intertwines these conjugation
transports.  Its identity-loop section is also preserved.  Taking the fibre
over that section consequently induces the raw conjugation transport on the
homotopy fibre $q^*\mathcal K_{\mathfrak G}$.  Under
Theorem~\ref{thm:formal-isotropy-descent}, that fibre is
$I^{\mathrm{for}}_{\mathfrak G}$, and the induced conjugation is precisely
the formal adjoint transport.
\end{proof}

\begin{corollary}[Formal adjoint group and conjugation descent]
\label{cor:adjoint-inertia-object}
For the constructed formal adjoint group
\[
    \operatorname{Ad}^{\mathrm{for}}(\mathfrak G)
    =
    \mathcal K_{\mathfrak G}
    \longrightarrow
    Q_{\mathfrak G},
\]
there is a canonical equivalence
\[
    q^*\operatorname{Ad}^{\mathrm{for}}(\mathfrak G)
    \simeq
    I^{\mathrm{for}}_{\mathfrak G}.
\]
Under this equivalence, the canonical \v{C}ech transport on the pullback of
$\operatorname{Ad}^{\mathrm{for}}(\mathfrak G)$ is the formal adjoint
transport of Theorem~\ref{thm:coherent-adjoint-transport}.  The raw isotropy
object separately is
\[
    I^{\mathrm{raw}}_{\mathfrak G}
    \simeq
    q^*\mathcal I_{Q_{\mathfrak G}/S}.
\]
\end{corollary}

\begin{proof}
The first equivalence is
Theorem~\ref{thm:formal-isotropy-descent} together with
Definition~\ref{def:formal-inertia-kernel}.  Proposition~\ref{prop:inertia-descent-conjugation}
identifies the canonical \v{C}ech transport on this pullback with groupoid
conjugation, which is exactly the formal adjoint transport.  The raw statement
is Proposition~\ref{prop:isotropy-inertia}.
\end{proof}

\begin{proposition}[Adjoint covariance of the nonlinear commutator]
\label{prop:adjoint-commutator}
Over $\mathfrak G_1$ there is a canonical equivalence
\[
    \operatorname{Ad}_{\mathfrak G}
    \left(
        \operatorname{Comm}_{\mathfrak G}(a,b)
    \right)
    \simeq
    \operatorname{Comm}_{\mathfrak G}
    \left(
        \operatorname{Ad}_{\mathfrak G}(a),
        \operatorname{Ad}_{\mathfrak G}(b)
    \right).
\]
\end{proposition}

\begin{proof}
Conjugation by an invertible arrow is a homomorphism between the corresponding
formal isotropy group objects and therefore preserves multiplication and
inversion.  Apply it to $aba^{-1}b^{-1}$.
\end{proof}

\section{The formal-groupoid disk monad and jet comonad}
\label{sec:flg-jet}

\subsection{Effective-image relation calculus for arbitrary morphisms}
\label{sec:flg-effective-relation-calculus}

The disk--jet mechanism used below is not an additional de Rham axiom and is
not restricted to effective quotient maps.  Every morphism in the ambient
$\infty$-topos has a canonical effective image, and the complete relation
calculus is forced by that image.

\begin{theorem}[Effective-Image Relation Disk--Jet Calculus]
\label{thm:effective-relation-calculus}
Let
\[
    q:X\longrightarrow B
\]
be an arbitrary morphism in $\mathcal E$, and factor it canonically as
\[
    X
    \xrightarrow{e_q}
    I_q
    \xrightarrow{m_q}
    B
\]
with $e_q$ an effective epimorphism and $m_q$ a monomorphism.  Put
\[
    R^q_\bullet:=\check C_\bullet(q),
    \qquad
    R^q:=R^q_1=X\times_BX.
\]
Then there are canonical equivalences
\[
    R^q_\bullet
    \simeq
    \check C_\bullet(e_q),
    \qquad
    |R^q_\bullet|\simeq I_q.
\]
Write
\[
    s,t:R^q\rightrightarrows X
\]
for the two projections, $u:X\to R^q$ for the diagonal, and
$\iota:R^q\xrightarrow{\sim}R^q$ for inversion.  Define endofunctors of
$\mathcal X_{/X}$ by
\[
    T_q:=q^*\Sigma_q,
    \qquad
    J_q:=q^*\Pi_q.
\]
Then the following structures and equivalences are canonical.
\begin{enumerate}
    \item Effective-image invariance gives equivalences of monads and comonads
    \[
        T_q\simeq e_q^*\Sigma_{e_q},
        \qquad
        J_q\simeq e_q^*\Pi_{e_q}.
    \]
    \item $T_q$ is the monad induced by $\Sigma_q\dashv q^*$ and $J_q$ is
    the comonad induced by $q^*\dashv\Pi_q$.  There are relation formulas
    \[
        T_q\simeq\Sigma_s t^*
        \simeq\Sigma_t s^*,
        \qquad
        J_q\simeq\Pi_s t^*.
    \]
    Under these formulas, the monad unit and comonad counit are induced by
    $u$, while the monad multiplication and comonad comultiplication are
    induced by composition
    \[
        R^q_2
        =
        X\times_{I_q}X\times_{I_q}X
        \longrightarrow
        X\times_{I_q}X,
        \qquad
        (x_0,x_1,x_2)\longmapsto(x_0,x_2).
    \]
    \item There is a canonical adjunction
    \[
        T_q\dashv J_q.
    \]
    More precisely,
    \[
    \begin{aligned}
        \operatorname{Map}_X(T_qA,E)
        &\simeq
        \operatorname{Map}_{R^q}(t^*A,s^*E)\\
        &\simeq
        \operatorname{Map}_{R^q}(s^*A,t^*E)\\
        &\simeq
        \operatorname{Map}_X(A,J_qE),
    \end{aligned}
    \]
    naturally in $A,E\in\mathcal X_{/X}$.
    \item For $E\to X$, a $T_q$-algebra structure, a $J_q$-coalgebra
    structure, and coherent transport along $R^q_\bullet$ are canonically
    equivalent.  The degree-one transport is an equivalence
    \[
        \theta:s^*E\xrightarrow{\sim}t^*E
    \]
    satisfying
    \[
        u^*\theta\simeq\operatorname{id}_E
    \]
    and, on $R^q_2$,
    \[
        p_{12}^*\theta\circ p_{01}^*\theta
        \simeq
        p_{02}^*\theta,
    \]
    together with the higher coherences on the higher terms of the
    \v{C}ech nerve.  Under these equivalences the algebra action and
    coalgebra coaction are mates under $T_q\dashv J_q$.
    \item Pullback along the effective image presentation
    $e_q:X\twoheadrightarrow I_q$ induces canonical equivalences
    \[
        \mathcal X_{/I_q}
        \simeq
        \operatorname{Alg}_{T_q}(\mathcal X_{/X})
        \simeq
        \operatorname{Coalg}_{J_q}(\mathcal X_{/X}),
    \]
    compatible with the forgetful functors to $\mathcal X_{/X}$.
\end{enumerate}
More precisely, for every commutative square
\[
\begin{tikzcd}
X \arrow[r,"f"] \arrow[d,"q"'] &
X' \arrow[d,"q'"]\\
B \arrow[r,"g"'] &
B'
\end{tikzcd}
\]
there are canonical mate transformations
\[
    T_qf^*\longrightarrow f^*T_{q'},
    \qquad
    f^*J_{q'}\longrightarrow J_qf^*,
\]
compatible with the monad and comonad structures, identities, composition,
and the induced morphism of effective-image relation groupoids.  When the
square is Cartesian, these transformations are Beck--Chevalley
equivalences.  In particular, for an arbitrary base change, that is, for the
corresponding Cartesian pullback square, all of the displayed constructions
commute with base change through these canonical equivalences.
\end{theorem}

\begin{proof}
Because $m_q:I_q\to B$ is a monomorphism,
\[
    I_q\xrightarrow{\sim}I_q\times_BI_q.
\]
Finite-limit pasting therefore gives
\[
    X\times_B\cdots\times_BX
    \simeq
    X\times_{I_q}\cdots\times_{I_q}X
\]
in every degree, compatibly with deletion and repetition of factors.  Hence
\[
    R^q_\bullet\simeq\check C_\bullet(e_q),
\]
and effectivity of $e_q$ gives
\[
    |R^q_\bullet|\simeq I_q.
\]

We next prove the effective-image invariance of the endofunctors.  Since
$q=m_qe_q$,
\[
    \Sigma_q\simeq\Sigma_{m_q}\Sigma_{e_q},
    \qquad
    \Pi_q\simeq\Pi_{m_q}\Pi_{e_q},
    \qquad
    q^*=e_q^*m_q^*.
\]
For a monomorphism $m_q$ there is a canonical equivalence
\[
    m_q^*\Sigma_{m_q}\simeq\operatorname{id}_{\mathcal X_{/I_q}}.
\]
Indeed, for $A\to I_q$,
\[
    m_q^*\Sigma_{m_q}A
    =
    A\times_BI_q
    \simeq
    A\times_{I_q}(I_q\times_BI_q)
    \simeq A.
\]
There is also a canonical equivalence
\[
    m_q^*\Pi_{m_q}\simeq\operatorname{id}_{\mathcal X_{/I_q}}.
\]
For $A,E\in\mathcal X_{/I_q}$,
\[
\begin{aligned}
    \operatorname{Map}_{I_q}(A,m_q^*\Pi_{m_q}E)
    &\simeq
    \operatorname{Map}_{B}(\Sigma_{m_q}A,\Pi_{m_q}E)\\
    &\simeq
    \operatorname{Map}_{I_q}(m_q^*\Sigma_{m_q}A,E)\\
    &\simeq
    \operatorname{Map}_{I_q}(A,E).
\end{aligned}
\]
Yoneda gives the second equivalence.  Substitution yields
\[
    T_q=q^*\Sigma_q
    \simeq
    e_q^*\Sigma_{e_q},
    \qquad
    J_q=q^*\Pi_q
    \simeq
    e_q^*\Pi_{e_q}.
\]
These equivalences arise from equivalences of the composing adjunctions and
therefore identify the monad and comonad structures.

The square
\[
\begin{tikzcd}
R^q \arrow[r,"t"] \arrow[d,"s"'] & X \arrow[d,"q"]\\
X \arrow[r,"q"'] & B
\end{tikzcd}
\]
is Cartesian.  Left and right Beck--Chevalley give
\[
    q^*\Sigma_q\simeq\Sigma_s t^*,
    \qquad
    q^*\Pi_q\simeq\Pi_s t^*.
\]
Pullback along inversion exchanges $s$ and $t$, giving
\[
    \Sigma_s t^*\simeq\Sigma_t s^*.
\]
Under the effective-image identification
$R^q_\bullet\simeq\check C_\bullet(e_q)$, the monad and comonad structure
maps are the images of the diagonal and the degree-two \v{C}ech composition
map.  Their identities and higher coherences are therefore the simplicial
identities of $R^q_\bullet$.

For the adjunction, use the relation formulas and the two dependent
adjunctions:
\[
\begin{aligned}
    \operatorname{Map}_X(T_qA,E)
    &\simeq
    \operatorname{Map}_{R^q}(t^*A,s^*E)\\
    &\simeq
    \operatorname{Map}_{R^q}(s^*A,t^*E)\\
    &\simeq
    \operatorname{Map}_X(A,J_qE).
\end{aligned}
\]
The middle equivalence is pullback along $\iota$.

Now let $E\to X$.  Under $T_q\simeq\Sigma_s t^*$, an algebra action
\[
    a:T_qE\longrightarrow E
\]
is adjoint to a morphism
\[
    \bar a:t^*E\longrightarrow s^*E.
\]
Transport along inversion gives
\[
    \theta_a:=\iota^*\bar a:
    s^*E\longrightarrow t^*E.
\]
The monad unit law is exactly
\[
    u^*\theta_a\simeq\operatorname{id}_E,
\]
and the monad multiplication law is exactly
\[
    p_{12}^*\theta_a\circ p_{01}^*\theta_a
    \simeq
    p_{02}^*\theta_a
\]
on $R^q_2$.  Pulling this cocycle equivalence back along
\[
    (x_0,x_1)\longmapsto(x_0,x_1,x_0)
    \quad\text{and}\quad
    (x_0,x_1)\longmapsto(x_1,x_0,x_1)
\]
and using the unit identity gives
\[
    \iota^*\theta_a\circ\theta_a\simeq\operatorname{id},
    \qquad
    \theta_a\circ\iota^*\theta_a\simeq\operatorname{id}.
\]
Hence $\theta_a$ is automatically an equivalence.

Similarly, under $J_q\simeq\Pi_s t^*$ a coalgebra coaction
\[
    \rho:E\longrightarrow J_qE
\]
is adjoint directly to
\[
    \theta_\rho:s^*E\longrightarrow t^*E.
\]
Its counit and coassociativity laws give the same unit and cocycle identities,
and the mapping-space construction of $T_q\dashv J_q$ shows that $a$ and
$\rho$ are mates precisely when $\theta_a=\theta_\rho$.

It remains to identify the full higher coherence rather than only its
degree-one and degree-two shadows.  Since
$e_q:X\twoheadrightarrow I_q$ is an effective epimorphism, effective descent
in the ambient $\infty$-topos gives
\[
    \mathcal X_{/I_q}
    \simeq
    \operatorname*{Tot}_{[n]\in\Delta}
    \mathcal X_{/R^q_n}.
\]
The right-hand side is the $\infty$-category of objects over $X$ equipped with
coherent transport along the complete \v{C}ech nerve
$R^q_\bullet\simeq\check C_\bullet(e_q)$.

Under
\[
    T_q\simeq e_q^*\Sigma_{e_q},
    \qquad
    J_q\simeq e_q^*\Pi_{e_q},
\]
the bar construction of the monad and the cobar construction of the comonad
are obtained from the same iterated pullback diagrams
\[
    R^q_n
    =
    \underbrace{
    X\times_{I_q}\cdots\times_{I_q}X
    }_{n+1}.
\]
The unit and counit correspond to the diagonal degeneracies, while
multiplication and comultiplication correspond to the face maps induced by
composition of the effective relation.  The dependent-adjunction mate
equivalences identify the algebra and coalgebra structure maps with the same
simplicial transport maps in every degree.  Consequently the full
Eilenberg--Moore mapping spaces, not merely the underlying objects, identify
with the same totalization.  Thus
\[
    \mathcal X_{/I_q}
    \simeq
    \operatorname{Alg}_{T_q}(\mathcal X_{/X})
    \simeq
    \operatorname{Coalg}_{J_q}(\mathcal X_{/X}),
\]
and the algebra action and coalgebra coaction are mates under
$T_q\dashv J_q$.  This supplies all higher transport coherences.

Finally, let
\[
\begin{tikzcd}
X \arrow[r,"f"] \arrow[d,"q"'] &
X' \arrow[d,"q'"]\\
B \arrow[r,"g"'] &
B'
\end{tikzcd}
\]
be a commutative square.  The locally Cartesian closed structure supplies the
canonical left and right mate transformations
\[
    \Sigma_q f^*\longrightarrow g^*\Sigma_{q'},
    \qquad
    g^*\Pi_{q'}\longrightarrow\Pi_qf^*.
\]
Pulling the first along $q$ and the second along $q$, and using
$q^*g^*\simeq f^*(q')^*$, gives
\[
    T_qf^*\longrightarrow f^*T_{q'},
    \qquad
    f^*J_{q'}\longrightarrow J_qf^*.
\]
Because these are mates of the morphism of dependent adjunctions determined
by the square, they commute with the monad and comonad structure maps and
their higher coherences.  Functorial effective-image factorization supplies
the compatible map of effective images and hence of the relation groupoids.
For a Cartesian square the two mate transformations are the
Beck--Chevalley equivalences.  This proves the asserted functoriality and, in
particular, compatibility with arbitrary base change through the associated
Cartesian pullback squares.
\end{proof}




The quotient presentation
\[
    q:X\twoheadrightarrow Q_{\mathfrak G}
\]
has the dependent adjoint triple
\[
    \Sigma_q
    \dashv
    q^*
    \dashv
    \Pi_q.
\]
Apply Theorem~\ref{thm:effective-relation-calculus} to the effective quotient
$q:X\twoheadrightarrow Q_{\mathfrak G}$.  The following constructions are its
formal-groupoid specialization.

\begin{definition}[Formal-groupoid disk and jet endofunctors]
\label{def:formal-groupoid-disk-jet}
Define endofunctors of $\mathcal X_{/X}$ by
\[
    T_{\mathfrak G}
    :=
    q^*\Sigma_q,
    \qquad
    J_{\mathfrak G}
    :=
    q^*\Pi_q.
\]
The first carries its canonical monad structure from
$\Sigma_q\dashv q^*$, and the second carries its canonical comonad structure
from $q^*\dashv\Pi_q$.
\end{definition}

\subsection{Relation-correspondence formulas}

Write
\[
    s,t:\mathfrak G_1\rightrightarrows X.
\]

\begin{proposition}[Relation formulas]
\label{prop:formal-groupoid-relation-formulas}
There are canonical equivalences
\[
    T_{\mathfrak G}
    \simeq
    \Sigma_s t^*,
    \qquad
    J_{\mathfrak G}
    \simeq
    \Pi_s t^*.
\]
Transport along inversion gives a further equivalence
\[
    T_{\mathfrak G}
    \simeq
    \Sigma_t s^*.
\]
Under these equivalences, the monad unit and comonad counit are induced by
$u:X\to\mathfrak G_1$, while monad multiplication and comonad
comultiplication are induced by groupoid composition.
\end{proposition}

\begin{proof}
The square
\[
\begin{tikzcd}
\mathfrak G_1 \arrow[r,"t"] \arrow[d,"s"'] & X \arrow[d,"q"]\\
X \arrow[r,"q"'] & Q_{\mathfrak G}
\end{tikzcd}
\]
is Cartesian because
$\mathfrak G_1\simeq X\times_{Q_{\mathfrak G}}X$.  Left and right
Beck--Chevalley therefore give
\[
    q^*\Sigma_q\simeq\Sigma_s t^*,
    \qquad
    q^*\Pi_q\simeq\Pi_s t^*.
\]
Inversion exchanges $s$ and $t$, producing
$\Sigma_s t^*\simeq\Sigma_t s^*$.  The monad and comonad structure maps are
the standard ones attached to the two adjunctions.  Under the Cech
correspondence, their degree-zero and degree-two structure maps are precisely
the groupoid unit and composition maps.
\end{proof}

\begin{theorem}[Formal-groupoid disk--jet adjunction]
\label{thm:formal-groupoid-disk-jet-adjunction}
There is a canonical adjunction
\[
    T_{\mathfrak G}
    \dashv
    J_{\mathfrak G}.
\]
More precisely, for $A,E\in\mathcal X_{/X}$ there is a natural chain of
equivalences
\[
\begin{aligned}
    \operatorname{Map}_X(T_{\mathfrak G}A,E)
    &\simeq
    \operatorname{Map}_{\mathfrak G_1}(t^*A,s^*E)
\\
    &\simeq
    \operatorname{Map}_{\mathfrak G_1}(s^*A,t^*E)
\\
    &\simeq
    \operatorname{Map}_X(A,J_{\mathfrak G}E),
\end{aligned}
\]
where the middle equivalence is transport along inversion.
\end{theorem}

\begin{proof}
Use
$T_{\mathfrak G}\simeq\Sigma_s t^*$ and
$\Sigma_s\dashv s^*$ for the first equivalence.  Pullback along
$\iota:\mathfrak G_1\simeq\mathfrak G_1$ exchanges $s$ and $t$, giving the
middle equivalence.  Finally use $s^*\dashv\Pi_s$ and
$J_{\mathfrak G}\simeq\Pi_s t^*$.
\end{proof}

\begin{proposition}[Fibrewise formal-disk formula]
\label{prop:formal-groupoid-fibrewise-jets}
Let $a:T\to X$.  For every $E\to X$ there is a canonical equivalence
\[
    a^*J_{\mathfrak G}E
    \simeq
    \Pi_{d^{\mathfrak G}_a}
    \left(
        \operatorname{ev}^{\mathfrak G}_a
    \right)^*E.
\]
Equivalently, the generalized fibre of $J_{\mathfrak G}E$ at $a$ is the object
of sections of $E$ over the $\mathfrak G$-formal orbit disk centered at $a$.
\end{proposition}

\begin{proof}
Pull back
$J_{\mathfrak G}\simeq\Pi_s t^*$ along $a:T\to X$ and apply right
Beck--Chevalley to the Cartesian square of
Proposition~\ref{prop:formal-disk-arrow-fibre}.
\end{proof}

\subsection{Comparison with the maximal de Rham disk and jet calculus}

Write
\[
    T^\infty_{X/S}
    :=
    \epsilon_{X/S}^*\Sigma_{\epsilon_{X/S}},
    \qquad
    J^\infty_{X/S}
    :=
    \epsilon_{X/S}^*\Pi_{\epsilon_{X/S}}
\]
for the already-constructed maximal formal-disk monad and jet comonad.  Since
$\epsilon_{X/S}=rq$, dependent adjunction supplies canonical comparison maps.

\begin{proposition}[Ambient disk and jet comparisons]
\label{prop:formal-groupoid-ambient-monad-comonad}
There are canonical natural transformations
\[
    T_{\mathfrak G}
    \longrightarrow
    T^\infty_{X/S},
    \qquad
    J^\infty_{X/S}
    \longrightarrow
    J_{\mathfrak G}.
\]
The first is induced by the unit
$\operatorname{id}\to r^*\Sigma_r$ of
$\Sigma_r\dashv r^*$, and the second by the counit
$r^*\Pi_r\to\operatorname{id}$ of $r^*\dashv\Pi_r$.
Under the fibrewise formulas they are respectively extension along
\[
    \mathbb D^{\mathfrak G}_a
    \longrightarrow
    \mathbb D^\infty_{X/S,a}
\]
and restriction of sections along the same map.
\end{proposition}

\begin{proof}
Using $\epsilon_{X/S}=rq$,
\[
\begin{aligned}
    T^\infty_{X/S}
    &=
    q^*r^*\Sigma_r\Sigma_q,\\
    J^\infty_{X/S}
    &=
    q^*r^*\Pi_r\Pi_q.
\end{aligned}
\]
Apply $q^*(-)\Sigma_q$ to the unit
$\operatorname{id}\to r^*\Sigma_r$ for the first map, and apply
$q^*(-)\Pi_q$ to the counit $r^*\Pi_r\to\operatorname{id}$ for the second.
The fibrewise interpretation follows from the universal properties of
dependent sum and dependent product.
\end{proof}

\section{Full infinitesimal targets and effective groupoids}
\label{sec:flg-full-effective}

The effective formal leaf space is the minimal quotient carrying the
kernel-pair geometry of a formal Lie groupoid.  It is sometimes useful to
retain a separate auxiliary target mapping to the full modal object, just as
the preceding de Rham theory retains both $Q_{X/S}$ and
$X_{\mathrm{dR}/S}$.

\begin{definition}[Auxiliary full infinitesimal comparison diagram]
\label{def:full-infinitesimal-target}
When an external construction supplies an object $\mathfrak X$ and morphisms
\[
    X
    \xrightarrow{a}
    \mathfrak X
    \xrightarrow{b}
    X_{\mathrm{dR}/S}
\]
with a specified identification
\[
    ba\simeq\eta_{X/S},
\]
we call this an \emph{auxiliary full infinitesimal comparison diagram}.
It is not part of the definition or construction of a formal Lie groupoid.
No effectivity of $a$, representability of $\mathfrak X$, or de Rham-locality
of $\mathfrak X$ is inferred from the diagram.
\end{definition}

\begin{remark}
The adjective ``auxiliary'' is deliberate.  The intrinsic theory neither
chooses nor requires $\mathfrak X$.  When such a comparison object has been
constructed for another purpose, it may contain excess information not seen
by the kernel pair of $X\to\mathfrak X$.  The intrinsic construction below
canonically discards that excess by passing to the effective image.
\end{remark}

Factor $a$ as
\[
    X
    \xrightarrow{\bar a}
    Q_a
    \xrightarrow{i_a}
    \mathfrak X,
\]
with $\bar a$ an effective epimorphism and $i_a$ a monomorphism.

\begin{proposition}[Effective-image kernel-pair invariance]
\label{prop:kernel-pair-image-invariance}
For every auxiliary comparison diagram there are canonical equivalences
\[
    \underbrace{
    X\times_{\mathfrak X}\cdots\times_{\mathfrak X}X
    }_{n+1}
    \simeq
    \underbrace{
    X\times_{Q_a}\cdots\times_{Q_a}X
    }_{n+1}
\]
for every $n\ge0$, compatibly with all simplicial operators.  Hence
\[
    \check C_\bullet(X\to\mathfrak X)
    \simeq
    \check C_\bullet(X\twoheadrightarrow Q_a).
\]
Moreover the composite to $X_{\mathrm{dR}/S}$ induces a canonical factorization
\[
    X\twoheadrightarrow Q_a\longrightarrow Q_{X/S}.
\]
Thus the kernel-pair groupoid of an auxiliary target is canonically anchored,
its unitwise completion is a brave-new formal Lie groupoid, and the
effective-image relation disk--jet calculus attached to
$a:X\to\mathfrak X$ is canonically the calculus attached to
$\bar a:X\twoheadrightarrow Q_a$.
\end{proposition}

\begin{proof}
Since $i_a$ is a monomorphism,
\[
    Q_a\simeq Q_a\times_{\mathfrak X}Q_a.
\]
Finite-limit pasting gives
\[
    X\times_{\mathfrak X}X
    \simeq
    X\times_{Q_a}
    \left(
        Q_a\times_{\mathfrak X}Q_a
    \right)
    \times_{Q_a}X
    \simeq
    X\times_{Q_a}X,
\]
and the same calculation in every degree gives the simplicial equivalence.

The composite
\[
    X\to Q_a\to\mathfrak X\to X_{\mathrm{dR}/S}
\]
is $\eta_{X/S}$.  Functoriality of effective-image factorization therefore
produces the map $Q_a\to Q_{X/S}$.  The anchored quotient classification
gives the corresponding anchor, and
Theorem~\ref{thm:unitwise-formal-groupoid-completion} gives its canonical
formalization.  Finally,
Theorem~\ref{thm:effective-relation-calculus} applied to
$a:X\to\mathfrak X$ identifies its relation, monad, comonad, and transport
calculus with those of the effective image
$\bar a:X\twoheadrightarrow Q_a$.
\end{proof}

\begin{remark}[Effective versus full target]
\label{rem:effective-full-leaf}
The kernel-pair groupoid sees $Q_a$, not arbitrary excess in
$\mathfrak X$.  Its formalization may further replace $Q_a$ by the
coreflective formal quotient $\widehat Q_{\bar a}$.  Thus three objects may
differ:
\[
    \mathfrak X,
    \qquad
    Q_a,
    \qquad
    \widehat Q_{\bar a}.
\]
Any later comparison with a represented or formal-moduli leaf space must
therefore identify which of these objects it models and prove the relevant
comparison rather than building that identification into the present
definition \cite{Nuiten2019,BrantnerMagidsonNuiten2025}.
\end{remark}


\section{Effective descent and formal Cartan objects}
\label{sec:flg-descent}

\begin{theorem}[Formal-groupoid transport and mate compatibility]
\label{thm:formal-groupoid-transport-mates}
For $E\to X$, the following data are canonically equivalent:
\begin{enumerate}
    \item a $T_{\mathfrak G}$-algebra structure on $E$;
    \item a $J_{\mathfrak G}$-coalgebra structure on $E$;
    \item coherent descent transport along $\mathfrak G_\bullet$, whose
    degree-one term is an equivalence
    \[
        \theta:s^*E\xrightarrow{\sim}t^*E
    \]
    satisfying
    \[
        u^*\theta\simeq\operatorname{id}_E
    \]
    and the Cech cocycle identity on $\mathfrak G_2$.
\end{enumerate}
Under these equivalences the algebra action and coalgebra coaction are mates
under
$T_{\mathfrak G}\dashv J_{\mathfrak G}$.
\end{theorem}

\begin{proof}
By groupoid effectivity,
\[
    \mathfrak G_\bullet
    \simeq
    \check C_\bullet(q).
\]
Apply Theorem~\ref{thm:effective-relation-calculus} to
$q:X\twoheadrightarrow Q_{\mathfrak G}$.  Its relation groupoid is
$\mathfrak G_\bullet$, and its monad and comonad are precisely
\[
    T_{\mathfrak G}=q^*\Sigma_q,
    \qquad
    J_{\mathfrak G}=q^*\Pi_q.
\]
The theorem therefore gives the stated equivalence between algebra structure,
coalgebra structure, and coherent transport, including the unit, cocycle, and
all higher \v{C}ech coherences.  It also identifies the algebra action and
coalgebra coaction as mates under
$T_{\mathfrak G}\dashv J_{\mathfrak G}$.
\end{proof}

\begin{theorem}[Effective formal-groupoid descent]
\label{thm:formal-groupoid-effective-descent}
Pullback along $q:X\twoheadrightarrow Q_{\mathfrak G}$ induces canonical
equivalences
\[
    \mathcal X_{/Q_{\mathfrak G}}
    \simeq
    \operatorname{Alg}_{T_{\mathfrak G}}(\mathcal X_{/X})
    \simeq
    \operatorname{Coalg}_{J_{\mathfrak G}}(\mathcal X_{/X}).
\]
Both equivalences commute with the forgetful functors to
$\mathcal X_{/X}$.
\end{theorem}

\begin{proof}
Apply Theorem~\ref{thm:effective-relation-calculus} to the effective quotient
\[
    q:X\twoheadrightarrow Q_{\mathfrak G}.
\]
Groupoid effectivity identifies its \v{C}ech nerve with
$\mathfrak G_\bullet$, so the theorem gives simultaneously
\[
    \mathcal X_{/Q_{\mathfrak G}}
    \simeq
    \operatorname*{Tot}_{[n]\in\Delta}
    \mathcal X_{/\mathfrak G_n}
    \simeq
    \operatorname{Alg}_{T_{\mathfrak G}}(\mathcal X_{/X})
    \simeq
    \operatorname{Coalg}_{J_{\mathfrak G}}(\mathcal X_{/X}),
\]
with the common underlying functor given by pullback along $q$.
\end{proof}

\begin{definition}[Formal Cartan objects]
\label{def:formal-cartan-objects}
Define the $\infty$-category of formal Cartan objects of $\mathfrak G$ by
\[
    \operatorname{Cart}(\mathfrak G)
    :=
    \mathcal X_{/Q_{\mathfrak G}}.
\]
Through Theorem~\ref{thm:formal-groupoid-effective-descent}, it may equivalently
be presented as
\[
    \operatorname{Alg}_{T_{\mathfrak G}}(\mathcal X_{/X})
    \quad\text{or}\quad
    \operatorname{Coalg}_{J_{\mathfrak G}}(\mathcal X_{/X}).
\]
\end{definition}

\begin{remark}[Cartan structure is derived]
A formal Cartan object is intrinsically an object over the formal leaf space.
Its disk action and jet coaction are derived by effective descent and are not
independent structures.  This is the exact analogue, for a chosen formal Lie
groupoid, of the intrinsic Cartan philosophy of the maximal jet chapter.
\end{remark}

\subsection{Invariant and descended sections}

Let $\bar E\to Q_{\mathfrak G}$ and set $E:=q^*\bar E$.

\begin{definition}[Coherent $\mathfrak G$-invariant sections]
\label{def:formal-groupoid-invariant-sections}
Define
\[
    \Gamma_{\mathfrak G}(E)
    :=
    \operatorname{Map}_{\mathcal X_{/Q_{\mathfrak G}}}
    \left(
        \operatorname{id}_{Q_{\mathfrak G}},
        \bar E
    \right).
\]
\end{definition}

\begin{proposition}[Cech presentation of invariant sections]
\label{prop:invariant-sections-totalization}
There is a canonical equivalence
\[
    \Gamma_{\mathfrak G}(E)
    \simeq
    \operatorname*{Tot}_{[n]\in\Delta}
    \operatorname{Map}_{\mathcal X_{/\mathfrak G_n}}
    \left(
        \operatorname{id}_{\mathfrak G_n},
        E|_{\mathfrak G_n}
    \right),
\]
where, for the zeroth vertex map
$p_0:\mathfrak G_n\to X$, we write
\[
    E|_{\mathfrak G_n}:=p_0^*E.
\]
The coherent descent transport canonically identifies this object with the
pullback along every other vertex and induces the displayed cosimplicial
structure.
\end{proposition}

\begin{proof}
Apply effective descent to mapping spaces in the slice over
$Q_{\mathfrak G}$.  Mapping spaces in a limit of $\infty$-categories are the
corresponding limits of mapping spaces.
\end{proof}


\section{Functoriality and base change}
\label{sec:flg-functoriality}

\subsection{Morphisms over a fixed object}

Let
\[
    F:\mathfrak G_\bullet\longrightarrow\mathfrak H_\bullet
\]
be a morphism of formal fixed-object groupoids on $X$.  Since
$R_{X/S,\bullet}$ is terminal among formal fixed-object groupoids,
\[
    \phi_{\mathfrak H}F
    \simeq
    \phi_{\mathfrak G}
\]
through a contractible space of choices.  Realization therefore gives a map
\[
    a:=Q_F:
    Q_{\mathfrak G}\longrightarrow Q_{\mathfrak H}
\]
under $X$ and over $Q_{X/S}$, with
\[
    q_{\mathfrak H}\simeq aq_{\mathfrak G}.
\]

\begin{proposition}[Functoriality over fixed $X$]
\label{prop:formal-gpd-functoriality-fixed-X}
A morphism
$F:\mathfrak G\to\mathfrak H$
induces canonically:
\begin{enumerate}
    \item for every $b:T\to X$, a morphism of formal orbit disks
    \[
        \mathbb D^{\mathfrak G}_b
        \longrightarrow
        \mathbb D^{\mathfrak H}_b;
    \]
    \item homomorphisms
    \[
        I^{\mathrm{raw}}_{\mathfrak G}
        \longrightarrow
        I^{\mathrm{raw}}_{\mathfrak H},
        \qquad
        I^{\mathrm{for}}_{\mathfrak G}
        \longrightarrow
        I^{\mathrm{for}}_{\mathfrak H};
    \]
    \item over $Q_{\mathfrak G}$, homomorphisms
    \[
        \mathcal I_{Q_{\mathfrak G}/S}
        \longrightarrow
        a^*\mathcal I_{Q_{\mathfrak H}/S},
        \qquad
        \mathcal K_{\mathfrak G}
        \longrightarrow
        a^*\mathcal K_{\mathfrak H};
    \]
    \item a natural transformation of monads
    \[
        T_{\mathfrak G}
        \longrightarrow
        T_{\mathfrak H};
    \]
    \item a natural transformation of comonads
    \[
        J_{\mathfrak H}
        \longrightarrow
        J_{\mathfrak G}.
    \]
\end{enumerate}
These constructions preserve identities, composition, commutators,
adjoint transport, bitorsor structures, and all higher functorial
coherences.
\end{proposition}

\begin{proof}
The quotient map $a$ induces
\[
    T\times_{Q_{\mathfrak G}}X
    \longrightarrow
    T\times_{Q_{\mathfrak H}}X,
\]
giving the disk map.  Raw isotropy is a finite-limit construction on
$F_1$, and formal isotropy is its centred completion, so
Proposition~\ref{prop:centred-completion-varying-centre} gives the two
isotropy homomorphisms.  Equivalently, formal isotropy functoriality follows
from the anchor-kernel formula.  Relative inertia is functorial in the
quotient, and the square over $Q_{X/S}$ carries the identity-loop fibre
defining $\mathcal K_{\mathfrak G}$ to that defining
$\mathcal K_{\mathfrak H}$.

For the disk monads,
\[
\begin{aligned}
    T_{\mathfrak H}
    &=
    q_{\mathfrak H}^*\Sigma_{q_{\mathfrak H}}\\
    &\simeq
    q_{\mathfrak G}^*a^*\Sigma_a\Sigma_{q_{\mathfrak G}}.
\end{aligned}
\]
Apply
$q_{\mathfrak G}^*(-)\Sigma_{q_{\mathfrak G}}$
to the unit
\[
    \operatorname{id}\longrightarrow a^*\Sigma_a
\]
of $\Sigma_a\dashv a^*$.

For jets,
\[
\begin{aligned}
    J_{\mathfrak H}
    &=
    q_{\mathfrak H}^*\Pi_{q_{\mathfrak H}}\\
    &\simeq
    q_{\mathfrak G}^*a^*\Pi_a\Pi_{q_{\mathfrak G}}.
\end{aligned}
\]
Apply
$q_{\mathfrak G}^*(-)\Pi_{q_{\mathfrak G}}$
to the counit
\[
    a^*\Pi_a\longrightarrow\operatorname{id}
\]
of $a^*\dashv\Pi_a$.  Compatibility with the monad and comonad structures
is naturality of units, counits, and composition of dependent adjunctions.
The remaining structures are formed functorially from groupoid
multiplication, inversion, finite limits, and the canonical anchor.
\end{proof}

\subsection{Functoriality in the geometric object}

\begin{definition}[Functor of formal Lie groupoids]
\label{def:functor-formal-lie-groupoids}
Let $f:X\to Y$ be a morphism over $S$.  A functor
\[
    (\mathfrak G,X)\longrightarrow(\mathfrak H,Y)
\]
of formal Lie groupoids over $S$ is a simplicial functor
\[
    F_\bullet:
    \mathfrak G_\bullet\longrightarrow\mathfrak H_\bullet
\]
whose degree-zero component is $f$.
\end{definition}

\begin{proposition}[Automatic anchor compatibility]
\label{prop:automatic-anchor-functoriality}
Every functor of formal Lie groupoids as above fits canonically into a
homotopy-commutative square
\[
\begin{tikzcd}
\mathfrak G_\bullet \arrow[r,"F_\bullet"] \arrow[d,"\phi_{\mathfrak G}"'] &
\mathfrak H_\bullet \arrow[d,"\phi_{\mathfrak H}"]\\
R_{X/S,\bullet} \arrow[r,"R(f)"'] &
R_{Y/S,\bullet}.
\end{tikzcd}
\]
The construction supplies a canonical compatible homotopy, functorial in
$(F,f)$ and coherently compatible with identities and composition.
\end{proposition}

\begin{proof}
For every fixed-object groupoid there is a canonical vertex morphism to the
relative pair groupoid.  Hence the square
\[
\begin{tikzcd}
\mathfrak G_\bullet \arrow[r,"F_\bullet"] \arrow[d] &
\mathfrak H_\bullet \arrow[d]\\
\operatorname{Pair}(X/S)_\bullet
\arrow[r,"\operatorname{Pair}(f)"'] &
\operatorname{Pair}(Y/S)_\bullet
\end{tikzcd}
\]
commutes by functoriality of the vertex maps.  Apply the varying-centre
completion of Proposition~\ref{prop:centred-completion-varying-centre}.
Using
Theorem~\ref{thm:pair-groupoid-completion} identifies the bottom completed
map with $R(f)$ and the vertical completed maps with the canonical anchors.
\end{proof}

\begin{proposition}[Quotient functoriality]
\label{prop:quotient-functoriality}
A functor
$(\mathfrak G,X)\to(\mathfrak H,Y)$
induces a canonical morphism
\[
    Q_{\mathfrak G}\longrightarrow Q_{\mathfrak H}
\]
compatible with
\[
    Q_{X/S}\longrightarrow Q_{Y/S}.
\]
The assignment preserves identities, composition, and all higher
functorial coherences.
\end{proposition}

\begin{proof}
Apply geometric realization to the simplicial functor and to the canonical
anchor square of
Proposition~\ref{prop:automatic-anchor-functoriality}.
\end{proof}

\subsection{Arbitrary change of spectral base}

Let $c:S'\to S$, put
\[
    X':=X\times_SS',
    \qquad
    \mathfrak G'_\bullet
    :=
    \mathfrak G_\bullet\times_SS'.
\]

\begin{theorem}[Base change of formal Lie groupoids]
\label{thm:formal-gpd-base-change}
There are canonical equivalences
\[
    Q_{\mathfrak G'}
    \simeq
    Q_{\mathfrak G}\times_SS',
    \qquad
    R_{X'/S',\bullet}
    \simeq
    R_{X/S,\bullet}\times_SS',
\]
under which $\mathfrak G'_\bullet$ is a formal Lie groupoid with the
base-changed canonical anchor.  Moreover,
\[
    \widehat{\mathfrak G'}^{\,\mathrm{dR}}
    \simeq
    \widehat{\mathfrak G}^{\,\mathrm{dR}}\times_SS',
\]
\[
    I^{\mathrm{raw}}_{\mathfrak G'}
    \simeq
    I^{\mathrm{raw}}_{\mathfrak G}\times_SS',
    \qquad
    I^{\mathrm{for}}_{\mathfrak G'}
    \simeq
    I^{\mathrm{for}}_{\mathfrak G}\times_SS',
\]
and
\[
    \mathcal K_{\mathfrak G'}
    \simeq
    \mathcal K_{\mathfrak G}\times_SS'.
\]
For every generalized point $a:T\to X$ and its base change
$a':T':=T\times_SS'\to X'$,
\[
    \mathbb D^{\mathfrak G'}_{a'}
    \simeq
    \mathbb D^{\mathfrak G}_a\times_SS'.
\]
The quotient formalization, disk monads, jet comonads, adjunction, transport
data, bitorsor structures, adjoint transport, and effective descent
equivalences commute with base change through their canonical
Beck--Chevalley maps.
\end{theorem}

\begin{proof}
Pullback between slices of an $\infty$-topos preserves finite limits and all
small colimits.  It therefore preserves groupoid realizations, effective
epimorphisms, monomorphisms, and the iterated fibre products of \v{C}ech nerves.

The relative de Rham base-change theorem and
Proposition~\ref{prop:centred-completion-base-change} give
\[
    \widehat{\mathfrak G'}^{\,\mathrm{dR}}
    \simeq
    \widehat{\mathfrak G}^{\,\mathrm{dR}}\times_SS'.
\]
Thus formality is preserved.  Base change for $R_{X/S,\bullet}$ follows
either from the inherited theorem of the preceding chapter or from
Theorem~\ref{thm:pair-groupoid-completion}.  Geometric realization gives
the quotient formula.

Raw isotropy, relative inertia, formal disks, and the two arrow images are
finite-limit or effective-image constructions, hence commute with pullback.
Formal isotropy commutes with base change by centred-completion base change,
and the formula for $\mathcal K$ follows from finite limits and the
base-changed quotient map.  Effective-image factorization also gives
\[
    \widehat Q_{q'}
    \simeq
    \widehat Q_q\times_SS'
\]
for the quotient formalization.  The dependent-adjoint and descent
statements are the canonical Beck--Chevalley and effective-descent
comparisons.
\end{proof}


\section{Fixed-object quotient invariance and atlas dependence}
\label{sec:flg-presentation-invariance}

The quotient classification gives complete presentation invariance while the
object space $X$ is fixed.  It does not give arbitrary-atlas Morita
invariance, because the formality condition itself depends on the chosen
centre.

\begin{definition}[Fixed-object quotient equivalence]
\label{def:quotient-equivalence}
A morphism
\[
    F:\mathfrak G\to\mathfrak H
\]
in $\operatorname{FormLieGpd}(X/S)$ is a
\emph{fixed-object quotient equivalence} if
\[
    Q_F:
    Q_{\mathfrak G}\longrightarrow Q_{\mathfrak H}
\]
is an equivalence.
\end{definition}

\begin{proposition}[Fixed-object presentation invariance]
\label{prop:presentation-invariance}
A fixed-object quotient equivalence is an equivalence of formal Lie
groupoids.  It induces equivalences of the formal orbit disks, raw and formal
isotropy groups, formal inertia kernels, raw and formal arrow images, adjoint
transport, bitorsors, disk monads, jet comonads, and Cartan categories.
\end{proposition}

\begin{proof}
By groupoid effectivity,
\[
    \mathfrak G_\bullet
    \simeq
    \check C_\bullet(X\to Q_{\mathfrak G}),
    \qquad
    \mathfrak H_\bullet
    \simeq
    \check C_\bullet(X\to Q_{\mathfrak H}).
\]
An equivalence of quotient objects under $X$ gives degreewise equivalences of
the \v{C}ech nerves.  Every remaining construction is obtained functorially from
the quotient, its relation groupoid, finite limits, effective images, centred
completion, or dependent adjunction.
\end{proof}

\begin{proposition}[Unitwise formality and formalization are not arbitrary-atlas invariant]
\label{prop:formality-not-morita-invariant}
Unitwise relative de Rham-completeness is not preserved by arbitrary
replacement of the object space by another effective atlas of the same
quotient.  Moreover the canonical unitwise formalization constructed in this
chapter does not restore arbitrary-atlas Morita invariance.
\end{proposition}

\begin{proof}
Let $X$ be noninitial and start with the zero formal groupoid
\[
    \check C_\bullet(\operatorname{id}_X).
\]
Change atlas by the split twofold effective cover
\[
    u:
    U:=X\amalg X
    \twoheadrightarrow
    X.
\]
The resulting groupoid is $\check C_\bullet(u)$.

By Theorem~\ref{thm:formal-quotient-monomorphism-recognition}, this groupoid is
formal if and only if
\[
    L_S(u):
    L_S(U)\longrightarrow L_S(X)
\]
is a monomorphism.  The preceding de Rham chapter proves finite-coproduct
compatibility, so
\[
    L_S(U)
    \simeq
    L_S(X)\amalg L_S(X),
\]
and $L_S(u)$ is the fold map.  We claim that $L_S(X)$ is noninitial.  If
$L_S(X)\simeq\varnothing$, then the relative de Rham unit would supply a
morphism
\[
    X\longrightarrow L_S(X)\simeq\varnothing.
\]
The initial object of an $\infty$-topos is strict, so this would force
$X\simeq\varnothing$, contrary to the choice of $X$.  Thus $L_S(X)$ is
noninitial.

Coproducts in an $\infty$-topos are disjoint.  The self-pullback of the fold
map
\[
    L_S(X)\amalg L_S(X)
    \longrightarrow
    L_S(X)
\]
contains the two off-diagonal coproduct summands, and these are noninitial.
Its diagonal is therefore not an equivalence.  Hence the fold map is not
monic and $\check C_\bullet(u)$ is not formal.

We now compute its canonical formalization.  The centred quotient completion
of $u:U\twoheadrightarrow X$ is
\[
\begin{aligned}
    P_u
    &=
    X\times_{L_S(X)}L_S(U)\\
    &\simeq
    X\times_{L_S(X)}
    \bigl(L_S(X)\amalg L_S(X)\bigr)\\
    &\simeq
    \bigl(X\times_{L_S(X)}L_S(X)\bigr)
    \amalg
    \bigl(X\times_{L_S(X)}L_S(X)\bigr)\\
    &\simeq
    X\amalg X
    =
    U.
\end{aligned}
\]
Under this equivalence, the completed centre
\[
    \delta_u=(u,\eta_U):U\longrightarrow P_u
\]
restricts to the identity on each coproduct component and is therefore an
equivalence.  Its effective image factorization is consequently
\[
    U
    \xrightarrow{\operatorname{id}_U}
    U
    \xrightarrow{\sim}
    P_u.
\]
Thus
\[
    \widehat Q_u\simeq U,
    \qquad
    \widehat u\simeq\operatorname{id}_U,
    \qquad
    c_u\simeq u.
\]
The quotient formula for unitwise completion now gives
\[
    \widehat{\check C_\bullet(u)}^{\,\mathrm{dR}}
    \simeq
    \check C_\bullet(\operatorname{id}_U)
    =
    0_{U,\bullet}.
\]
Its quotient is $U$, not $X$.  Hence the chapter's canonical formalization,
although a coreflection at fixed object space, does not turn this atlas
replacement back into the original zero groupoid on $X$.
\end{proof}

\begin{remark}[Scope of presentation invariance]
The constructions are presentation-independent under equivalences of quotient
presentations with fixed object space $X$.  No stronger arbitrary-atlas
invariance is asserted.  Unitwise formalization is itself a fixed-centre
coreflection and, as the preceding calculation shows, is not a Morita
replacement procedure.  If a later recognition theorem uses another atlas,
its comparison with the present formal groupoid must therefore construct and
prove the relevant comparison rather than identify the two presentations
solely from a common quotient.
\end{remark}


\section{Basic geometric examples}
\label{sec:flg-examples}

\subsection{The zero formal Lie groupoid}

\begin{example}[Zero groupoid]
The identity groupoid
\[
    0_{X,\bullet}
    =
    \check C_\bullet(\operatorname{id}_X)
\]
is a formal Lie groupoid with quotient $X$.  For every $a:T\to X$,
\[
    \mathbb D^{0_X}_a\simeq T,
    \qquad
    T_{0_X}\simeq\operatorname{id},
    \qquad
    J_{0_X}\simeq\operatorname{id}.
\]
Its raw isotropy is the relative inertia of $X$,
\[
    I^{\mathrm{raw}}_{0_X}
    \simeq
    \mathcal I_{X/S},
\]
while its formal isotropy is
\[
    I^{\mathrm{for}}_{0_X}
    \simeq
    \mathcal I_{X/S}
    \times_{\epsilon_{X/S}^*\mathcal I_{Q_{X/S}/S}}
    X.
\]
\end{example}

\begin{proof}
The quotient is $\operatorname{id}_X$, so the orbit disk and the two
dependent-adjoint composites are identities.  The raw isotropy formula is
Proposition~\ref{prop:isotropy-inertia}.  Here the formal quotient map
$r:X\to Q_{X/S}$ is $\epsilon_{X/S}$, so
Theorem~\ref{thm:formal-isotropy-descent} gives the displayed formal
isotropy kernel.
\end{proof}

\begin{remark}
The raw isotropy of the zero groupoid need not be trivial in an
$\infty$-topos: it records the ambient relative inertia of the object $X$.
The term ``zero groupoid'' refers to the simplicial relation structure, not
to the vanishing of ambient higher loops.
\end{remark}

\subsection{The maximal relative de Rham groupoid}

\begin{example}[Maximal groupoid]
The terminal formal Lie groupoid is
\[
    R_{X/S,\bullet}.
\]
Its formal leaf space is $Q_{X/S}$, its formal orbit disks are
\[
    \mathbb D^\infty_{X/S,a},
\]
and
\[
    T_{R_{X/S}}=T^\infty_{X/S},
    \qquad
    J_{R_{X/S}}=J^\infty_{X/S}.
\]
Its raw isotropy is
\[
    I^{\mathrm{raw}}_{R_{X/S}}
    \simeq
    \epsilon_{X/S}^*\mathcal I_{Q_{X/S}/S},
\]
whereas its formal isotropy is the trivial group object
\[
    I^{\mathrm{for}}_{R_{X/S}}
    \simeq X.
\]
\end{example}

\begin{proof}
The quotient and jet assertions are immediate from the definitions.  The raw
isotropy formula is
Proposition~\ref{prop:isotropy-inertia}.  The canonical anchor of the maximal
relation is the identity, so the anchor-kernel formula gives
\[
    I^{\mathrm{for}}_{R_{X/S}}
    \simeq
    R_{X/S}\times_{R_{X/S}}X
    \simeq X.
\]
\end{proof}

\subsection{Formal completion of the pair groupoid}

\begin{example}[Pair groupoid]
For the relative pair groupoid
\[
    \operatorname{Pair}(X/S)_n
    =
    X^{\times_S(n+1)},
\]
Theorem~\ref{thm:pair-groupoid-completion} gives
\[
    \widehat{\operatorname{Pair}(X/S)}^{\,\mathrm{dR}}_\bullet
    \simeq
    R_{X/S,\bullet}.
\]
Thus the maximal infinitesimal relation is precisely the unitwise
relative de Rham completion of the codiscrete relative groupoid.
\end{example}

\subsection{Formal action groupoids}

Let $G\to S$ be a group object and define
\[
    \mu^{\mathrm{mod}}_{G/S}
    :=
    G\times_{G_{\mathrm{dR}/S}}S,
\]
where $S\to G_{\mathrm{dR}/S}$ is the relative de Rham image of the identity
section.

\begin{proposition}[Constructed formal identity group]
\label{prop:formal-identity-group}
The object $\mu^{\mathrm{mod}}_{G/S}$ carries a canonical group-object
structure over $S$.  The projection
\[
    \mu^{\mathrm{mod}}_{G/S}\longrightarrow G
\]
is a group homomorphism, $\mu^{\mathrm{mod}}_{G/S}$ is complete along its
identity section, and the construction is compatible with arbitrary base
change in $S$.
\end{proposition}

\begin{proof}
Relative de Rham localization is left exact, hence carries the finite-limit
diagram defining the group object $G$ to a group object
$G_{\mathrm{dR}/S}$, and naturality makes
\[
    \eta_G:G\longrightarrow G_{\mathrm{dR}/S}
\]
a group homomorphism.  The terminal object $S$ of $\mathcal X_{/S}$ carries
its unique group structure, and the map
$S\to G_{\mathrm{dR}/S}$ is the identity section of the localized group.
Therefore the pullback
\[
    \mu^{\mathrm{mod}}_{G/S}
    =
    G\times_{G_{\mathrm{dR}/S}}S
\]
is a group object, with projection to $G$ a homomorphism.

Moreover
\[
    \mu^{\mathrm{mod}}_{G/S}
    =
    C_S(G,e).
\]
Theorem~\ref{thm:centred-completion-idempotent} therefore makes it complete
along the identity section.  Base-change compatibility is
Proposition~\ref{prop:centred-completion-base-change}, together with
base change of the group-object finite-limit diagrams.
\end{proof}

Now let
\[
    \alpha:G\times_SX\longrightarrow X
\]
be a left action.  Write $G\ltimes X$ for the action groupoid, with
\[
    s(g,x)=x,
    \qquad
    t(g,x)=\alpha(g,x).
\]

\begin{theorem}[Completion of an action groupoid]
\label{thm:action-groupoid-completion}
The $G$-action precomposes canonically along the homomorphism
$\mu^{\mathrm{mod}}_{G/S}\to G$ to an action
\[
    \mu^{\mathrm{mod}}_{G/S}\times_SX
    \longrightarrow X,
\]
and there is a canonical equivalence
\[
    \widehat{(G\ltimes X)}^{\,\mathrm{dR}}
    \simeq
    \mu^{\mathrm{mod}}_{G/S}\ltimes X.
\]
\end{theorem}

\begin{proof}
Proposition~\ref{prop:formal-identity-group} gives a canonical group
homomorphism
\[
    \mu^{\mathrm{mod}}_{G/S}\longrightarrow G.
\]
Precomposing the $G$-action with this homomorphism constructs
\[
    \mu^{\mathrm{mod}}_{G/S}\times_SX\longrightarrow X,
\]
and its unit, associativity, and higher coherences are the pullbacks of those
of the original action.

The degree-one centre is
\[
    X\longrightarrow G\times_SX,
    \qquad
    x\longmapsto(e,x).
\]
By left exactness,
\[
    (G\times_SX)_{\mathrm{dR}/S}
    \simeq
    G_{\mathrm{dR}/S}
    \times_S
    X_{\mathrm{dR}/S}.
\]
Hence
\[
\begin{aligned}
    C_X(G\times_SX)
    &\simeq
    (G\times_SX)
    \times_{
        G_{\mathrm{dR}/S}\times_SX_{\mathrm{dR}/S}
    }
    X_{\mathrm{dR}/S}\\
    &\simeq
    \left(
        G\times_{G_{\mathrm{dR}/S}}S
    \right)
    \times_SX\\
    &=
    \mu^{\mathrm{mod}}_{G/S}\times_SX.
\end{aligned}
\]
The completed target map is the action obtained by precomposition along
$\mu^{\mathrm{mod}}_{G/S}\to G$.  Pullback preservation of centred
completion carries the Segal pullbacks of the action nerve to the Segal
pullbacks of this precomposed action, and hence identifies the completed
simplicial groupoid with $\mu^{\mathrm{mod}}_{G/S}\ltimes X$.
\end{proof}

\subsection{Formal gauge groupoids}

\begin{definition}[Classifying object of a group object]
\label{def:classifying-object-group}
Let $B_\bullet G$ be the one-object groupoid with
\[
    (B_\bullet G)_0=S,
    \qquad
    (B_\bullet G)_n=G^{\times_S n}
    \quad(n\geq1).
\]
Define
\[
    BG:=|B_\bullet G|.
\]
By groupoid effectivity, the canonical basepoint
\[
    b_G:S\twoheadrightarrow BG
\]
is an effective epimorphism and
\[
    \check C_\bullet(b_G)\simeq B_\bullet G.
\]
We call $b_G$ the \emph{universal classifying basepoint}.  For every base
change $c:S'\to S$, with $G':=G\times_SS'$, pullback preserves the finite
limits defining $B_\bullet G$ and the geometric realization, hence gives
canonical equivalences
\[
    B_\bullet G'
    \simeq
    B_\bullet G\times_SS',
    \qquad
    BG'\simeq BG\times_SS',
\]
under which $b_{G'}$ is the pullback of $b_G$.
\end{definition}

\begin{definition}[Constructed principal right group object]
\label{def:principal-group-object}
Let $\chi:X\to BG$ be a morphism.  Define
\[
    P_\chi
    :=
    X\times_{BG}S
\]
and write
\[
    p_\chi:P_\chi\longrightarrow X
\]
for the first projection.  Pulling back
$\check C_\bullet(b_G)\simeq B_\bullet G$ along $\chi$ gives canonical
equivalences
\[
    P_\chi\times_XP_\chi
    \simeq
    P_\chi\times_SG
\]
and their higher iterates.  Under the degree-one equivalence, the second
projection
$P_\chi\times_XP_\chi\to P_\chi$ defines a right action
\[
    \alpha_\chi:
    P_\chi\times_SG\longrightarrow P_\chi.
\]
The higher pulled-back \v{C}ech identities supply its unit, associativity, and
all higher action coherences.  We call
$P_\chi\to X$ with this constructed action the \emph{principal right
$G$-object classified by $\chi$}.
\end{definition}

\begin{theorem}[Construction and recognition of principal right group objects]
\label{prop:principal-classifying-map}
For every $\chi:X\to BG$, the constructed principal right $G$-object
$p_\chi:P_\chi\to X$ satisfies:
\begin{enumerate}
    \item $p_\chi$ is an effective epimorphism;
    \item the shear morphism
    \[
        (\operatorname{pr}_1,\alpha_\chi):
        P_\chi\times_SG
        \xrightarrow{\sim}
        P_\chi\times_XP_\chi
    \]
    is an equivalence.
\end{enumerate}

Conversely, let $p:P\to X$ carry a coherent right $G$-action.  If $p$ is an
effective epimorphism and its shear morphism
\[
    (\operatorname{pr}_1,\alpha):
    P\times_SG
    \longrightarrow
    P\times_XP
\]
is an equivalence, then the action constructs canonically a classifying
morphism
\[
    \chi_P:X\longrightarrow BG
\]
and a canonical equivalence of right $G$-objects over $X$
\[
    P\xrightarrow{\sim}P_{\chi_P}
    =
    X\times_{BG}S.
\]
Thus effective epimorphism and invertible shear are recognition properties of
the constructed principal objects, not additional principal-object axioms.
The construction is natural under pullback in $X$ and arbitrary base change
in $S$.
\end{theorem}

\begin{proof}
For a given $\chi$, the projection
\[
    p_\chi:
    X\times_{BG}S\longrightarrow X
\]
is the base change of the effective epimorphism
$b_G:S\twoheadrightarrow BG$, hence is effective.  Moreover
\[
\begin{aligned}
    P_\chi\times_XP_\chi
    &\simeq
    X\times_{BG}(S\times_{BG}S)\\
    &\simeq
    P_\chi\times_SG,
\end{aligned}
\]
using
$S\times_{BG}S\simeq G$.  By construction this equivalence is the shear
morphism, and its higher iterates are the pullbacks of the higher terms of
$\check C_\bullet(b_G)\simeq B_\bullet G$.

Conversely, let $P\to X$ have effective projection and invertible shear.
Write $P/\!/G$ for its action groupoid.  Iterating the shear equivalence gives
a canonical simplicial equivalence
\[
    P/\!/G
    \xrightarrow{\sim}
    \check C_\bullet(p).
\]
Since $p$ is effective,
\[
    |P/\!/G|\simeq X.
\]
Forgetting the $P$-coordinate defines a simplicial morphism
\[
    P/\!/G\longrightarrow B_\bullet G.
\]
Realization constructs
\[
    \chi_P:X\longrightarrow BG.
\]

In degree zero the simplicial morphism gives the structural map $P\to S$,
and compatibility with realization gives a commutative square
\[
\begin{tikzcd}
P \arrow[r] \arrow[d,"p"'] & S \arrow[d,"b_G"]\\
X \arrow[r,"\chi_P"'] & BG.
\end{tikzcd}
\]
Hence there is a canonical comparison
\[
    j_P:
    P\longrightarrow
    X\times_{BG}S
    =
    P_{\chi_P}.
\]
Pull this map back along the effective epimorphism $p:P\to X$.  On the source,
the shear equivalence gives
\[
    P\times_XP
    \simeq
    P\times_SG.
\]
On the target,
\[
\begin{aligned}
    P\times_XP_{\chi_P}
    &\simeq
    P\times_{BG}S\\
    &\simeq
    P\times_S(S\times_{BG}S)\\
    &\simeq
    P\times_SG,
\end{aligned}
\]
where the middle equivalence uses the commutative square just constructed.
Tracing the simplicial forgetful morphism identifies the pullback of $j_P$
with the shear equivalence.  It is therefore an equivalence.  Pullback along
an effective epimorphism is conservative in an $\infty$-topos, so
$j_P$ itself is an equivalence.  The construction respects the action because
both action structures are the degree-one structures of the same identified
\v{C}ech nerves.

Every step is formed by pullback, realization, groupoid effectivity, or the
canonical action nerve.  Naturality under pullback and arbitrary base change
follows from preservation of these constructions.
\end{proof}

\begin{definition}[Space of principal right group objects]
\label{def:principal-object-space}
For $X\to S$, let
\[
    \operatorname{Prin}_G(X)
\]
denote the maximal $\infty$-groupoid whose objects are coherent right
$G$-objects
\[
    p:P\longrightarrow X
\]
such that $p$ is an effective epimorphism and the shear morphism
\[
    (\operatorname{pr}_1,\alpha):
    P\times_SG
    \longrightarrow
    P\times_XP
\]
is an equivalence.  Morphisms are equivariant equivalences over $X$, with
all higher homotopies inherited from the $\infty$-category of coherent right
$G$-objects.
\end{definition}

\begin{theorem}[Universal classification of principal right group objects]
\label{thm:principal-objects-classified-by-bg}
Pullback of the universal classifying basepoint induces a canonical
equivalence of spaces
\[
\boxed{
    \operatorname{Map}_{\mathcal E}(X,BG)
    \xrightarrow{\sim}
    \operatorname{Prin}_G(X),
    \qquad
    \chi\longmapsto
    P_\chi=X\times_{BG}S.
}
\]
The equivalence is natural in $X$ and compatible with arbitrary base change
in $S$.  Under this equivalence, the inverse sends a recognized principal
right $G$-object $P\to X$ to the classifying morphism
$\chi_P:X\to BG$ constructed in
Theorem~\ref{prop:principal-classifying-map}.  Consequently the words
``universal'' and ``principal'' introduce no additional classification or
torsor axiom beyond the constructions and recognition properties already
proved.
\end{theorem}

\begin{proof}
Theorem~\ref{prop:principal-classifying-map} shows that pullback of
$b_G:S\twoheadrightarrow BG$ sends every
$\chi:X\to BG$ to an object of $\operatorname{Prin}_G(X)$.  Pullback is
functorial in $\chi$ and its higher homotopies, so this defines a morphism of
spaces
\[
    \Phi_X:
    \operatorname{Map}_{\mathcal E}(X,BG)
    \longrightarrow
    \operatorname{Prin}_G(X).
\]

Conversely, for a recognized principal right $G$-object $P\to X$, the same
theorem identifies the action groupoid with the \v{C}ech groupoid:
\[
    P/\!/G
    \xrightarrow{\sim}
    \check C_\bullet(P\to X).
\]
The projection $P\to S$ and the action define the simplicial forgetful
morphism
\[
    P/\!/G\longrightarrow B_\bullet G.
\]
Since $P\to X$ is effective, realization identifies
$|P/\!/G|$ with $X$, and realization of the forgetful morphism gives
\[
    \chi_P:X\longrightarrow BG.
\]
Equivariant equivalences over $X$ induce equivalences of action groupoids
compatible with the forgetful morphisms, so this construction is functorial
on the maximal $\infty$-groupoid and defines
\[
    \Psi_X:
    \operatorname{Prin}_G(X)
    \longrightarrow
    \operatorname{Map}_{\mathcal E}(X,BG).
\]

Let $\chi:X\to BG$.  The action groupoid of
$P_\chi=X\times_{BG}S$ is the pullback of
\[
    B_\bullet G
    \simeq
    \check C_\bullet(S\twoheadrightarrow BG)
\]
along $\chi$.  Its simplicial forgetful morphism to $B_\bullet G$ is the
pullback projection.  Geometric realization in an $\infty$-topos is
universal, so realizing this pulled-back \v{C}ech diagram recovers the
original morphism $\chi$.  Hence
\[
    \Psi_X\Phi_X(\chi)\simeq\chi
\]
naturally in $\chi$.

Now start with $P\in\operatorname{Prin}_G(X)$.  In the proof of
Theorem~\ref{prop:principal-classifying-map}, the classifying map $\chi_P$
comes with a canonical comparison
\[
    j_P:
    P\longrightarrow
    X\times_{BG}S
    =
    P_{\chi_P}.
\]
That theorem proves $j_P$ is an equivalence of right $G$-objects over $X$.
The construction of $j_P$ is functorial under equivariant equivalences and
their higher homotopies because it is induced by the simplicial action nerve
and the universal pullback square.  Therefore
\[
    \Phi_X\Psi_X(P)\simeq P
\]
naturally in $P$.

The two natural equivalences exhibit $\Phi_X$ and $\Psi_X$ as inverse
equivalences of spaces.  Naturality under pullback in $X$ and arbitrary base
change in $S$ follows from the corresponding functoriality already proved for
$b_G$, action groupoids, effective descent, and
Theorem~\ref{prop:principal-classifying-map}.
\end{proof}

\begin{definition}[Gauge groupoid]
\label{def:gauge-groupoid}
For a principal right $G$-object
\[
    P_\chi=X\times_{BG}S\longrightarrow X
\]
classified by $\chi:X\to BG$, define its gauge groupoid by
\[
    \operatorname{Gauge}(P_\chi)
    :=
    \check C_\bullet(\chi)
    \rightrightarrows X.
\]
Thus
\[
    \operatorname{Gauge}(P_\chi)_1
    =
    X\times_{BG}X.
\]
\end{definition}

\begin{proposition}[Effective quotient formula for gauge arrows]
\label{prop:gauge-arrow-quotient}
Let $G$ act diagonally on $P_\chi\times_SP_\chi$ by
\[
    (p_0,p_1)\cdot g:=(p_0g,p_1g).
\]
Define the effective quotient
\[
    (P_\chi\times_SP_\chi)/G
\]
to be the realization of this diagonal action groupoid.  Then there is a
canonical equivalence
\[
    (P_\chi\times_SP_\chi)/G
    \xrightarrow{\sim}
    X\times_{BG}X
    =
    \operatorname{Gauge}(P_\chi)_1.
\]
Under this equivalence the source, target, unit, inversion, and composition
are those of the \v{C}ech groupoid of $\chi$.
\end{proposition}

\begin{proof}
Effective descent along $b_G:S\twoheadrightarrow BG$ is an equivalence of
$\infty$-categories and therefore preserves finite limits.  By construction,
the pullback of $X\to BG$ is $P_\chi\to S$.  Hence the pullback of
\[
    X\times_{BG}X\longrightarrow BG
\]
is
\[
    P_\chi\times_SP_\chi\longrightarrow S,
\]
and its induced descent transport is the diagonal $G$-action.  The inverse of
effective descent therefore identifies $X\times_{BG}X$ with the effective
quotient of that diagonal action.  The simplicial structure maps descend from
the corresponding finite-limit maps.
\end{proof}

\begin{definition}[Formal gauge groupoid]
\label{def:formal-gauge-groupoid}
Define
\[
    \operatorname{Gauge}^{\mathrm{for}}(P_\chi)
    :=
    \widehat{\operatorname{Gauge}(P_\chi)}^{\,\mathrm{dR}}.
\]
\end{definition}

\begin{proposition}[Gauge groupoids give formal Lie groupoids]
\label{prop:gauge-formal-lie-groupoid}
The formal gauge groupoid is canonically an object of
\[
    \operatorname{FormLieGpd}(X/S),
\]
natural under pullback of the classifying morphism $\chi$ and arbitrary
spectral base change.
\end{proposition}

\begin{proof}
The gauge groupoid is the explicit \v{C}ech groupoid of
$\chi:X\to BG$.  Apply
Theorem~\ref{thm:unitwise-formal-groupoid-completion}.  Pullback of $\chi$
pulls back the universal principal object by
Definition~\ref{def:principal-group-object}, and arbitrary spectral base
change commutes with $BG$, its universal basepoint, the defining pullbacks,
and unitwise completion.  Naturality therefore follows from
Proposition~\ref{prop:centred-completion-varying-centre} and
Theorem~\ref{thm:formal-gpd-base-change}.
\end{proof}

\begin{remark}[Preview of the Atiyah sequence]
\label{rem:gauge-preview-atiyah}
The formal gauge groupoid is the nonlinear object to which the next
differentiation stage may apply in the geometric range where the relevant
cotangent or tangent controller has been constructed.  In the
characteristic-zero formal-moduli range this is parallel to the relation
between completed gauge groupoids and Atiyah Lie algebroids
\cite{Nuiten2019}.  The tangent fibre sequence and any resulting brave-new
Atiyah sequence are not asserted in the present chapter.
\end{remark}


\section{Group objects as the one-object special case}
\label{sec:flg-one-object-groups}

Let $G\to S$ be a group object and regard it as the one-object groupoid
\[
    B_\bullet G,
    \qquad
    (B_\bullet G)_0=S,
    \qquad
    (B_\bullet G)_1=G.
\]

\begin{theorem}[One-object completion]
\label{thm:one-object-completion}
The unitwise completion of $B_\bullet G$ is the reduced bar groupoid of
\[
    \mu^{\mathrm{mod}}_{G/S}
    :=
    G\times_{G_{\mathrm{dR}/S}}S.
\]
In particular
\[
    \widehat G^{\,\mathrm{dR}}_e
    \simeq
    \mu^{\mathrm{mod}}_{G/S}
\]
in degree one and
\[
    \widehat{B_\bullet G}^{\,\mathrm{dR}}
    \simeq
    B_\bullet\mu^{\mathrm{mod}}_{G/S}.
\]
\end{theorem}

\begin{proof}
The centre in degree one is the identity section
$e:S\to G$.  Since $S$ is terminal in $\mathcal X_{/S}$,
\[
    S_{\mathrm{dR}/S}\simeq S.
\]
Therefore
\[
    C_S(G,e)
    =
    G\times_{G_{\mathrm{dR}/S}}S
    =
    \mu^{\mathrm{mod}}_{G/S}.
\]
Pullback preservation of centred completion carries the higher Segal
pullbacks over $S$ to the reduced bar construction of this completed identity
group.
\end{proof}

\begin{corollary}[Recognition of one-object formality]
\label{cor:one-object-formality-recognition}
The one-object groupoid $B_\bullet G$ is formal if and only if
\[
    \mu^{\mathrm{mod}}_{G/S}
    \longrightarrow
    G
\]
is an equivalence.  Equivalently, the localized identity section
\[
    S
    \longrightarrow
    G_{\mathrm{dR}/S}
\]
is an equivalence.
\end{corollary}

\begin{proof}
The first criterion is
Proposition~\ref{prop:degree-one-formality} together with
Theorem~\ref{thm:one-object-completion}.  The second is
Proposition~\ref{prop:centred-completion-recognition} applied to
$e:S\to G$.
\end{proof}

\begin{corollary}[Classifying quotient of the formal identity group]
\label{cor:formal-identity-classifying-quotient}
The quotient of the completed one-object groupoid is
\[
    B\mu^{\mathrm{mod}}_{G/S},
\]
and the canonical classifying basepoint
\[
    S\twoheadrightarrow B\mu^{\mathrm{mod}}_{G/S}
\]
is an effective epimorphism whose \v{C}ech nerve is
$B_\bullet\mu^{\mathrm{mod}}_{G/S}$.  Pullback along this basepoint identifies
objects over $B\mu^{\mathrm{mod}}_{G/S}$ with objects over $S$ equipped with
coherent $\mu^{\mathrm{mod}}_{G/S}$-transport.
\end{corollary}

\begin{proof}
Apply groupoid effectivity to
$B_\bullet\mu^{\mathrm{mod}}_{G/S}$ for the effective epimorphism and \v{C}ech
nerve statement.  Then apply
Theorem~\ref{thm:effective-relation-calculus} to that effective epimorphism;
its relation groupoid is precisely the one-object groupoid
$B_\bullet\mu^{\mathrm{mod}}_{G/S}$, so its effective-descent equivalence is
the asserted classification by coherent transport.
\end{proof}

\begin{proposition}[One-object isotropy and adjoint structure]
\label{prop:one-object-isotropy}
For the completed one-object groupoid associated to $G/S$,
\[
    I^{\mathrm{raw}}
    \simeq
    I^{\mathrm{for}}
    \simeq
    \mu^{\mathrm{mod}}_{G/S}.
\]
The nonlinear formal-isotropy commutator is the ordinary internal group
commutator of $\mu^{\mathrm{mod}}_{G/S}$ and formal adjoint transport is
conjugation.
\end{proposition}

\begin{proof}
The object space is the terminal object $S$ of the slice, so the source-target
diagonal fibre is the entire arrow group
$\mu^{\mathrm{mod}}_{G/S}$.  This arrow group is already the output of centred
completion and is therefore formal.  The commutator and adjoint statements are
the general definitions specialized to a one-object groupoid.
\end{proof}

\begin{remark}
An arbitrary group object $G/S$ therefore determines a canonical one-object
formal Lie groupoid by completion.  Only when
$\mu^{\mathrm{mod}}_{G/S}\to G$ is an equivalence is the original
one-object groupoid itself formal.  Any later Lie-algebraic object is to be
obtained by differentiating the completed groupoid in a range where a
differentiation theorem has been established.
\end{remark}

\section{The modal relation of a group as an action groupoid}
\label{sec:flg-group-translation}

We now recover the nonlinear translation theory of a group object as a theorem
inside the general groupoid framework.

Apply the canonical relative infinitesimal groupoid construction to the object
$G\to S$:
\[
    R_{G/S}
    =
    G\times_{G_{\mathrm{dR}/S}}G
    \rightrightarrows G.
\]
Let
\[
    \mu:=\mu^{\mathrm{mod}}_{G/S}.
\]

\subsection{Translation trivializations}

\begin{theorem}[Group Translation Theorem]
\label{thm:group-translation}
Left and right multiplication define canonical equivalences over $G$
\[
\begin{aligned}
    \lambda_L:
    R_{G/S}
    &\xrightarrow{\sim}
    G\times_S\mu,
    &
    \lambda_L(g,h)&=(g,g^{-1}h),
\\
    \lambda_R:
    R_{G/S}
    &\xrightarrow{\sim}
    G\times_S\mu,
    &
    \lambda_R(g,h)&=(g,hg^{-1}).
\end{aligned}
\]
Their inverses are
\[
    \lambda_L^{-1}(g,u)=(g,gu),
    \qquad
    \lambda_R^{-1}(g,u)=(g,ug).
\]
\end{theorem}

\begin{proof}
If $(g,h)\in R_{G/S}$, then
$\eta_{G/S}(g)=\eta_{G/S}(h)$, and since the relative de Rham functor preserves
the group structure,
\[
    \eta_{G/S}(g^{-1}h)
    =
    \eta_{G/S}(g)^{-1}\eta_{G/S}(h)
    =
    e_{\mathrm{dR}}.
\]
Thus $g^{-1}h\in\mu$.  The same calculation gives
$hg^{-1}\in\mu$.  Conversely, if $u\in\mu$, then
\[
    \eta_{G/S}(gu)
    =
    \eta_{G/S}(g)e_{\mathrm{dR}}
    =
    \eta_{G/S}(g),
\]
so $(g,gu)\in R_{G/S}$, and similarly $(g,ug)\in R_{G/S}$.  The displayed
formulas are mutually inverse by associativity and the inverse identities.
\end{proof}

\begin{theorem}[Higher translation coordinates]
\label{thm:higher-translation-coordinates}
For every $n\geq0$ there is a canonical equivalence
\[
    R_{G/S,n}
    \xrightarrow{\sim}
    G\times_S\mu^{\times_S n}
\]
given by
\[
    (g_0,\ldots,g_n)
    \longmapsto
    \left(
        g_0,
        g_0^{-1}g_1,
        g_1^{-1}g_2,
        \ldots,
        g_{n-1}^{-1}g_n
    \right).
\]
Its inverse is
\[
    (g,u_1,\ldots,u_n)
    \longmapsto
    \left(
        g,
        gu_1,
        gu_1u_2,
        \ldots,
        gu_1\cdots u_n
    \right).
\]
Under these equivalences,
\[
    R_{G/S,\bullet}
    \simeq
    (G/\!/\mu)_\bullet
\]
for the right action $(g,u)\mapsto gu$.
\end{theorem}

\begin{proof}
Every tuple in $R_{G/S,n}$ has all vertices with the same relative de Rham
image, so each successive quotient
$g_{i-1}^{-1}g_i$ lies in $\mu$.  Conversely, multiplying successively by
elements of $\mu$ does not change the relative de Rham image.  The two formulas
are inverse by associativity.  In these coordinates, deleting an interior
vertex multiplies the two adjacent increments, deleting the initial vertex
acts by the first increment on $g$, deleting the terminal vertex forgets the
last increment, and degeneracies insert the identity.  These are exactly the
nerve maps of the right action groupoid.
\end{proof}

\begin{corollary}[Principal formal identity torsor]
\label{cor:principal-formal-identity-torsor}
There is a canonical quotient equivalence
\[
    \left|G/\!/\mu\right|
    \simeq
    Q_{G/S}.
\]
Thus
\[
    G\twoheadrightarrow Q_{G/S}
\]
is a principal right $\mu$-object in the ambient $\infty$-topos.
\end{corollary}

\begin{proof}
The map $G\twoheadrightarrow Q_{G/S}$ is an effective epimorphism by the
definition of the effective infinitesimal quotient.  By
Theorem~\ref{thm:higher-translation-coordinates}, the degree-one translation
equivalence identifies the shear map for the right $\mu$-action,
\[
    G\times_S\mu
    \longrightarrow
    G\times_{Q_{G/S}}G,
    \qquad
    (g,u)\longmapsto(g,gu),
\]
with an equivalence.  The recognition direction of
Theorem~\ref{prop:principal-classifying-map} therefore identifies
$G\to Q_{G/S}$ with the principal right $\mu$-object classified by the
realization of the action nerve.  The same theorem identifies
the entire action groupoid with the \v{C}ech groupoid of this effective
epimorphism, and realization gives
\[
    |G/\!/\mu|\simeq Q_{G/S}.
\]
\end{proof}

\subsection{The integrated Maurer--Cartan cocycle}

\begin{definition}[Formal Maurer--Cartan projections]
\label{def:formal-mc-projections}
Define
\[
    \vartheta_L:
    R_{G/S}\to\mu,
    \qquad
    \vartheta_L(g,h):=g^{-1}h,
\]
and
\[
    \vartheta_R:
    R_{G/S}\to\mu,
    \qquad
    \vartheta_R(g,h):=hg^{-1}.
\]
\end{definition}

\begin{theorem}[Integrated group-valued Maurer--Cartan cocycle]
\label{thm:integrated-mc-cocycle}
For every $n\geq0$, the successive left increments define
\[
    \Theta_{L,n}:
    R_{G/S,n}\longrightarrow B_n\mu=\mu^{\times_Sn}
\]
by
\[
    \Theta_{L,n}(g_0,\ldots,g_n)
    =
    \left(
        g_0^{-1}g_1,
        g_1^{-1}g_2,
        \ldots,
        g_{n-1}^{-1}g_n
    \right).
\]
These maps assemble into a simplicial morphism
\[
    \Theta_L:
    R_{G/S,\bullet}\longrightarrow B_\bullet\mu.
\]
Its degree-one component is $\vartheta_L$, and its realization is the
classifying morphism
\[
    Q_{G/S}\longrightarrow B\mu
\]
of the principal right $\mu$-object $G\to Q_{G/S}$.

The successive right increments similarly define a simplicial morphism
\[
    \Theta_R:
    R_{G/S,\bullet}
    \longrightarrow
    B_\bullet(\mu^{\mathrm{op}}).
\]
\end{theorem}

\begin{proof}
For left increments $u_i=g_{i-1}^{-1}g_i$, deleting an interior vertex $g_i$
replaces $(u_i,u_{i+1})$ by
\[
    g_{i-1}^{-1}g_{i+1}
    =
    (g_{i-1}^{-1}g_i)(g_i^{-1}g_{i+1})
    =
    u_iu_{i+1}.
\]
Repeating a vertex inserts the identity.  These are the reduced bar face and
degeneracy maps, so the maps are simplicial.  By
Theorem~\ref{thm:higher-translation-coordinates}, the source simplicial object
is the action groupoid of the principal right $\mu$-object
$G\to Q_{G/S}$.  The morphism $\Theta_L$ is exactly the simplicial
forgetful map from this action groupoid to $B_\bullet\mu$; hence its
realization is the classifying morphism constructed in
Proposition~\ref{prop:principal-classifying-map}.  For right increments the
multiplication order reverses, giving the opposite group.
\end{proof}

\begin{proposition}[Integrated curvature and Bianchi coherence]
\label{prop:integrated-curvature}
Define
\[
\begin{aligned}
    F_L(g_0,g_1,g_2)
    &:=
    \vartheta_L(g_0,g_1)
    \vartheta_L(g_1,g_2)
    \vartheta_L(g_0,g_2)^{-1},\\
    F_R(g_0,g_1,g_2)
    &:=
    \vartheta_R(g_1,g_2)
    \vartheta_R(g_0,g_1)
    \vartheta_R(g_0,g_2)^{-1}.
\end{aligned}
\]
Both admit canonical coherent trivializations at the identity of $\mu$.  The
higher simplices of $\Theta_L$ and $\Theta_R$ provide the corresponding higher
Bianchi coherences.
\end{proposition}

\begin{proof}
For the left expression,
\[
    (g_0^{-1}g_1)(g_1^{-1}g_2)
    =
    g_0^{-1}g_2,
\]
and for the right expression,
\[
    (g_2g_1^{-1})(g_1g_0^{-1})
    =
    g_2g_0^{-1}.
\]
Thus both curvatures are canonically the identity.  On a $3$-simplex the two
ways of composing three successive increments are related by associativity,
and the complete hierarchy is exactly the hierarchy of simplicial identities
of the bar cocycle.
\end{proof}

\begin{remark}[Why there is no general fixed-group Maurer--Cartan map]
\label{rem:no-general-fixed-group-mc}
The maps $\Theta_L$ and $\Theta_R$ exist because translation trivializes every
formal source fibre of the group object $G$ by the single identity fibre
$\mu$.  For a general formal Lie groupoid the anchor fibres form only the varying
formal-isotropy bitorsors of Theorem~\ref{thm:formal-bitorsor}, and no canonical map to the bar construction of one fixed group
exists.  The general replacement is the groupoid nerve itself together with
its Cartan and adjoint transport.  Additive and operator-valued
Maurer--Cartan objects arise after differentiation and Spencer normalization.
\end{remark}

\section{Represented and adic interfaces}
\label{sec:flg-represented-interface}


The intrinsic formal Lie groupoid theory is complete without choosing an
all-orders represented formal completion, principal-parts algebra, adic model,
or coordinate system.  No such general represented completion of a groupoid
unit has been constructed in the manuscript, and the intrinsic theory below
does not use one.
The represented interface actually constructed in this chapter is the
first-order structured neighbourhood and its comparison with the diagonal,
recorded in the handoff section below.

\begin{remark}[What is postponed]
Finite principal-parts towers, completed function algebras, distribution
objects, refined Koszul duals, partition-Lie objects, PBW comparisons, adic
formal-completion models, and completion exchange morphisms are postponed to
the differentiation and recognition chapters.  Any future represented or
formal-moduli model must come with its own construction and a proved comparison
with the intrinsic formal groupoid.  The intrinsic package through
Theorem~\ref{thm:formal-lie-groupoid-package} uses only the nonlinear formal
groupoid and its intrinsic descent calculus; the final handoff uses the
already constructed first structured neighbourhood machinery to record the
represented first-order anchor.
\end{remark}


\section{The formal Lie groupoid package}
\label{sec:flg-synthesis}

\begin{theorem}[Brave-New Formal Lie Groupoid Package]
\label{thm:formal-lie-groupoid-package}
Let $X\to S$ be an object of the native big affine spectral Nisnevich
$\infty$-topos.  The constructions of this chapter canonically provide:
\begin{enumerate}
    \item the centred relative de Rham completion
    \[
        C_X(Y)
        =
        Y\times_{Y_{\mathrm{dR}/S}}X_{\mathrm{dR}/S}
    \]
    on $\mathcal E_{X/}$, preserving pullbacks and fixing
    $(X,\operatorname{id}_X)$, with fixed points forming a coreflective
    subcategory and associated idempotent comonad;
    \item the induced idempotent unitwise completion
    \[
        \widehat{(-)}^{\,\mathrm{dR}}:
        \operatorname{Gpd}_X(\mathcal X_{/S})
        \longrightarrow
        \operatorname{Gpd}_X(\mathcal X_{/S}),
    \]
    with degree-one and localized-unit recognition of its fixed points;
    \item the canonical effective quotient map
    \[
        \epsilon_{X/S}:X\twoheadrightarrow Q_{X/S},
        \qquad
        L_S(\epsilon_{X/S}):
        L_S(X)\xrightarrow{\sim}L_S(Q_{X/S}),
    \]
    together with the maximal relative infinitesimal groupoid
    \[
        R_{X/S,\bullet}
        \simeq
        \check C_\bullet
        \left(
            X\twoheadrightarrow Q_{X/S}
        \right)
        \simeq
        \widehat{\operatorname{Pair}(X/S)}^{\,\mathrm{dR}}_\bullet,
    \]
    terminal among formal fixed-object groupoids;
    \item formal Lie groupoids defined without anchor data by
    \[
        \operatorname{FormLieGpd}(X/S)
        =
        \operatorname{Gpd}^{\mathrm{for}}_X(\mathcal X_{/S}),
    \]
    together with their contractibly unique constructed canonical anchors
    \[
        \mathfrak G_\bullet\longrightarrow R_{X/S,\bullet};
    \]
    \item the preformal anchored category
    \[
        \operatorname{AnchGpd}(X/S)
        =
        \left(
            \operatorname{Gpd}_X(\mathcal X_{/S})
        \right)_{/R_{X/S,\bullet}}
    \]
    and its quotient classification
    \[
        \operatorname{AnchGpd}(X/S)
        \simeq
        \operatorname{Fact}_{\mathrm{eff}}(\epsilon_{X/S}),
    \]
    where every factorization is automatically a chain
    \[
        X\twoheadrightarrow Q\twoheadrightarrow Q_{X/S}
    \]
    of effective epimorphisms;
    \item for every effective quotient $q:X\twoheadrightarrow Q$, the
    quotient formalization
    \[
        P_q
        =
        Q\times_{Q_{\mathrm{dR}/S}}X_{\mathrm{dR}/S},
        \qquad
        X\twoheadrightarrow\widehat Q_q\hookrightarrow P_q,
    \]
    its right-adjoint coreflection, and the equivalent recognition criteria
    \[
    \begin{aligned}
        \check C_\bullet(q)\text{ formal}
        &\Longleftrightarrow
        q_{\mathrm{dR}/S}\text{ is a monomorphism}\\
        &\Longleftrightarrow
        q_{\mathrm{dR}/S}\text{ is an equivalence}\\
        &\Longleftrightarrow
        \widehat Q_q\xrightarrow{\sim}Q;
    \end{aligned}
    \]
    \item the formal quotient classification
    \[
        \operatorname{FormLieGpd}(X/S)
        \simeq
        \operatorname{Epi}^{\mathrm{for}}(X)
        \simeq
        \operatorname{Fact}^{\mathrm{for}}_{\mathrm{eff}}
        (\epsilon_{X/S}),
    \]
    with a canonically recovered factorization
    \[
        X\twoheadrightarrow Q_{\mathfrak G}
        \twoheadrightarrow Q_{X/S}
    \]
    in which both arrows are relative de Rham equivalences;
    \item for every generalized point $a:T\to X$, the intrinsically complete
    formal orbit disk
    \[
        \mathbb D^{\mathfrak G}_a
        =
        T\times_{Q_{\mathfrak G}}X
        \longrightarrow
        \mathbb D^\infty_{X/S,a};
    \]
    \item for every fixed-object groupoid, the raw isotropy group and its
    finite-limit inertia formula
    \[
        I^{\mathrm{raw}}_{\mathfrak G}
        \simeq
        q_{\mathfrak G}^*
        \mathcal I_{Q_{\mathfrak G}/S};
    \]
    \item for every formal groupoid, the formal isotropy group
    \[
        I^{\mathrm{for}}_{\mathfrak G}
        =
        C_X(I^{\mathrm{raw}}_{\mathfrak G})
        \simeq
        \mathfrak G_1\times_{R_{X/S}}X,
    \]
    together with the constructed formal adjoint group
    \[
        \operatorname{Ad}^{\mathrm{for}}(\mathfrak G)
        :=
        \mathcal K_{\mathfrak G},
        \qquad
        I^{\mathrm{for}}_{\mathfrak G}
        \simeq
        q_{\mathfrak G}^*\mathcal K_{\mathfrak G},
    \]
    characterized as the contractibly unique effective descent of formal
    isotropy;
    \item the raw source-target bitorsor over
    $M^{\mathrm{raw}}_{\mathfrak G}$ and the formal anchor bitorsor over
    $M^{\mathrm{for}}_{\mathfrak G}$, with
    \[
        t^*I^{\mathrm{for}}_{\mathfrak G}
        \simeq
        \mathfrak G_1\times_{R_{X/S}}\mathfrak G_1
        \simeq
        s^*I^{\mathrm{for}}_{\mathfrak G};
    \]
    \item the nonlinear formal-isotropy commutator, coherent formal adjoint
    transport, and the theorem identifying its \v{C}ech transport with
    conjugation on the constructed inertia kernel;
    \item for every morphism $q:X\to B$, its effective image
    \[
        X\twoheadrightarrow I_q\lhook\joinrel\longrightarrow B
    \]
    and the effective-image relation disk--jet calculus
    \[
        T_q=q^*\Sigma_q\simeq e_q^*\Sigma_{e_q},
        \qquad
        J_q=q^*\Pi_q\simeq e_q^*\Pi_{e_q},
    \]
    with
    \[
        T_q\simeq\Sigma_s t^*\simeq\Sigma_t s^*,
        \qquad
        J_q\simeq\Pi_s t^*,
        \qquad
        T_q\dashv J_q,
    \]
    and
    \[
        \mathcal X_{/I_q}
        \simeq
        \operatorname{Alg}_{T_q}(\mathcal X_{/X})
        \simeq
        \operatorname{Coalg}_{J_q}(\mathcal X_{/X});
    \]
    \item functoriality in the groupoid and geometric object, automatic
    compatibility of canonical anchors, arbitrary spectral base change,
    fixed-object quotient invariance, and the explicit split-atlas
    calculation
    \[
        \widehat{\check C_\bullet(X\amalg X\twoheadrightarrow X)}^{\,\mathrm{dR}}
        \simeq
        0_{X\amalg X,\bullet},
    \]
    showing both failure of arbitrary-atlas invariance and failure of
    unitwise formalization to restore it;
    \item the canonical formalization
    \[
        \mathcal H_\bullet
        \longmapsto
        \widehat{\mathcal H}^{\,\mathrm{dR}}_\bullet
        \in
        \operatorname{FormLieGpd}(X/S)
    \]
    for arbitrary internal groupoids on $X$;
    \item for every group object $G/S$, the constructed formal identity group
    \[
        \mu^{\mathrm{mod}}_{G/S}
        =
        G\times_{G_{\mathrm{dR}/S}}S,
    \]
    its functorial group structure, and for every $G$-action the equivalence
    \[
        \widehat{(G\ltimes X)}^{\,\mathrm{dR}}
        \simeq
        \mu^{\mathrm{mod}}_{G/S}\ltimes X;
    \]
    \item the universal classifying object
    \[
        b_G:S\twoheadrightarrow BG,
        \qquad
        \check C_\bullet(b_G)\simeq B_\bullet G,
    \]
    principal right $G$-objects constructed as
    $P_\chi=X\times_{BG}S$, recognition of these objects by effective
    epimorphism plus invertible shear, the classification equivalence
    \[
        \operatorname{Map}_{\mathcal E}(X,BG)
        \simeq
        \operatorname{Prin}_G(X),
    \]
    and the gauge groupoid
    \[
        \operatorname{Gauge}(P_\chi)
        =
        \check C_\bullet(\chi);
    \]
    \item the one-object specialization
    \[
        \widehat{B_\bullet G}^{\,\mathrm{dR}}
        \simeq
        B_\bullet\mu^{\mathrm{mod}}_{G/S},
    \]
    the translation equivalence
    \[
        R_{G/S,\bullet}
        \simeq
        (G/\!/\mu^{\mathrm{mod}}_{G/S})_\bullet,
    \]
    and the integrated group-valued Maurer--Cartan cocycles.
\end{enumerate}
Every coherence in this package is inherited from the relative de Rham
left-exact localization, finite limits, the centred completion comonad,
groupoid effectivity, universal colimits, dependent adjunction,
Beck--Chevalley, functorial effective-image factorization, or effective
descent.  The terms ``formal'', ``principal'', ``gauge'', and ``Cartan''
refer respectively to fixed points of the constructed coreflection, objects
classified by the proved equivalence
$\operatorname{Map}_{\mathcal E}(X,BG)\simeq\operatorname{Prin}_G(X)$,
\v{C}ech groupoids of constructed classifying morphisms, and objects over the
constructed effective formal quotient.  They introduce no separate
admissibility, presentation, torsor, anchor, or descent axioms.  No additive
bracket or Lie-algebraic presentation has been introduced.
\end{theorem}

\begin{proof}
Items (1)--(2) are
Propositions~\ref{prop:centred-completion-left-exact} and
\ref{prop:centred-completion-recognition},
Theorems~\ref{thm:centred-completion-idempotent},
\ref{thm:centred-formalization-coreflection},
\ref{thm:unitwise-completion-idempotent}, and
\ref{thm:formalization-coreflection}.
Items (3)--(4) are
Lemma~\ref{lem:effective-quotient-dr-equivalence},
Theorems~\ref{thm:pair-groupoid-completion} and
\ref{thm:maximal-formal-terminal},
Corollary~\ref{cor:canonical-formal-anchor}, and
Definition~\ref{def:formal-lie-groupoid}.
Item (5) is
Definition~\ref{def:infinitesimal-quotient-factorizations},
Lemma~\ref{lem:second-factor-effective}, and
Theorem~\ref{thm:anchored-gpd-quotient-classification}.
Items (6)--(7) are
Theorems~\ref{thm:quotient-formula-unitwise-completion},
\ref{thm:formal-quotient-monomorphism-recognition},
\ref{thm:quotient-formalization-coreflection}, and
\ref{thm:formal-gpd-quotient-classification}, together with
Proposition~\ref{prop:formal-quotient-canonical-anchor}.
Item (8) is
Definition~\ref{def:formal-orbit-disk} and
Propositions~\ref{prop:formal-disk-unitwise-complete} and
\ref{prop:formal-disk-ambient-comparison}.
Item (9) is
Definition~\ref{def:relative-inertia-object} and
Propositions~\ref{prop:relative-inertia-group} and
\ref{prop:isotropy-inertia}.
Item (10) is
Definition~\ref{def:formal-inertia-kernel} and
Theorems~\ref{thm:isotropy-group} and
\ref{thm:formal-isotropy-descent}.
Item (11) is
Theorems~\ref{thm:raw-bitorsor} and
\ref{thm:formal-bitorsor}.
Item (12) is
Definition~\ref{def:isotropy-commutator},
Theorem~\ref{thm:coherent-adjoint-transport},
Proposition~\ref{prop:inertia-descent-conjugation}, and
Corollary~\ref{cor:adjoint-inertia-object}.
Item (13) is
Theorem~\ref{thm:effective-relation-calculus}.
Item (14) is
Propositions~\ref{prop:automatic-anchor-functoriality},
\ref{prop:formal-gpd-functoriality-fixed-X},
\ref{prop:presentation-invariance}, and
\ref{prop:formality-not-morita-invariant}, together with
Theorem~\ref{thm:formal-gpd-base-change}.
Item (15) is
Theorem~\ref{thm:unitwise-formal-groupoid-completion}.
Item (16) is
Proposition~\ref{prop:formal-identity-group} and
Theorem~\ref{thm:action-groupoid-completion}.
Item (17) is
Definitions~\ref{def:classifying-object-group},
\ref{def:principal-group-object},
\ref{def:principal-object-space},
\ref{def:gauge-groupoid}, and
\ref{def:formal-gauge-groupoid}, together with
Theorems~\ref{prop:principal-classifying-map} and
\ref{thm:principal-objects-classified-by-bg}, and
Propositions~\ref{prop:gauge-arrow-quotient} and
\ref{prop:gauge-formal-lie-groupoid}.
Item (18) is
Theorems~\ref{thm:one-object-completion},
\ref{thm:group-translation},
\ref{thm:higher-translation-coordinates}, and
\ref{thm:integrated-mc-cocycle}.
\end{proof}


\section{Handoff to brave-new differentiation}
\label{sec:flg-handoff}

The nonlinear input for a later differentiation theorem is now completely
intrinsic.  A formal Lie groupoid has a canonical anchor
\[
    \phi_{\mathfrak G}:
    \mathfrak G_\bullet
    \longrightarrow
    R_{X/S,\bullet},
\]
a formal quotient
\[
\begin{tikzcd}
X \arrow[r,two heads,"q_{\mathfrak G}"]
  \arrow[dr,two heads,"\epsilon_{X/S}"'] &
Q_{\mathfrak G}
  \arrow[d,two heads,"r_{\mathfrak G}"]\\
&Q_{X/S},
\end{tikzcd}
\]
and the equivalent quotient-side conditions
\[
    (q_{\mathfrak G})_{\mathrm{dR}/S}
    \text{ monic},
    \qquad
    (q_{\mathfrak G})_{\mathrm{dR}/S}
    \text{ an equivalence},
    \qquad
    \widehat Q_{q_{\mathfrak G}}
    \xrightarrow{\sim}
    Q_{\mathfrak G}.
\]
The canonically recovered second quotient map is also a relative de Rham
equivalence:
\[
    L_S(r_{\mathfrak G}):
    L_S(Q_{\mathfrak G})
    \xrightarrow{\sim}
    L_S(Q_{X/S}).
\]
For every generalized point there is a morphism of intrinsically complete
formal orbit disks
\[
    \mathbb D^{\mathfrak G}_a
    \longrightarrow
    \mathbb D^\infty_{X/S,a}.
\]
The infinitesimal stabilizer of the anchor is the formal isotropy group
\[
    I^{\mathrm{for}}_{\mathfrak G}
    \simeq
    \mathfrak G_1\times_{R_{X/S}}X
    \simeq
    q_{\mathfrak G}^*\mathcal K_{\mathfrak G},
\]
with its commutator, formal bitorsor action, and coherent adjoint transport.

The present manuscript has not constructed a tangent object
$\mathbb T_{X/S}$ for an arbitrary object of the big spectral Nisnevich
$\infty$-topos.  The first-order comparison of the preceding de Rham chapter
is established for represented morphisms of spectral schemes.  Accordingly,
the first linearization of the formal-groupoid comparison must be made in the
represented range using the already constructed structured square-zero
neighbourhood machinery.

When $X\to S$ is represented by a spectral morphism, the preceding chapters
provide the cotangent complex and the first structured neighbourhood of the
diagonal.  To linearize a formal groupoid itself, the arrow geometry must also
lie in the represented range where its relative cotangent complex and
structured square-zero extensions have been constructed.  We first make this
comparison without using smoothness.

\begin{proposition}[Representability from degrees zero and one]
\label{prop:represented-groupoid-segal}
Let
\[
    \mathfrak G_\bullet\in
    \operatorname{Gpd}_X(\mathcal X_{/S})
\]
be a fixed-object groupoid.  If $X=\mathfrak G_0$ and $\mathfrak G_1$ are
represented by spectral schemes, then every $\mathfrak G_n$ is represented by
a spectral scheme.  Under these representations all face and degeneracy maps
are morphisms of spectral schemes, so $\mathfrak G_\bullet$ is a simplicial
spectral scheme.  No separate higher-degree representability condition is
required.
\end{proposition}

\begin{proof}
For $n\geq2$, the Segal equivalence gives
\[
    \mathfrak G_n
    \simeq
    \underbrace{
    \mathfrak G_1\times_X\cdots\times_X\mathfrak G_1
    }_{n\text{ factors}}.
\]
Spectral schemes admit finite fibre products.  Moreover the native
functor-of-points realization preserves these fibre products: for spectral
schemes $A\to C\leftarrow B$ and every affine spectral test object $T$,
\[
\begin{aligned}
    h_{A\times_C B}(T)
    &=
    \operatorname{Map}(T,A\times_C B)\\
    &\simeq
    \operatorname{Map}(T,A)
    \times_{\operatorname{Map}(T,C)}
    \operatorname{Map}(T,B)\\
    &=
    \left(h_A\times_{h_C}h_B\right)(T).
\end{aligned}
\]
Thus the ambient Segal fibre product displayed above is the native realization
of the corresponding finite fibre product of spectral schemes.  Hence
$\mathfrak G_n$ is represented for every $n\geq2$.  Degree zero and degree one
are represented by hypothesis.

The native Yoneda realization of spectral schemes in $\mathcal X$ is fully
faithful.  Hence every simplicial operator between the represented native
objects is induced by a unique morphism of spectral schemes, and the
simplicial identities are preserved by full faithfulness.
\end{proof}

\subsection{The first structured neighbourhood of the groupoid unit}

Work now in the represented range: let $S$, $X$, and $\mathfrak G_1$ be
spectral schemes representing the corresponding native objects.  Write
\[
    \mathfrak G_1
    \mathop{\rightrightarrows}^{s}_{t}
    X,
    \qquad
    u:X\longrightarrow\mathfrak G_1
\]
for the source, target, and unit.

\begin{definition}[First structured neighbourhood of the groupoid unit]
\label{def:first-unit-neighbourhood}
Define the first-order source coefficient
\[
    \mathfrak a_{\mathfrak G}^{\vee}
    :=
    u^*L_{\mathfrak G_1/X,s}.
\]
Since
\[
    s u=\operatorname{id}_X,
\]
the transitivity triangle for
\[
    X\xrightarrow{u}\mathfrak G_1\xrightarrow{s}X
\]
has the form
\[
    \mathfrak a_{\mathfrak G}^{\vee}
    \longrightarrow
    L_{X/X}\simeq0
    \longrightarrow
    L_{X/\mathfrak G_1}.
\]
Using the same rotation convention for transitivity triangles that the
preceding de Rham chapter fixed for the diagonal, define the resulting
connecting equivalence
\[
    \partial_u:
    L_{X/\mathfrak G_1}
    \xrightarrow{\sim}
    \mathfrak a_{\mathfrak G}^{\vee}[1].
\]
This fixes the sign compatibly with the previously constructed
$\partial_\Delta$.
The global structured square-zero classification applied to
$\partial_u$ produces a canonical structured square-zero thickening
\[
    X
    \xrightarrow{i_u}
    D^{(1)}_{\mathfrak G}
    \xrightarrow{\nu_u}
    \mathfrak G_1
\]
with coefficient $\mathfrak a_{\mathfrak G}^{\vee}$.  We call
$D^{(1)}_{\mathfrak G}$ the \emph{first structured infinitesimal
neighbourhood of the groupoid unit}.
\end{definition}

\begin{remark}
No smoothness or dualizability is used in
Definition~\ref{def:first-unit-neighbourhood}.  The global cotangent complex is
connective for represented spectral morphisms, so the coefficient
$\mathfrak a_{\mathfrak G}^{\vee}$ lies in the connective range required by
the structured square-zero construction.
\end{remark}

\subsection{First-order comparison with the diagonal}

Put
\[
    Y:=X\times_SX,
    \qquad
    f:=(s,t):\mathfrak G_1\longrightarrow Y,
    \qquad
    \Delta:=fu:X\longrightarrow Y.
\]
Regard $Y$ over $X$ through the first projection.  The preceding de Rham
chapter constructed the first structured neighbourhood of the diagonal
\[
    X
    \xrightarrow{i_1}
    D^{(1)}_{X/S}
    \xrightarrow{\nu_1}
    X\times_SX
\]
with coefficient $L_{X/S}$, together with the canonical morphism
\[
    \ell_1:
    D^{(1)}_{X/S}
    \longrightarrow
    R_{X/S}.
\]

\begin{lemma}[Naturality of structured square-zero neighbourhoods]
\label{lem:structured-square-zero-naturality}
Let
\[
    X\xrightarrow{u}Y\xrightarrow{f}Z
\]
be morphisms of spectral schemes.  Let $M,N\in\operatorname{QCoh}(X)_{\geq0}$,
let
\[
    \eta_Y:L_{X/Y}\longrightarrow M[1],
    \qquad
    \eta_Z:L_{X/Z}\longrightarrow N[1]
\]
be structured square-zero classifying morphisms, and let
\[
    \alpha:N\longrightarrow M
\]
be a coefficient morphism.  Suppose the natural cotangent map
\[
    \lambda_f:L_{X/Z}\longrightarrow L_{X/Y}
\]
and the two classifying morphisms fit into a specified homotopy-commutative
square
\[
\begin{tikzcd}
L_{X/Z} \arrow[r,"{\eta_Z}"] \arrow[d,"\lambda_f"'] &
N[1] \arrow[d,"{\alpha[1]}"]\\
L_{X/Y} \arrow[r,"\eta_Y"'] & M[1].
\end{tikzcd}
\]
Let
\[
    X\xrightarrow{i_Y}D_{\eta_Y}\xrightarrow{\nu_Y}Y,
    \qquad
    X\xrightarrow{i_Z}D_{\eta_Z}\xrightarrow{\nu_Z}Z
\]
be the structured square-zero thickenings constructed from $\eta_Y$ and
$\eta_Z$.  Then the displayed compatibility induces a canonical morphism
\[
    D_{\eta_Y}\longrightarrow D_{\eta_Z}
\]
which restricts to $\operatorname{id}_X$, lies over $f:Y\to Z$, and has
coefficient morphism $\alpha$.  The construction is functorial in compatible
morphisms of such diagrams and commutes with arbitrary represented base
change.
\end{lemma}

\begin{proof}
Write
\[
    \mathcal A_Z:=(fu)^{-1}\mathcal O_Z,
    \qquad
    \mathcal A_Y:=u^{-1}\mathcal O_Y.
\]
The structural morphism
$\mathcal A_Z\to\mathcal A_Y\to\mathcal O_X$ lets us regard both structured
extensions as extensions of $\mathcal O_X$ in commutative
$\mathcal A_Z$-algebras on the Zariski $\infty$-topos of $X$.  The
restriction-of-scalars functor from commutative $\mathcal A_Y$-algebras to
commutative $\mathcal A_Z$-algebras creates limits.  Consequently the finite
pullback defining the $Y$-relative structured square-zero extension remains
the same pullback after restriction of scalars to $\mathcal A_Z$.

Under the relative cotangent--derivation correspondence, the classifying
morphisms determine derivations
\[
    d_{\eta_Z}:
    \mathcal O_X\longrightarrow
    \mathcal O_X\oplus N[1],
    \qquad
    d_{\eta_Y}:
    \mathcal O_X\longrightarrow
    \mathcal O_X\oplus M[1].
\]
The first is $\mathcal A_Z$-linear.  The second is
$\mathcal A_Y$-linear and hence also $\mathcal A_Z$-linear.  The coefficient
morphism $\alpha:N\to M$ induces a morphism of split square-zero
$\mathcal A_Z$-algebras
\[
    \operatorname{id}\oplus\alpha[1]:
    \mathcal O_X\oplus N[1]
    \longrightarrow
    \mathcal O_X\oplus M[1].
\]
Naturality of the cotangent--derivation correspondence identifies the
specified homotopy-commutative square of classifying morphisms with a
specified homotopy
\[
    (\operatorname{id}\oplus\alpha[1])d_{\eta_Z}
    \simeq
    d_{\eta_Y}
\]
in the $\infty$-category of commutative $\mathcal A_Z$-algebras under
$\mathcal O_X$.  The zero derivations are compatible with
$\operatorname{id}\oplus\alpha[1]$ as well.

The global structured square-zero construction forms the two structure
sheaves as the specified pullbacks
\[
    \mathcal O_{D_{\eta_Z}}
    \simeq
    \mathcal O_X
    \times_{\mathcal O_X\oplus N[1]}
    \mathcal O_X,
    \qquad
    \mathcal O_{D_{\eta_Y}}
    \simeq
    \mathcal O_X
    \times_{\mathcal O_X\oplus M[1]}
    \mathcal O_X,
\]
using respectively the structured derivation and the zero derivation.  The
preceding compatible morphisms give a morphism between these two cospan
diagrams.  Functoriality of limits therefore produces a canonical morphism
of commutative $\mathcal A_Z$-algebras
\[
    \mathcal O_{D_{\eta_Z}}
    \longrightarrow
    \mathcal O_{D_{\eta_Y}}.
\]
Contravariance of spectral schemes yields
\[
    D_{\eta_Y}\longrightarrow D_{\eta_Z}.
\]
Because the morphism of cospans is the identity on both copies of
$\mathcal O_X$, the resulting morphism restricts to
$\operatorname{id}_X$.  Because the entire construction was performed in
commutative $\mathcal A_Z$-algebras and the $Y$-structure is retained on the
target pullback, it lies over $f:Y\to Z$.  The induced map on the
square-zero fibres is the original coefficient morphism
$\alpha:N\to M$.

For fixed structured derivations, the global structured square-zero
classification identifies the space of compatible pullback models and
specified pullback identifications with a contractible space.  Hence the
morphism just constructed introduces no additional extension datum or
choice.  Compatible morphisms of the input diagrams induce compatible
morphisms of the defining cospans, so identities, composition, and all higher
functorial coherences follow from functoriality of limits.

Finally, represented base change carries the cotangent map
$\lambda_f$, the classifying morphisms, and the coefficient morphism to their
base changes by cotangent base change.  Inverse image preserves split
square-zero algebras and the finite pullbacks defining the structured
extensions.  Applying the same cospan construction after base change
therefore identifies the pulled-back morphism with the morphism constructed
from the base-changed data.
\end{proof}


\begin{proposition}[First-order source--target comparison]
\label{prop:first-order-source-target-comparison}
There is a canonical cotangent morphism
\[
    a_{\mathfrak G}^{\vee}:
    L_{X/S}
    \longrightarrow
    \mathfrak a_{\mathfrak G}^{\vee}
\]
and a canonical morphism of structured first neighbourhoods
\[
    d^{(1)}(s,t):
    D^{(1)}_{\mathfrak G}
    \longrightarrow
    D^{(1)}_{X/S}
\]
fitting into a commutative diagram
\[
\begin{tikzcd}
X \arrow[r,"i_u"] \arrow[d,equal] &
D^{(1)}_{\mathfrak G}
    \arrow[r,"\nu_u"]
    \arrow[d,"{d^{(1)}(s,t)}"']
&
\mathfrak G_1
    \arrow[d,"{(s,t)}"]
\\
X \arrow[r,"i_1"'] &
D^{(1)}_{X/S}
    \arrow[r,"\nu_1"']
&
X\times_SX .
\end{tikzcd}
\]
The coefficient morphism of this structured square-zero comparison is
$a_{\mathfrak G}^{\vee}$.  These constructions are compatible with arbitrary
represented base change.
\end{proposition}

\begin{proof}
Cotangent base change for the first projection gives
\[
    L_{(X\times_SX)/X,\operatorname{pr}_1}
    \simeq
    \operatorname{pr}_2^*L_{X/S}.
\]
The morphism
\[
    f=(s,t):
    (\mathfrak G_1,s)
    \longrightarrow
    (X\times_SX,\operatorname{pr}_1)
\]
therefore induces
\[
    t^*L_{X/S}
    \simeq
    f^*L_{(X\times_SX)/X,\operatorname{pr}_1}
    \longrightarrow
    L_{\mathfrak G_1/X,s}.
\]
Pulling back along the unit and using
$tu=\operatorname{id}_X$ gives
\[
    a_{\mathfrak G}^{\vee}:
    L_{X/S}
    \longrightarrow
    u^*L_{\mathfrak G_1/X,s}
    =
    \mathfrak a_{\mathfrak G}^{\vee}.
\]

Naturality of the transitivity triangles for the two sections
\[
    \Delta:X\to X\times_SX,
    \qquad
    u:X\to\mathfrak G_1
\]
under the morphism $f$ gives a commutative square
\[
\begin{tikzcd}
L_{X/(X\times_SX)}
    \arrow[r,"{\partial_\Delta}"']
    \arrow[d]
&
L_{X/S}[1]
    \arrow[d,"{a_{\mathfrak G}^{\vee}[1]}"]
\\
L_{X/\mathfrak G_1}
    \arrow[r,"{\partial_u}"']
&
\mathfrak a_{\mathfrak G}^{\vee}[1],
\end{tikzcd}
\]
where $\partial_\Delta$ is the fixed connecting equivalence used to construct
$D^{(1)}_{X/S}$.  Apply
Lemma~\ref{lem:structured-square-zero-naturality} to the composable maps
\[
    X\xrightarrow{u}\mathfrak G_1
    \xrightarrow{(s,t)}X\times_SX,
\]
the coefficient map
\[
    a_{\mathfrak G}^{\vee}:
    L_{X/S}\longrightarrow\mathfrak a_{\mathfrak G}^{\vee},
\]
and the displayed homotopy of classifying morphisms.  It produces the
canonical morphism
\[
    d^{(1)}(s,t):
    D^{(1)}_{\mathfrak G}
    \longrightarrow
    D^{(1)}_{X/S}.
\]
The lemma also shows that it preserves the centre sections, lies over
$(s,t):\mathfrak G_1\to X\times_SX$, has coefficient morphism
$a_{\mathfrak G}^{\vee}$, and is compatible with arbitrary represented base
change.
\end{proof}

\begin{remark}[Equivalent formula for the cotangent anchor]
\label{rem:cotangent-anchor-equivalent-formula}
The morphism of
Proposition~\ref{prop:first-order-source-target-comparison} is equivalently the
composite obtained from the target map and transitivity:
\[
    L_{X/S}
    \simeq
    u^*t^*L_{X/S}
    \longrightarrow
    u^*L_{\mathfrak G_1/S}
    \longrightarrow
    u^*L_{\mathfrak G_1/X,s}.
\]
Thus the first-neighbourhood construction recovers the usual cotangent
source--target morphism, but now with its nonlinear origin retained.
\end{remark}

\subsection{Linearization of the canonical formal anchor}

\begin{theorem}[First-order linearization of the canonical anchor]
\label{thm:first-order-linearization-canonical-anchor}
Let
\[
    \mathfrak G_\bullet
    \in
    \operatorname{FormLieGpd}(X/S)
\]
be represented in degrees zero and one as above.  Put
\[
    D:=D^{(1)}_{\mathfrak G},
    \qquad
    Y:=X\times_SX,
    \qquad
    f_D:=(s,t)\nu_u:D\longrightarrow Y,
\]
and let
\[
    \kappa_Y:R_{X/S}=C_X(Y,\Delta)\longrightarrow Y
\]
be the completion counit.  The lift space
\[
    \operatorname{Lift}(f_D)
    :=
    \operatorname{Map}_{\mathcal E_{X/}}(D,R_{X/S})
    \times_{
        \operatorname{Map}_{\mathcal E_{X/}}(D,Y)
    }
    \{f_D\}
\]
is contractible.

The two composites in
\[
\begin{tikzcd}
D^{(1)}_{\mathfrak G}
    \arrow[r,"{d^{(1)}(s,t)}"]
    \arrow[d,"\nu_u"']
&
D^{(1)}_{X/S}
    \arrow[d,"\ell_1"]
\\
\mathfrak G_1
    \arrow[r,"{\phi_{\mathfrak G,1}}"']
&
R_{X/S}
\end{tikzcd}
\]
define points of this lift space.  Consequently they are canonically
homotopic over the specified identification of their composites with
$f_D$, and the space of such compatible homotopies is contractible.
The coefficient morphism
\[
    a_{\mathfrak G}^{\vee}:
    L_{X/S}
    \longrightarrow
    \mathfrak a_{\mathfrak G}^{\vee}
\]
is therefore the first-order coefficient of the canonical nonlinear
infinitesimal anchor
\[
    \phi_{\mathfrak G,1}:
    \mathfrak G_1\longrightarrow R_{X/S}.
\]
\end{theorem}

\begin{proof}
The centre
\[
    i_u:
    X\longrightarrow D^{(1)}_{\mathfrak G}
\]
is a global structured square-zero thickening.  The preceding de Rham chapter
proves that every such thickening is inverted by relative de Rham
localization.  Hence
\[
    L_S(i_u):
    L_S(X)
    \xrightarrow{\sim}
    L_S(D)
\]
is an equivalence, and
Proposition~\ref{prop:centred-completion-recognition} shows that $D$ is
complete along its centre.

The centred formalization coreflection therefore gives an equivalence
\[
    \kappa_Y\circ(-):
    \operatorname{Map}_{\mathcal E_{X/}}(D,R_{X/S})
    \xrightarrow{\sim}
    \operatorname{Map}_{\mathcal E_{X/}}(D,Y).
\]
Its homotopy fibre over the specified point
$f_D=(s,t)\nu_u$ is precisely $\operatorname{Lift}(f_D)$, so
\[
    \operatorname{Lift}(f_D)\simeq *.
\]

By Proposition~\ref{prop:canonical-anchor-completed-source-target}, the
degree-one canonical anchor is the completed source--target map
\[
    \phi_{\mathfrak G,1}
    \simeq
    C_X(s,t):
    \mathfrak G_1
    \longrightarrow
    R_{X/S}.
\]
Therefore
\[
    \phi_{\mathfrak G,1}\nu_u:
    D\longrightarrow R_{X/S}
\]
is a lift of $f_D$.

On the other hand,
Proposition~\ref{prop:first-order-source-target-comparison} gives
\[
    \nu_1d^{(1)}(s,t)
    \simeq
    (s,t)\nu_u=f_D.
\]
The defining property of
\[
    \ell_1:
    D^{(1)}_{X/S}\longrightarrow R_{X/S}
\]
identifies $\kappa_Y\ell_1$ with $\nu_1$.  Thus
\[
    \ell_1d^{(1)}(s,t):
    D\longrightarrow R_{X/S}
\]
is a second point of $\operatorname{Lift}(f_D)$.

Since this lift space is contractible, the space of paths between these two
points inside $\operatorname{Lift}(f_D)$ is contractible.  Equivalently, the
two composites are homotopic through a contractible space of homotopies whose
projection to $Y$ is the specified source--target homotopy.  The comparison
$d^{(1)}(s,t)$ was constructed from the coefficient morphism
$a_{\mathfrak G}^{\vee}$, so that morphism is precisely the first-order
coefficient of the canonical nonlinear anchor.
\end{proof}

\subsection{The represented internal-dual tangent anchor}

\begin{theorem}[Represented internal-dual first-order anchor]
\label{thm:represented-internal-dual-anchor}
Let
\[
    \mathfrak G_1
    \mathop{\rightrightarrows}^{s}_{t}
    X
\]
be a brave-new formal Lie groupoid on $X/S$ with $S$, $X$, and
$\mathfrak G_1$ represented by spectral schemes.  For
$M\in\operatorname{QCoh}(X)$ write
\[
    M^\vee
    :=
    \underline{\operatorname{Hom}}_X(M,\mathcal O_X).
\]
Define
\[
    \mathfrak a_{\mathfrak G}
    :=
    \left(
        \mathfrak a_{\mathfrak G}^{\vee}
    \right)^\vee,
    \qquad
    \mathbb T^{\mathrm{rep}}_{X/S}
    :=
    L_{X/S}^{\vee}.
\]
Then the cotangent anchor
\[
    a_{\mathfrak G}^{\vee}:
    L_{X/S}\longrightarrow\mathfrak a_{\mathfrak G}^{\vee}
\]
canonically induces by internal duality a first-order tangent anchor
\[
\boxed{
    a_{\mathfrak G}:
    \mathfrak a_{\mathfrak G}
    \longrightarrow
    \mathbb T^{\mathrm{rep}}_{X/S}.
}
\]
No smoothness or dualizability is required for this construction.

For every represented base change $S'\to S$, with pulled-back groupoid
$\mathfrak G'$ on $X':=X\times_SS'$ and projection $c_X:X'\to X$, there are
canonical comparison morphisms
\[
    c_X^*\mathfrak a_{\mathfrak G}
    \longrightarrow
    \mathfrak a_{\mathfrak G'},
    \qquad
    c_X^*\mathbb T^{\mathrm{rep}}_{X/S}
    \longrightarrow
    \mathbb T^{\mathrm{rep}}_{X'/S'},
\]
and these comparisons carry $c_X^*a_{\mathfrak G}$ to
$a_{\mathfrak G'}$.
\end{theorem}

\begin{proof}
The category $\operatorname{QCoh}(X)$ is presentably symmetric monoidal and
its tensor product preserves colimits separately in each variable.  Hence,
for every $M$, the functor
\[
    -\otimes M:
    \operatorname{QCoh}(X)\longrightarrow\operatorname{QCoh}(X)
\]
admits a right adjoint by the presentable adjoint functor theorem.  Evaluating
that right adjoint at $\mathcal O_X$ gives the internal dual
$M^\vee=\underline{\operatorname{Hom}}_X(M,\mathcal O_X)$ for arbitrary
$M$.

Precomposition with
\[
    a_{\mathfrak G}^{\vee}:
    L_{X/S}\longrightarrow\mathfrak a_{\mathfrak G}^{\vee}
\]
therefore induces
\[
    \underline{\operatorname{Hom}}_X
    (\mathfrak a_{\mathfrak G}^{\vee},\mathcal O_X)
    \longrightarrow
    \underline{\operatorname{Hom}}_X
    (L_{X/S},\mathcal O_X),
\]
which is the displayed $a_{\mathfrak G}$.

Under represented base change, cotangent base change and the pulled-back unit
give canonical equivalences
\[
    c_X^*L_{X/S}\simeq L_{X'/S'},
    \qquad
    c_X^*\mathfrak a_{\mathfrak G}^{\vee}
    \simeq
    \mathfrak a_{\mathfrak G'}^{\vee}.
\]
Strong symmetric monoidality of $c_X^*$ carries the evaluation morphism
\[
    M^\vee\otimes M\longrightarrow\mathcal O_X
\]
to an evaluation pairing
\[
    c_X^*M^\vee\otimes c_X^*M\longrightarrow\mathcal O_{X'}.
\]
Adjunction therefore constructs a canonical comparison
\[
    c_X^*M^\vee\longrightarrow(c_X^*M)^\vee.
\]
Apply this first to $M=\mathfrak a_{\mathfrak G}^{\vee}$ and then to
$M=L_{X/S}$.  Naturality of evaluation and of
$a_{\mathfrak G}^{\vee}$ gives the asserted compatibility square for tangent
anchors.
\end{proof}

\subsection{The smooth finite-locally-free upgrade}

\begin{theorem}[Smooth finite-locally-free first-order anchor]
\label{thm:represented-first-order-anchor}
Let
\[
    \mathfrak G_1
    \mathop{\rightrightarrows}^{s}_{t}
    X
\]
be a brave-new formal Lie groupoid on $X/S$ with $S$, $X$, and
$\mathfrak G_1$ represented by spectral schemes.  When $X\to S$ and the
source morphism
\[
    s:\mathfrak G_1\to X
\]
are smooth, the modules
\[
    L_{X/S},
    \qquad
    \mathfrak a_{\mathfrak G}^{\vee}
    =
    u^*L_{\mathfrak G_1/X,s}
\]
are finite locally free and hence dualizable.  The represented internal-dual
anchor of Theorem~\ref{thm:represented-internal-dual-anchor} therefore takes
the finite-locally-free form
\[
\boxed{
    a_{\mathfrak G}:
    \mathfrak a_{\mathfrak G}
    \longrightarrow
    \mathbb T_{X/S},
}
\]
where
\[
    \mathfrak a_{\mathfrak G}
    :=
    \left(
        \mathfrak a_{\mathfrak G}^{\vee}
    \right)^{\vee},
    \qquad
    \mathbb T_{X/S}
    :=
    L_{X/S}^{\vee}.
\]
For arbitrary represented base change, the comparison morphisms of
Theorem~\ref{thm:represented-internal-dual-anchor} are equivalences in this
finite-locally-free range.  Thus the displayed tangent anchor commutes with
represented base change by canonical equivalences.
\end{theorem}

\begin{proof}
Smoothness of $s$ makes
$L_{\mathfrak G_1/X,s}$ finite locally free by the standing brave-new
smoothness theorem.  Pullback along the unit preserves finite locally free
modules, so
\[
    \mathfrak a_{\mathfrak G}^{\vee}
    =
    u^*L_{\mathfrak G_1/X,s}
\]
is finite locally free on $X$.  Smoothness of $X\to S$ similarly makes
$L_{X/S}$ finite locally free.

The coefficient morphism
\[
    a_{\mathfrak G}^{\vee}:
    L_{X/S}
    \longrightarrow
    \mathfrak a_{\mathfrak G}^{\vee}
\]
was constructed in
Proposition~\ref{prop:first-order-source-target-comparison}, and
Theorem~\ref{thm:first-order-linearization-canonical-anchor} identifies it as
the first-order coefficient of the canonical nonlinear anchor.
Theorem~\ref{thm:represented-internal-dual-anchor} dualizes this map without
requiring smoothness.  In the present range finite locally free modules are
dualizable, so their internal duals are the dualizable duals and biduality
holds.

Finally, pullback preserves finite locally free modules and their evaluation
and coevaluation maps.  Hence for either of the two finite locally free
modules $M$ above, the canonical base-change comparison
\[
    c_X^*M^\vee\longrightarrow(c_X^*M)^\vee
\]
is an equivalence.  The base-change square for $a_{\mathfrak G}$ from
Theorem~\ref{thm:represented-internal-dual-anchor} is therefore a square of
equivalences on the two vertical comparisons.
\end{proof}

The represented constructions above produce the first-order source module and
anchor in the full represented range already supported by the manuscript's cotangent
theory and identify
that anchor as the first-order shadow of the canonical nonlinear
formal-groupoid anchor.  A bracket, isotropy Lie object, shifted partition-Lie
structure, and higher operations still require a separate differentiation and
recognition theorem.  They are to be derived from the nonlinear formal
groupoid, its formal isotropy kernel, formal disks, and Cartan descent rather
than supplied as independent low-level structure.

In the characteristic-zero formal-moduli setting over a commutative
dg-algebra, dg Lie algebroids provide the relevant algebraic recognition
objects \cite{Nuiten2019}.  In general characteristic, partition Lie
algebras and partition Lie algebroids provide the corresponding
deformation-theoretic recognition objects in the base and finiteness ranges
of their construction and formal-integration theorems
\cite{BrantnerMathew2025,BrantnerCamposNuiten2025,
BrantnerMagidsonNuiten2025,Fu2024}.  The present chapter does not identify
its differentiated object with any of these algebraic models.

When $X=S$, the relative tangent target vanishes in every represented range.
An arbitrary group object $G/S$ first determines the completed one-object
group
\[
    \mu^{\mathrm{mod}}_{G/S}
    =
    G\times_{G_{\mathrm{dR}/S}}S,
\]
and $B_\bullet G$ itself is formal exactly when
$\mu^{\mathrm{mod}}_{G/S}\to G$ is an equivalence.  Any later
Lie-algebraic specialization must therefore differentiate this completed
one-object groupoid.  Further integration or recognition may then compare it,
under the hypotheses of the corresponding theorems, with spectral Lie and
derived Milnor--Moore frameworks originating in the spectral Lie operad and
its Hopf-theoretic Koszul duality \cite{Ching2005,Heine2024}.

